% Mechanically extracted after \mainmatter from the frozen standalone master.

\markboth{EXPOSÉ XXIV. AUTOMORPHISMS OF REDUCTIVE GROUPS}
  {EXPOSÉ XXIV. AUTOMORPHISMS OF REDUCTIVE GROUPS}
% BEGIN inputs/XXIV/tex/components/01_expose_XXIV_title_intro.tex
% SGA 3, Exposé XXIV, Loop-1 English reconstruction.
% Current scope: authority local pages 1--3 / combined-reader pages 1246--1248.
% Local GPU OCR and Reinhold's English Markdown are locator/comparison
% witnesses only. The Polo--Gille page image controls text and formulas.

\setcounter{footnote}{0}
\SGARefTarget{sga3:xxiv:chapter:xxiv}\chapter*{Exposé XXIV\\
Automorphisms of Reductive Groups}
\addcontentsline{toc}{chapter}{Exposé XXIV. Automorphisms of Reductive
Groups}
\markboth{EXPOSÉ XXIV. AUTOMORPHISMS OF REDUCTIVE GROUPS}
  {EXPOSÉ XXIV. AUTOMORPHISMS OF REDUCTIVE GROUPS}

\begin{center}
  \textit{by Michel Demazure}
\end{center}

\begingroup
\renewcommand{\thefootnote}{0}
\footnote{Editor's note: Version dated October 13, 2024.}
\addtocounter{footnote}{-1}
\endgroup
The first part of this Exposé (nos.~1 to 5) is a direct consequence of
the existence, for a reductive group, of ``sufficiently many outer
automorphisms'', a result which is itself a consequence of the weakest
form of the isomorphism theorem for pinned groups. The second part
(nos.~6 and 7) sets out two applications of the more precise results of
the preceding Exposé; in particular, no.~7 uses the generators-and-
relations theorem in its explicit form. Finally, we have included in
an appendix (no.~8) results on ``Galois'' cohomology used in the text.

Let us specify our cohomological notation. If \(S\) is a scheme and
\(G\) an \(S\)-group scheme, we write \(H^1(S,G)\) for the first
cohomology set of \(S\) with coefficients in \(G\), computed for the
\textup{(fpqc)} topology; this is also the set of isomorphism classes
of principal homogeneous \textup{(fpqc)} sheaves under \(G\). We write
\(H^1_{\mathrm{\acute et}}(S,G)\) for the corresponding set for the
étale topology; it is therefore the part of \(H^1(S,G)\) formed by
classes of homogeneous sheaves under \(G\) which are quasi-isotrivial
\textup{(that is, locally trivial for the étale topology)}. We write
\(\Fib(S,G)\) for the part of \(H^1(S,G)\) formed by classes of
representable sheaves \textup{(principal homogeneous bundles)}. Thus
we have the inclusions
\[
  H^1_{\mathrm{\acute et}}(S,G)\subset H^1(S,G),
  \qquad
  \Fib(S,G)\subset H^1(S,G).
\]

If every principal homogeneous sheaf under \(G\) is representable
\textup{(for example, if \(G\) is quasi-affine over \(S\); cf.\ SGA~1,
VIII, \SGARefLink{sga3:viii:proposition:7:9}{7.9})}, then
\[
  \Fib(S,G)=H^1(S,G).
\]

If \(S'\to S\) is a covering morphism for the \textup{(fpqc)}
topology, we write \(H^1(S'/S,G)\) for the kernel of the canonical map
\[
  H^1(S,G)\longrightarrow H^1(S',G_{S'}).
\]
It is known that \(H^1(S'/S,G)\) can be computed simplicially
\textup{(TDTE~I, \S A.4)}, which implies that, when \(S'\to S\) is a
covering for the étale topology, \(H^1(S'/S,G)\) is also the kernel of
\[
  H^1_{\mathrm{\acute et}}(S,G)
  \longrightarrow
  H^1_{\mathrm{\acute et}}(S',G_{S'}).
\]

Finally, following \SGARefLink{sga3:viii:corollary:4:5}{Exp.~VIII, 4.5}, the name ``\SGARefLink{sga3:viii:corollary:4:5}{Theorem~90}'' is given to
the following assertion: ``every principal homogeneous sheaf under
\(G_{m,S}\) is representable and locally trivial.'' This assertion is
equivalent to
\[
  H^1(S,G_{m,S})=\operatorname{Pic}(S),
\]
or again to
\[
  H^1(S,G_{m,S})=0
\]
when \(S\) is local \textup{(or, more generally, semilocal)}.

\SGARefTarget{sga3:xxiv:section:1}
\section*{1. Scheme of Automorphisms of a Reductive Group}
\addcontentsline{toc}{section}{1. Scheme of Automorphisms of a Reductive Group}

\SGARefTarget{sga3:xxiv:1.0}
\subsection*{1.0.}
We must first make precise certain definitions from the preceding
Exposé. Let \(S\) be a nonempty scheme, let \(G\) be a reductive
\(S\)-group, and let
\[
  \mathscr R=(M,M^*,R,R^*,\Delta)
\]
be a pinned reduced root datum \textup{(\SGARefLink{sga3:XXIII:1.5}{Exp.~XXIII, 1.5})}. A
\emph{pinning of \(G\) of type \(\mathscr R\)}, or an
\(\mathscr R\)-\emph{pinning of \(G\)}, consists of:
\begin{enumerate}
\item[(i)] an isomorphism of \(D_S(M)\) onto a maximal torus \(T\) of
  \(G\) \textup{(or, equivalently, a monomorphism
  \(D_S(M)\to G\) whose image is a maximal torus \(T\) of \(G\))},
  identifying \(R\) with a root system of \(G\) relative to \(T\)
  \textup{(\SGARefLink{sga3:xix:definition:3-6}{Exp.~XIX, 3.6})} and \(R^*\) with the corresponding set of
  coroots;
\item[(ii)] for each \(\alpha\in\Delta\), an element
  \[
    X_\alpha\in\Gamma(S,\mathfrak g^\alpha)^\times.
  \]
\end{enumerate}
The group \(G\) possesses an \(\mathscr R\)-pinning if and only if it
is split and of type \(\mathscr R\) \textup{(\SGARefLink{sga3:xxii:def:2.7}{Exp.~XXII, 2.7})}.

If \(u:G\to G'\) is an isomorphism of reductive \(S\)-groups, then
every \(\mathscr R\)-pinning \(\mathcal E\) of \(G\) determines, by
``transport of structure'', an \(\mathscr R\)-pinning
\(u(\mathcal E)\) of \(G'\). If
\(v:\mathscr R'\to\mathscr R\) is an isomorphism of pinned root data,
then every \(\mathscr R'\)-pinning \(\mathcal E'\) of \(G'\)
determines, by transport of structure, an \(\mathscr R\)-pinning
\(v(\mathcal E')\) of \(G\).

Call a \emph{pinned group} a triple \((G,\mathscr R,\mathcal E)\),
where \(G\) is a reductive \(S\)-group, \(\mathscr R\) is a pinned
reduced root datum, and \(\mathcal E\) is a pinning of \(G\) of type
\(\mathscr R\). An isomorphism from the pinned group
\((G,\mathscr R,\mathcal E)\) onto the pinned group
\((G',\mathscr R',\mathcal E')\) is a pair \((u,v)\), where
\(u:G\to G'\) is an isomorphism and
\(v:\mathscr R'\to\mathscr R\) is an isomorphism of pinned root data,
such that
\[
  u(\mathcal E)=v(\mathcal E').
\]
\footnote{Editor's note: Thus
\(\Lie(u)(X_{v(\alpha)})=X'_\alpha\) for every
\(\alpha\in R'\).}

\textit{N.B.} If \(S\) is nonempty, \(v\) is uniquely determined by
\(u\), and by an abuse of language \(u\) will also be called an
isomorphism of the pinned groups. In particular, if
\((G,\mathscr R,\mathcal E)\) is a pinned group, an automorphism of
\((G,\mathscr R,\mathcal E)\) is therefore an automorphism \(u\) of
\(G\) for which there exists an automorphism \(v\) of \(\mathscr R\)
such that \(u(\mathcal E)=v(\mathcal E)\). It is thus an automorphism
of \(G\) which normalizes \(T\), induces on \(T\) an automorphism of
the form \(D_S(h)\), where \(h\) is an automorphism of \(M\), and
permutes the elements \(X_\alpha\), \(\alpha\in\Delta\).
\textup{(As is readily seen, the preceding conditions in fact
characterize the automorphisms of
\((G,\mathscr R,\mathcal E)\).)}

There is an evident contravariant functor
\[
  \Phi:\quad
  (G,\mathscr R,\mathcal E)\longmapsto\mathscr R,
  \qquad
  (u,v)\longmapsto v,
\]
and the principal result of the preceding Exposé
\textup{(\SGARefLink{sga3:xxiii:theorem:4-1}{Exp.~XXIII, 4.1})} shows that this functor is fully faithful
\textup{(we shall moreover see in the following Exposé that it is an
equivalence of categories)}. In particular, it follows that the
automorphism group of \((G,\mathscr R,\mathcal E)\) is canonically
isomorphic to the automorphism group of the pinned root datum
\(\mathscr R\) \textup{(cf.\ \SGARefLink{sga3:XXIII:5.5}{Exp.~XXIII, 5.5})}.

\SGARefTarget{sga3:xxiv:1.1}
\subsection*{1.1.}
Let \(S\) be a scheme; endow \((\mathbf{Sch})/S\) with the
\textup{(fpqc)} topology, and consider the \(S\)-sheaf of groups
\(\underline{\Aut}_{S\text{-}\mathrm{gr.}}(G)\), where \(G\) is an
\(S\)-group scheme. There is an exact sequence of \(S\)-sheaves of
groups
\[
  1\longrightarrow \underline{\Centr}(G)
  \longrightarrow G
  \xrightarrow{\mathrm{int}}
  \underline{\Aut}_{S\text{-}\mathrm{gr.}}(G),
\]
which defines a monomorphism
\[
  j:G/\underline{\Centr}(G)
  \longrightarrow \underline{\Aut}_{S\text{-}\mathrm{gr.}}(G).
\]
The image sheaf of \(j\) is the sheaf of inner automorphisms of \(G\).
An automorphism \(u\) of \(G\) is inner if and only if there is a
covering family \(\{S_i\to S\}\) and, for every \(i\), an element
\(g_i\in G(S_i)\) such that
\[
  \operatorname{int}(g_i)=u_{S_i}.
\]
In this case, if \(v\) is another automorphism of \(G\), one sees at
once that
\[
  \operatorname{int}(v)u=vuv^{-1}
\]
is the inner automorphism defined by the family
\(g'_i=v(g_i)\). It follows that the image of \(j\) is normal in
\(\underline{\Aut}_{S\text{-}\mathrm{gr.}}(G)\). The quotient sheaf of
groups, denoted \(\underline{\Autext}(G)\), is the sheaf of outer
automorphisms of \(G\). Thus there is an exact sequence
\[
  1\longrightarrow G/\underline{\Centr}(G)
  \longrightarrow \underline{\Aut}_{S\text{-}\mathrm{gr.}}(G)
  \longrightarrow \underline{\Autext}(G)
  \longrightarrow 1.
\]
All the preceding definitions are compatible with base change. They
are naturally valid in any site.

\SGARefTarget{sga3:xxiv:1.2}
\subsection*{1.2.}
Let \(S\) be a scheme and let \((G,\mathscr R,\mathcal E)\) be a
pinned reductive \(S\)-group. Let \(E\) be the \textup{(abstract)}
automorphism group of the pinned root datum \(\mathscr R\), that is,
the group of automorphisms of \(\mathscr R\) which normalize
\(\Delta\). By \SGARefLink{sga3:XXIII:5.5}{Exp.~XXIII, 5.5}, there is a canonical monomorphism
\[
  E\longrightarrow \underline{\Aut}_{S\text{-}\mathrm{gr.}}(G)
\]
which associates with \(h\in E\) the unique automorphism \(u\) of the
pinned group \(G\) such that \(\Phi(u)=h\). This monomorphism
canonically defines a monomorphism of sheaves
\[
  \tag{*}\label{sga3:xxiv:subsection:1-2:equation:star}
  a:E_S\longrightarrow \underline{\Aut}_{S\text{-}\mathrm{gr.}}(G).
\]
An automorphism \(u\) of \(G\) is a section of the image sheaf of \(a\)
if and only if the following conditions hold:
\begin{enumerate}
\item[(i)] \(u\) normalizes \(T\). It is then known that \(u\)
  permutes the roots of \(G\) relative to \(T\). If \(\alpha\in R\),
  then
  \[
    u(\alpha):t\longmapsto\alpha\bigl(u^{-1}(t)\bigr)
  \]
  is therefore locally on \(S\) of the form
  \(t\longmapsto\beta(t)\), with \(\beta\in R\). The second condition
  may then be written as follows:
\item[(ii)] If \(\alpha\in\Delta\) and \(U\) is an open subset of
  \(S\) such that \(u(\alpha)_U\in R\), then
  \(u(\alpha)_U\in\Delta\) and
  \[
    \Lie(u_U)(X_\alpha)_U=(X_{u(\alpha)})_U.
  \]
\end{enumerate}
It follows at once from the definitions that the sections of
\(a(E_S)\) normalize the subgroups of \(G\) defined by the pinning:
\[
  T,\quad B,\quad B^{-},\quad U,\quad U^{-}.
\]

With these definitions in place, we have:

% END inputs/XXIV/tex/components/01_expose_XXIV_title_intro.tex

% BEGIN inputs/XXIV/tex/components/02_expose_XXIV_theorem13_lemma15_uniqueness.tex
% SGA 3, Exposé XXIV, authority local page 4 / combined-reader page 1249.
% Personally checked against the Polo--Gille page image, band by band.

\SGARefTarget{sga3:xxiv:1.3}
\paragraph{\SGARefLink{sga3:xxiv:1.3}{Theorem 1.3}.}
\emph{Let \(S\) be a scheme and \(G\) a reductive \(S\)-group. Consider
the canonical exact sequence of \(S\)-sheaves of groups}
\footnote{Editor's note: Recall that
\(\operatorname{ad}(G)=G/\underline{\Centr}(G)\) denotes the adjoint
group of \(G\); cf.\ \SGARefLink{sga3:xxii:def:4.3.6}{Exp.~XXII, 4.3.6}.}
\[
  1\longrightarrow \operatorname{ad}(G)
  \longrightarrow \underline{\Aut}_{S\text{-}\mathrm{gr.}}(G)
  \xrightarrow{p}\underline{\Autext}(G)
  \longrightarrow 1.
\]
\emph{Then:}
\begin{enumerate}
\item[\SGAInlineRefTarget{sga3:xxiv:1.3:item:i}(i)] \(\underline{\Aut}_{S\text{-}\mathrm{gr.}}(G)\) is
  representable by a smooth separated \(S\)-scheme;
\item[\SGAInlineRefTarget{sga3:xxiv:1.3:item:ii}(ii)] \(\underline{\Autext}(G)\) is representable by a twisted
  constant \(S\)-scheme of finite generation \textup{(\SGARefLink{sga3:x:definition:5-1}{Exp.~X, 5.1})};
\item[\SGAInlineRefTarget{sga3:xxiv:1.3:item:iii}(iii)] if \(G\) is split, the preceding exact sequence is split.
  More precisely, for every pinning of \(G\), the morphism
  \textup{(cf.\ \SGARefLink{sga3:xxiv:1.2}{1.2}~\((*)\))}
  \[
    p\circ a:E_S\longrightarrow\underline{\Autext}(G)
  \]
  is an isomorphism.
\end{enumerate}

Let us first show how the theorem follows from the following lemma.

\SGARefTarget{sga3:xxiv:1.4}
\paragraph{\SGARefLink{sga3:xxiv:1.4}{Lemma 1.4}.}
\emph{Under the hypotheses of \textup{\SGARefLink{sga3:xxiv:1.3:item:iii}{(iii)}},
\(\underline{\Aut}_{S\text{-}\mathrm{gr.}}(G)\) is the semidirect
product}
\[
  a(E_S)\cdot\operatorname{ad}(G).
\]
\footnote{Editor's note: Here and below, the notation
\(\operatorname{int}(G)\) has been replaced by
\(\operatorname{ad}(G)\).}

The lemma immediately implies the theorem when \(G\) is split. Since
\(G\) is locally split for the étale topology \textup{(\SGARefLink{sga3:xxii:cor:2.3}{Exp.~XXII,
2.3})}, hence also for the \textup{(fppf)} topology, and since the
latter is ``of effective descent'' for the fibred category of twisted
constant morphisms \textup{(\SGARefLink{sga3:x:corollary:5-5}{Exp.~X, 5.5})}, assertion \textup{\SGARefLink{sga3:xxiv:1.3:item:ii}{(ii)}}
follows in the general case \textup{(cf.\ \SGARefLink{sga3:iv:proposition:4:6:8}{Exp.~IV, 4.6.8})}. To deduce
\textup{\SGARefLink{sga3:xxiv:1.3:item:i}{(i)}}, observe that \(\operatorname{ad}(G)\) is affine over
\(S\); thus \(p\) is affine whenever
\(\underline{\Aut}_{S\text{-}\mathrm{gr.}}(G)\) is representable, and
the result follows by descent for affine schemes.
\footnote{Editor's note: Cf.\ SGA~1, VIII, \SGARefLink{sga3:viii:proposition:2:1}{2.1}.}

It remains only to prove \SGARefLink{sga3:xxiv:1.4}{1.4}. For this, it suffices to prove:

\SGARefTarget{sga3:xxiv:1.5}
\paragraph{\SGARefLink{sga3:xxiv:1.5}{Lemma 1.5}.}
\emph{If \((\mathscr R,\mathcal E)\) and
\((\mathscr R',\mathcal E')\) are two pinnings of the reductive
\(S\)-group \(G\), there exists a unique inner automorphism \(u\) of
\(G\) over \(S\) carrying one pinning to the other; that is, such that
there exists an isomorphism
\[
  v:\mathscr R'\xrightarrow{\sim}\mathscr R
\]
for which \(u(\mathcal E)=v(\mathcal E')\); cf.\ \SGARefLink{sga3:xxiv:1.0}{1.0}.}

\SGARefTarget{sga3:xxiv:1.5.1}
\subsubsection*{1.5.1. Uniqueness.}
It suffices to prove that, if \(G\) is a pinned \(S\)-group and
\(\operatorname{int}(g)\) is an automorphism of the pinned group
\textup{(\(g\in G(S)\))}, then \(\operatorname{int}(g)=\operatorname{id}\).
First,
\[
  \operatorname{int}(g)T=T,\qquad
  \operatorname{int}(g)B=B,
\]
and hence
\[
  g\in
  \underline{\operatorname{Norm}}_G(T)(S)
  \cap\underline{\operatorname{Norm}}_G(B)(S)
  =T(S)
\]
\textup{(cf., for example, \SGARefLink{sga3:xxii:prop:5.6.1}{Exp.~XXII, 5.6.1})}. It follows that
\(\operatorname{int}(g)\) normalizes every \(U_\alpha\) and that
\[
  \Lie(\operatorname{int}(g))X_\alpha
  =\operatorname{Ad}(g)X_\alpha
  =\alpha(g)X_\alpha
\]
for every \(\alpha\in\Delta\). Thus \(\alpha(g)=1\) for
\(\alpha\in\Delta\), and therefore
\[
  g\in
  \bigcap_{\alpha\in\Delta}(\ker\alpha)(S)
  =\underline{\Centr}(G)(S)
\]
\textup{(\SGARefLink{sga3:xxii:rem:4.1.8}{Exp.~XXII, 4.1.8})}. \qed

% END inputs/XXIV/tex/components/02_expose_XXIV_theorem13_lemma15_uniqueness.tex

% BEGIN inputs/XXIV/tex/components/03_expose_XXIV_lemma15_existence_corollaries16_18.tex
% SGA 3, Exposé XXIV, authority local page 5 / combined-reader page 1250.
% Personally checked against the Polo--Gille page image, band by band.

\SGARefTarget{sga3:xxiv:1.5.2}
\subsubsection*{1.5.2. Existence.}
It suffices to prove existence locally for the \textup{(fpqc)}
topology. Let
\[
  (T,M,R,\Delta,(X_\alpha)_{\alpha\in\Delta})
  \quad\text{and}\quad
  (T',M',R',\Delta',(X'_{\alpha'})_{\alpha'\in\Delta'})
\]
be the two pinnings. By conjugacy of maximal tori, we may assume
\(T=T'\). After restricting \(S\), we may assume that the isomorphism
\[
  D_S(M)\simeq D_S(M')
\]
comes from an isomorphism \(M\simeq M'\) carrying \(R\) onto \(R'\);
we are thus reduced to the situation
\[
  T=T',\qquad M=M',\qquad R=R'.
\]
Since the systems of simple roots are conjugate under the Weyl group
\textup{(\SGARefLink{sga3:XXI:3.3.7}{Exp.~XXI, 3.3.7})}, we may also assume
\(\Delta=\Delta'\). For every \(\alpha\in\Delta\), there is then a
scalar \(z_\alpha\in\mathbb G_m(S)\) such that
\[
  X'_\alpha=z_\alpha X_\alpha,
\]
and it is enough to construct, locally for the \textup{(fpqc)}
topology, a section \(t\) of \(T\) such that
\(\alpha(t)=z_\alpha\) for every \(\alpha\in\Delta\). But the
morphism
\[
  T\longrightarrow(\mathbb G_{m,S})^\Delta
\]
whose components are the \(\alpha\), \(\alpha\in\Delta\), is dual to
an injection \(\mathbb Z^\Delta\to M\), and is therefore faithfully
flat. This completes the proof of \SGARefLink{sga3:xxiv:1.5.2}{1.5.2}, and hence of \SGARefLink{sga3:xxiv:1.3}{1.3}.

\SGARefTarget{sga3:xxiv:1.6}
\paragraph{\SGARefLink{sga3:xxiv:1.6}{Corollary 1.6}.}
\emph{Let \(S\) be a scheme and \(G\) a reductive \(S\)-group. The
following conditions are equivalent:}
\begin{enumerate}
\item[\SGAInlineRefTarget{sga3:xxiv:1.6:item:i}(i)] \(\underline{\Aut}_{S\text{-}\mathrm{gr.}}(G)\) is affine
  \emph{\textup{(respectively of finite type, respectively of finite
  presentation, respectively quasi-compact)}} over \(S\);
\item[\SGAInlineRefTarget{sga3:xxiv:1.6:item:ii}(ii)] \(\underline{\Autext}(G)\) is finite over \(S\);
\item[\SGAInlineRefTarget{sga3:xxiv:1.6:item:iii}(iii)] for every \(s\in S\),
  \[
    \operatorname{rg}_{\mathrm{red}}(G_s)
    -\operatorname{rg}_{\mathrm{ss}}(G_s)\leq 1.
  \]
\end{enumerate}

Indeed, since \(\operatorname{ad}(G)\) is affine, flat, and of finite
presentation over \(S\), the morphism
\[
  p:\underline{\Aut}_{S\text{-}\mathrm{gr.}}(G)
  \longrightarrow\underline{\Autext}(G)
\]
is affine, faithfully flat, and of finite presentation.

If \(\underline{\Autext}(G)\) is finite over \(S\), it is affine over
\(S\), and hence so is
\(\underline{\Aut}_{S\text{-}\mathrm{gr.}}(G)\); this proves
\textup{\SGARefLink{sga3:xxiv:1.6:item:ii}{(ii)}}\(\Rightarrow\)\textup{\SGARefLink{sga3:xxiv:1.6:item:i}{(i)}}. If
\(\underline{\Aut}_{S\text{-}\mathrm{gr.}}(G)\) is quasi-compact over
\(S\), it is of finite presentation over \(S\), since it is in any
case locally of finite presentation and separated over \(S\). By
\SGARefLink{sga3:v:proposition:9:1}{Exp.~V, 9.1}, \(\underline{\Autext}(G)\) is then of finite presentation
over \(S\), and hence finite; this proves
\textup{\SGARefLink{sga3:xxiv:1.6:item:i}{(i)}}\(\Rightarrow\)\textup{\SGARefLink{sga3:xxiv:1.6:item:ii}{(ii)}}. Finally, to prove the
equivalence of \textup{\SGARefLink{sga3:xxiv:1.6:item:ii}{(ii)}} and \textup{\SGARefLink{sga3:xxiv:1.6:item:iii}{(iii)}}, one may suppose that
\(G\) is split, and the assertion reduces to \SGARefLink{sga3:XXI:6.7.8}{Exp.~XXI, 6.7.8}.

\SGARefTarget{sga3:xxiv:1.7}
\paragraph{\SGARefLink{sga3:xxiv:1.7}{Corollary 1.7}.}
\emph{Let \(S\) be a scheme and \(G\) a reductive \(S\)-group. Then}
\[
  \underline{\Aut}_{S\text{-}\mathrm{gr.}}(G)^0
  \simeq\operatorname{ad}(G).
\]

\SGARefTarget{sga3:xxiv:1.8}
\paragraph{\SGARefLink{sga3:xxiv:1.8}{Corollary 1.8}.}
\emph{Let \(S\) be a scheme, \(G\) a reductive \(S\)-group, and
\(H\) an \(S\)-group scheme that is smooth, affine, and has connected
fibres over \(S\). Then the \(S\)-functor}
\[
  \underline{\operatorname{Isom}}_{S\text{-}\mathrm{gr.}}(G,H)
\]
\emph{is representable by a smooth separated \(S\)-scheme, which is
affine over \(S\) if \(G\) is semisimple.}

Indeed, let \(U\) be the set of points \(s\in S\) for which \(H_s\)
is reductive; this is an open subset \textup{(\SGARefLink{sga3:xix:corollary:2:6}{Exp.~XIX, 2.6})}. If
\(S'\) is an \(S\)-scheme, \(H_{S'}\) is reductive if and only if
\(S'\to S\) factors through \(U\). It follows that the canonical
morphism
\[
  \underline{\operatorname{Isom}}_{S\text{-}\mathrm{gr.}}(G,H)
  \longrightarrow S
\]
factors through \(U\). We may therefore assume \(S=U\), and we are
reduced to the following corollary.

% END inputs/XXIV/tex/components/03_expose_XXIV_lemma15_existence_corollaries16_18.tex

% BEGIN inputs/XXIV/tex/components/04_expose_XXIV_corollaries19_110_remark111.tex
% SGA 3, Exposé XXIV, authority local page 6 / combined-reader page 1251.
% Personally checked against the Polo--Gille page image, band by band.

\SGARefTarget{sga3:xxiv:1.9}
\paragraph{\SGARefLink{sga3:xxiv:1.9}{Corollary 1.9}.}
\emph{Let \(S\) be a scheme, and let \(G\) and \(G'\) be two
reductive \(S\)-groups. Then}
\[
  F=\underline{\operatorname{Isom}}_{S\text{-}\mathrm{gr.}}(G,G')
\]
\emph{is representable by a smooth separated \(S\)-scheme, affine
if \(G\) or \(G'\) is semisimple. Moreover, \(S\) is the disjoint
union of two open subschemes \(S_1\) and \(S_2\) such that
\(F_{S_1}=\varnothing\), while \(F_{S_2}\) is a principal homogeneous
bundle on the left \textup{(respectively, on the right)} under
\(\underline{\Aut}_{S_2\text{-}\mathrm{gr.}}(G'_{S_2})\)
\textup{(respectively,
\(\underline{\Aut}_{S_2\text{-}\mathrm{gr.}}(G_{S_2})\)).}}

Indeed, let \(S_2\) be the set of points of \(S\) at which \(G\) and
\(G'\) have the same type, and let \(S_1\) be its complement. Since
the type of a reductive group is a locally constant function, \(S_1\)
and \(S_2\) are open. Clearly \(F_{S_1}=\varnothing\), and we may
assume \(S=S_2\). By \SGARefLink{sga3:XXIII:5.6}{Exp.~XXIII, 5.6}, \(F\) is a principal
homogeneous sheaf under
\(\underline{\Aut}_{S\text{-}\mathrm{gr.}}(G)\), locally trivial for
the étale topology. It follows that
\[
  F_0=F/\operatorname{ad}(G)
\]
is a principal homogeneous sheaf under \(\underline{\Autext}(G)\),
locally trivial for the étale topology, and hence representable
\textup{(\SGARefLink{sga3:x:corollary:5-5}{Exp.~X, 5.5})}. Since \(\operatorname{ad}(G)\) is affine,
\(F\) is therefore representable as well.

We note that the proof also gives the following result.

\SGARefTarget{sga3:xxiv:1.10}
\paragraph{\SGARefLink{sga3:xxiv:1.10}{Corollary 1.10}.}
\emph{Let \(S\) be a scheme, and let \(G\) and \(G'\) be two
reductive \(S\)-groups having the same type at every \(s\in S\).
Then \(\operatorname{ad}(G)\) acts freely on the right on
\(\underline{\operatorname{Isom}}_{S\text{-}\mathrm{gr.}}(G,G')\).
The quotient sheaf}
\[
  \underline{\operatorname{Isomext}}(G,G')
  =
  \underline{\operatorname{Isom}}_{S\text{-}\mathrm{gr.}}(G,G')
  /\operatorname{ad}(G)
\]
\emph{is representable by a twisted constant \(S\)-scheme; it is a
principal homogeneous bundle under \(\underline{\Autext}(G)\), and
is therefore finite over \(S\) if \(G\) is semisimple. Moreover, the
isomorphism}
\[
  \underline{\operatorname{Isom}}_{S\text{-}\mathrm{gr.}}(G,G')
  \simeq
  \underline{\operatorname{Isom}}_{S\text{-}\mathrm{gr.}}(G',G),
  \qquad u\longmapsto u^{-1},
\]
\emph{induces an isomorphism}
\[
  \underline{\operatorname{Isomext}}(G,G')
  \simeq
  \underline{\operatorname{Isomext}}(G',G).
\]

\SGARefTarget{sga3:xxiv:1.11}
\paragraph{\SGARefLink{sga3:xxiv:1.11}{Remark 1.11}.}
If
\(\underline{\operatorname{Isomext}}(G,G')(S)\neq\varnothing\), one
says that \(G\) is an \emph{inner twisted form} of \(G'\). Then \(G'\)
is an inner twisted form of \(G\), and the structure group of
\(\underline{\operatorname{Isom}}_{S\text{-}\mathrm{gr.}}(G,G')\)
can be reduced to \(\operatorname{ad}(G)\). More precisely, let
\[
  u\in\underline{\operatorname{Isomext}}(G,G')(S)
\]
and regard \(u\) as a section
\(u:S\to\underline{\operatorname{Isomext}}(G,G')\). Denote by
\[
  \underline{\operatorname{Isomint}}_u(G,G')
\]
the inverse image, under the canonical morphism
\[
  \underline{\operatorname{Isom}}_{S\text{-}\mathrm{gr.}}(G,G')
  \longrightarrow
  \underline{\operatorname{Isomext}}(G,G'),
\]
of the closed subscheme of
\(\underline{\operatorname{Isomext}}(G,G')\) that is the image of
\(u\). The natural action of \(\operatorname{ad}(G)\) on
\(\underline{\operatorname{Isomint}}_u(G,G')\) makes this scheme a
principal homogeneous bundle. By extension of the structure
group\footnote{The source's editorial note corrects
\(\operatorname{ad}(G)\to\underline{\Autext}(G)\) to
\(\operatorname{ad}(G)\to\underline{\Aut}(G)\).}
\[
  \operatorname{ad}(G)\longrightarrow\underline{\Aut}(G),
\]
one recovers
\(\underline{\operatorname{Isom}}_{S\text{-}\mathrm{gr.}}(G,G')\)
from \(\underline{\operatorname{Isomint}}_u(G,G')\).

By Hensel's lemma \textup{(\SGARefLink{sga3:xi:corollary:1:11}{Exp.~XI, 1.11})}, \SGARefLink{sga3:xxiv:1.8}{Corollary~1.8} immediately
gives:

% END inputs/XXIV/tex/components/04_expose_XXIV_corollaries19_110_remark111.tex

% BEGIN inputs/XXIV/tex/components/05_expose_XXIV_corollaries112_115.tex
% SGA 3, Exposé XXIV, authority local page 7 / combined-reader page 1252.
% Personally checked against the Polo--Gille page image, band by band.

\SGARefTarget{sga3:xxiv:1.12}
\paragraph{\SGARefLink{sga3:xxiv:1.12}{Corollary 1.12}.}
\emph{Let \(S\) be a henselian local scheme, let \(G\) be a
reductive \(S\)-group, let \(G'\) be a smooth affine \(S\)-group with
connected fibres, and let \(s\) be the closed point of \(S\). If
\(G_s\) and \(G'_s\) are isomorphic algebraic groups over
\(\kappa(s)\), then \(G\) and \(G'\) are isomorphic. More precisely,
every \(\kappa(s)\)-isomorphism}
\[
  G_s\simeq G'_s
\]
\emph{comes from an \(S\)-isomorphism \(G\simeq G'\).}

Applying \SGARefLink{sga3:xxiv:1.7}{Corollary~1.7} \textup{(respectively, \SGARefLink{sga3:xxiv:1.12}{Corollary~1.12})} to the
scheme of dual numbers over a field, \SGARefLink{sga3:iii:remark:2-10}{Exp.~III, 2.10}
\textup{(respectively, \SGARefLink{sga3:xxiv:3.10}{3.10})} yields part \textup{(i)}
\textup{(respectively, part \textup{(ii)})} of the next corollary.

\SGARefTarget{sga3:xxiv:1.13}
\paragraph{\SGARefLink{sga3:xxiv:1.13}{Corollary 1.13}.}
\emph{Let \(k\) be a field and \(G\) a reductive \(k\)-group.}
\begin{enumerate}
\item[\SGAInlineRefTarget{sga3:xxiv:1.13:item:i}(i)] \emph{If \(G\) is adjoint, then}
\[
  H^1(G,\Lie(G/k))=0.
\]
\footnote{The source's editorial note corrects the original, which
stated \textup{\SGARefLink{sga3:xxiv:1.13:item:i}{(i)}} without assuming \(G\) adjoint. By the
description of \(H^i(G,-)\) as the derived functors of
\(\operatorname{Hom}_G(k,-)\) \textup{(\SGARefLink{sga3:i:theorem:5-3-1}{Exp.~I, 5.3.1}; see also
[Ja87], I, 4.16)}, one has
\(H^i(G,V)=\operatorname{Ext}^i_G(k,V)\) for every \(G\)-module
\(V\). If \(\operatorname{char}(k)=p>0\) and
\(G=\operatorname{SL}_{p,k}\), then
\(\Lie(\operatorname{GL}_p/k)\) is a nontrivial extension of \(k\)
by \(\mathfrak g=\Lie(\operatorname{SL}_p/k)\), so
\(H^1(G,\mathfrak g)\neq0\). See also the addition \SGARefLink{sga3:xxiv:1.15.1}{1.15.1} below.}
\item[\SGAInlineRefTarget{sga3:xxiv:1.13:item:ii}(ii)] \emph{One has}
\[
  H^2(G,\Lie(G/k))=0.
\]
\end{enumerate}

\SGARefTarget{sga3:xxiv:1.14}
\paragraph{\SGARefLink{sga3:xxiv:1.14}{Remark 1.14}.}
\begin{enumerate}
\item[(i)] The assertion concerning \(H^1\) was known to Chevalley;
the assertion concerning \(H^2\) had been proved in most cases of
the classification by Chevalley.
\item[(ii)] In fact, the conjunction of \SGARefLink{sga3:xxiv:1.13}{Corollary~1.13} and the
uniqueness theorem over an algebraically closed field is essentially
equivalent to the uniqueness theorem. A direct proof of
\SGARefLink{sga3:xxiv:1.13}{Corollary~1.13} would therefore give a way to deduce the general
uniqueness theorem from Chevalley's uniqueness theorem over a field.
\end{enumerate}

The existence of reductive groups of every type over every scheme
\textup{(Exp.~XXV)} shows that the obstructions to lifting a
reductive \(k\)-group \(G\) over artinian rings with residue field
\(k\), which by \SGARefLink{sga3:iii:corollary:3-8}{Exp.~III, 3.8} are elements of
\[
  H^3(G,\Lie(G/k)\otimes V)
  \simeq H^3(G,\Lie(G/k))\otimes V,
\]
where \(V\) is a certain \(k\)-vector space endowed with the trivial
action of \(G\), vanish. This suggests that
\[
  H^3(G,\Lie(G/k))=0.
\]
A direct proof, if the assertion is true, would likewise doubtless
give a way to deduce the general existence theorem from the
existence theorem over a field \textup{(Chevalley's Tôhoku)}.

\SGARefTarget{sga3:xxiv:1.15}
\paragraph{\SGARefLink{sga3:xxiv:1.15}{Corollary 1.15}.}
\emph{Let \(k\) be a field,\footnote{The source's editorial note
replaces ``scheme'' with ``field.''} and let \(G\) be a reductive
\(k\)-group. Regard \(k\) as a trivial \(G\)-module. Then}
\[
  H^1(G,k)=H^2(G,k)=0.
\]
\footnote{Because of the correction made in \SGARefLink{sga3:xxiv:1.13}{1.13}, the source's
editors supply another proof for \(H^1\). Moreover, a theorem of
G.~Kempf gives \(H^i(G,k)=0\) for every \(i>0\); cf.\ [Ja87],
II, \SGARefLink{sga3:ii:proposition:4-5}{4.5} and 4.11.}
Since
\[
  H^i(G_{\bar k},\bar k)=H^i(G,k)\otimes_k\bar k,
\]
we may assume \(k\) algebraically closed. An element of \(H^1(G,k)\)
is simply a morphism of \(k\)-groups
\[
  \varphi:G\longrightarrow\mathbb G_{a,k}.
\]
Then
\[
  \varphi(G)=G/\ker(\varphi)
\]
is a smooth connected reductive subgroup of
\(\mathbb G_{a,k}\) \textup{(cf.\ XIX, \SGARefLink{sga3:xix:paragraph:1-7}{1.7})}, and is therefore
trivial. Thus \(H^1(G,k)=0\).

% END inputs/XXIV/tex/components/05_expose_XXIV_corollaries112_115.tex

% BEGIN inputs/XXIV/tex/components/06_expose_XXIV_corollary115_through_proposition28_en.tex
% Exposé XXIV, authority local pages 8--12 / combined-reader pages 1253--1257.
% Authority local p. 8 / combined-reader p. 1253.
Now consider the reductive \(k\)-group
\[
  H=G\times_k\mathbb G_{m,k}.
\]
There is an \(H\)-stable decomposition
\[
  \Lie(H/k)=\Lie(G/k)\oplus k.
\]
For every \(H\)-module \(V\),
\[
  H^i(H,V)=H^i\bigl(G,H^0(\mathbb G_{m,k},V)\bigr).
\]
Indeed, the \(H^i(H,-)\) are the derived functors of
\[
  H^0(H,-)
  =H^0\bigl(G,H^0(\mathbb G_{m,k},-)\bigr),
\]
and the functor \(H^0(\mathbb G_{m,k},-)\) is exact
(cf.\ \SGARefLink{sga3:i:theorem:5-3-1}{Exp.~I, 5.3.1} and \SGARefLink{sga3:i:theorem:5-3-3}{5.3.3}). In particular, for every \(i\),
\[
  H^i(H,\Lie(H/k))
  =
  H^i(G,\Lie(G/k))\oplus H^i(G,k).
\]
Applying \SGARefLink{sga3:xxiv:1.13:item:ii}{1.13(ii)} to the reductive group \(H\) gives
\(H^2(G,k)=0\).

\SGARefTarget{sga3:xxiv:1.15.1}
\medskip\noindent\textbf{Corollary 1.15.1.} ---
\footnote{Editor's note: This corollary, drawn from remarks of
O.~Gabber, has been added to make \SGARefLink{sga3:xxiv:1.13:item:i}{1.13(i)} more precise.}
{\itshape Let \(k\) be a field, \(G\) a reductive \(k\)-group, \(Z\)
its center, and \(\mathscr R=(M,M^*,R,R^*)\) the type of \(G\). Then
\[
  H^1(G,\Lie(G/k))
  \simeq
  \operatorname{Ext}^1_{\mathbb Z}(M/\Gamma_0(R),k).
\]
In particular, \(H^1(G,\Lie(G/k))=0\) if and only if \(Z\) is smooth
over \(k\).}

Let \(\mathfrak g\), \(\mathfrak z\), and
\(\mathfrak g_{\mathrm{ad}}\) be the Lie algebras of \(G\), \(Z\), and
\(G_{\mathrm{ad}}=G/Z\), respectively. It follows from \SGARefLink{sga3:xxiv:1.15}{1.15} and its
proof that
\[
  H^1(G,\mathfrak g)
  =
  H^1(G_{\mathrm{ad}},\mathfrak g)
  \simeq
  H^1(G_{\mathrm{ad}},\mathfrak g/\mathfrak z).
\]
Set \(C=\operatorname{Coker}(\mathfrak g\to
\mathfrak g_{\mathrm{ad}})\). There is an exact sequence
\[
  0\longrightarrow\mathfrak g/\mathfrak z
  \longrightarrow\mathfrak g_{\mathrm{ad}}
  \longrightarrow C\longrightarrow0.
\]
Since \(H^1(G_{\mathrm{ad}},\mathfrak g_{\mathrm{ad}})=0\) by \SGARefLink{sga3:xxiv:1.13}{1.13}
and
\[
  H^0(G_{\mathrm{ad}},\mathfrak g_{\mathrm{ad}})=0
\]
(cf.\ \SGARefLink{sga3:ii:theorem:5-2-3}{Exp.~II, 5.2.3}), we obtain
\[
  H^1(G_{\mathrm{ad}},\mathfrak g/\mathfrak z)
  \simeq H^0(G_{\mathrm{ad}},C).
\]
To calculate the right-hand side, we may assume \(k\) algebraically
closed. Let \((G,T,M,R)\) be a splitting of \(G\). Then
\[
  Z=D_k(M/\Gamma_0(R))
\]
and \(C\) is identified with
\[
  \operatorname{Coker}\!\left(
    \operatorname{Hom}_{\mathbb Z}(M,k)
    \longrightarrow
    \operatorname{Hom}_{\mathbb Z}(\Gamma_0(R),k)
  \right)
  \simeq
  \operatorname{Ext}^1_{\mathbb Z}(M/\Gamma_0(R),k),
\]
with trivial \(G_{\mathrm{ad}}\)-action. The corollary follows, since
\(Z\) is not smooth over \(k\) exactly when
\(\operatorname{char}(k)=p>0\) and \(M/\Gamma_0(R)\) has
\(p\)-torsion (\SGARefLink{sga3:ix:proposition:2:1}{Exp.~IX, 2.1}).

\SGARefTarget{sga3:xxiv:1.16}
\medskip\noindent\textbf{Definition 1.16.} ---
{\itshape Let \(S\) be a scheme and \(G\) a reductive \(S\)-group. A
\emph{form of \(G\) over \(S\)} is an \(S\)-group scheme \(G'\)
locally isomorphic to \(G\) for the \textup{(fpqc)} topology. It is
equivalent to require that \(G'\) be locally isomorphic to \(G\) for
the étale topology (cf.\ \SGARefLink{sga3:XXIII:5.6}{Exp.~XXIII, 5.6}), or that \(G'\) be a
reductive \(S\)-group of the same type as \(G\) at every point of
\(S\).}

\SGARefTarget{sga3:xxiv:1.17}
\medskip\noindent\textbf{Corollary 1.17.} ---
{\itshape Let \(S\) be a scheme and \(G\) a reductive \(S\)-group.
\begin{enumerate}
\item[(i)] The functor
\[
  G'\longmapsto
  \underline{\operatorname{Isom}}_{S\text{-}\mathrm{gr}}(G,G')
\]
is an equivalence between the category of forms of \(G\) over \(S\)
and the category of principal homogeneous bundles under
\(\underline{\Aut}_{S\text{-}\mathrm{gr}}(G)\).
\item[(ii)] If \(S'\to S\) is a covering morphism, forms of \(G\)
trivialized by \(S'\) correspond to bundles trivialized by \(S'\).
\item[\SGAInlineRefTarget{sga3:xxiv:1.17:item:iii}(iii)] Every principal homogeneous sheaf under
\(\underline{\Aut}_{S\text{-}\mathrm{gr}}(G)\) is representable and
quasi-isotrivial, that is, locally trivial for the étale topology.
\end{enumerate}}

% Authority local p. 9 / combined-reader p. 1254.
The first assertion is formal in the category of sheaves, for example
for the \textup{(fpqc)} topology. Every sheaf locally isomorphic to
\(G\) for the \textup{(fpqc)} topology is representable, since \(G\)
is affine over \(S\), and is locally isomorphic to \(G\) for the
étale topology. Finally, for every form \(G'\) of \(G\), the
\(S\)-sheaf
\(\underline{\operatorname{Isom}}_{S\text{-}\mathrm{gr}}(G,G')\) is
representable by \SGARefLink{sga3:xxiv:1.8}{1.8}. The corollary follows.

\SGARefTarget{sga3:xxiv:1.18}
\medskip\noindent\textbf{Corollary 1.18.} ---
{\itshape The set of isomorphism classes of forms of the reductive
group \(G\) over \(S\) is isomorphic to
\[
  H^1\!\left(S,\underline{\Aut}_{S\text{-}\mathrm{gr}}(G)\right)
  =
  H^1_{\mathrm{\acute et}}\!
    \left(S,\underline{\Aut}_{S\text{-}\mathrm{gr}}(G)\right)
  =
  \Fib\!\left(S,\underline{\Aut}_{S\text{-}\mathrm{gr}}(G)\right).
\]
If \(S'\to S\) is a covering morphism, the subset consisting of forms
trivialized by \(S'\) is isomorphic to
\[
  H^1\!\left(S'/S,\underline{\Aut}_{S\text{-}\mathrm{gr}}(G)\right).
\]}

\SGARefTarget{sga3:xxiv:1.19}
\medskip\noindent\textbf{Corollary 1.19.} ---
{\itshape Let \(S\) be a scheme and \(\mathscr R\) a pinned reduced root
datum such that \(\operatorname{\acute{E}p}_S(\mathscr R)\)
exists,\footnote{Editor's note: Cf.\ XXIII, \SGARefLink{sga3:XXIII:5.11}{Definition 5.11}.} a
condition which is automatically satisfied by Exp.~XXV. Put
\[
  A_S(\mathscr R)
  =
  \underline{\Aut}_{S\text{-}\mathrm{gr}}
    (\operatorname{\acute{E}p}_S(\mathscr R))
  =
  \operatorname{ad}(\operatorname{\acute{E}p}_S(\mathscr R))
  \cdot E(\mathscr R)_S.
\]
\begin{enumerate}
\item[\SGAInlineRefTarget{sga3:xxiv:1.19:item:i}(i)] The set of isomorphism classes of reductive \(S\)-groups
of type \(\mathscr R\) (\SGARefLink{sga3:xxii:def:2.7}{Exp.~XXII, 2.7}) is isomorphic, by
\SGARefLink{sga3:XXIII:5.12}{Exp.~XXIII, 5.12}, to
\[
  H^1(S,A_S(\mathscr R))
  =
  H^1_{\mathrm{\acute et}}(S,A_S(\mathscr R))
  =
  \Fib(S,A_S(\mathscr R)).
\]
\item[\SGAInlineRefTarget{sga3:xxiv:1.19:item:ii}(ii)] If \(S'\to S\) is a covering morphism, the subset
consisting of classes of groups splittable over \(S'\) is isomorphic
to \(H^1(S'/S,A_S(\mathscr R))\).
\end{enumerate}}

\SGARefTarget{sga3:xxiv:1.20}
\medskip\noindent\textbf{Remark 1.20.} ---
With the preceding notation, every reductive \(S\)-group \(G\) of
type \(\mathscr R\) has a canonically associated right principal
homogeneous bundle under \(A_S(\mathscr R)\):
\[
  \underline{\operatorname{Isom}}_{S\text{-}\mathrm{gr}}
    (\operatorname{\acute{E}p}_S(\mathscr R),G)=P.
\]
The scheme \(P\) may be interpreted as the ``scheme of pinnings of
\(G\) of type \(\mathscr R\)'' (cf.\ \SGARefLink{sga3:xxiv:1.0}{1.0}). It is also a principal
homogeneous bundle on the left under
\(\underline{\Aut}_{S\text{-}\mathrm{gr}}(G)\), a structure visible
immediately from the description above.

\SGARefTarget{sga3:xxiv:1.21}
\medskip\noindent\textbf{Proposition 1.21.} ---
{\itshape Let \(S\) be a henselian local scheme and \(s\) its closed
point. The functor
\[
  G\longmapsto G_s
\]
induces a bijection from the set of isomorphism classes of reductive
\(S\)-groups onto the set of isomorphism classes of reductive
\(\kappa(s)\)-groups. In particular, for every reductive \(S\)-group
\(G\), there exists a finite étale surjective morphism \(S'\to S\)
such that \(G_{S'}\) is splittable.}

% Authority local p. 10 / combined-reader p. 1255.
Using the existence of \(\operatorname{\acute{E}p}_S(\mathscr R)\)
(Exp.~XXV), it is enough to prove that, with
\(H=A_S(\mathscr R)\), the canonical map
\[
  \Fib(S,H)\longrightarrow\Fib(\kappa(s),H_s)
\]
is bijective, and that every element of \(\Fib(S,H)\) has the property
just stated. Every finite subset of \(H\) is contained in an affine
open subscheme: this is clear for a constant group, and \(H\) is
affine over a constant group. The result proved in the Appendix,
\SGARefLink{sga3:xxiv:proposition:8-1}{8.1}, therefore applies.

\SGARefTarget{sga3:xxiv:section:2}
\section*{2. Automorphisms and Subgroups}

Introduce the following notation. If
\[
  \mathcal H=\underline{\Aut}_{S\text{-}\mathrm{gr}}(G)
\]
and \(X\) is a subfunctor of \(G\), set
\[
  \underline{\Aut}_{S\text{-}\mathrm{gr}}(G,X)
  =\operatorname{Norm}_{\mathcal H}(X),
  \qquad
  \underline{\Aut}_{S\text{-}\mathrm{gr}}(G,\operatorname{id}_X)
  =\underline{\Centr}_{\mathcal H}(X).
\]
If \(Y\) is a second subfunctor of \(G\), define
\[
  \underline{\Aut}_{S\text{-}\mathrm{gr}}(G,X,Y)
  =
  \underline{\Aut}_{S\text{-}\mathrm{gr}}(G,X)
  \cap
  \underline{\Aut}_{S\text{-}\mathrm{gr}}(G,Y).
\]
If \(G'\) is a second \(S\)-group and \(X'\) a subfunctor of \(G'\),\footnote{Translator's note: The French text prints \(G\) here; \(G'\) is required by the typing of the definition.}
write
\[
  \underline{\operatorname{Isom}}_{S\text{-}\mathrm{gr}}
    (G,X;G',X')
\]
for the subfunctor of
\(\underline{\operatorname{Isom}}_{S\text{-}\mathrm{gr}}(G,G')\)
defined, for every \(S'\to S\), by
\[
  \left\{
    u\in
    \underline{\operatorname{Isom}}_{S\text{-}\mathrm{gr}}(G,G')(S')
    \,\middle|\,
    u(X_{S'})=X'_{S'}
  \right\}.
\]
The notation
\(\underline{\operatorname{Isom}}_{S\text{-}\mathrm{gr}}
(G,X,Y;G',X',Y')\), and so on, is defined similarly.\footnote{Editor's
note: The preceding definitions have been made explicit. The original
only said that
\(\underline{\Aut}_{S\text{-}\mathrm{gr}}(G,X,Y),\ldots\) and
\(\underline{\operatorname{Isom}}_{S\text{-}\mathrm{gr}}
(G,X;G',X'),\ldots\) were defined similarly.}

\SGARefTarget{sga3:xxiv:2.1}
\medskip\noindent\textbf{Proposition 2.1.} ---
{\itshape Let \(S\) be a scheme, \(G\) a reductive \(S\)-group, and
\(T\) a maximal torus of \(G\) (respectively, \(B\) a Borel subgroup
of \(G\), or \(B\supset T\) a Killing couple of \(G\)). Let
\(T_{\mathrm{ad}}\) (respectively, \(B_{\mathrm{ad}}\)) denote the
corresponding maximal torus (respectively, Borel subgroup) of
\(\operatorname{ad}(G)\):
\[
  B_{\mathrm{ad}}
  \simeq B/\underline{\Centr}(G)
  =B/\underline{\Centr}(B),
  \qquad
  T_{\mathrm{ad}}\simeq T/\underline{\Centr}(G).
\]
Then
\(\underline{\Aut}_{S\text{-}\mathrm{gr}}(G,T)\)
(respectively,
\(\underline{\Aut}_{S\text{-}\mathrm{gr}}(G,B)\), or
\(\underline{\Aut}_{S\text{-}\mathrm{gr}}(G,B,T)\))
is represented by a closed subscheme of
\(\underline{\Aut}_{S\text{-}\mathrm{gr}}(G)\), smooth over \(S\),
and the exact sequence of \SGARefLink{sga3:xxiv:1.3}{Theorem 1.3} induces the exact sequences
\[
  1\longrightarrow
  \operatorname{Norm}_{\operatorname{ad}(G)}(T_{\mathrm{ad}})
  \longrightarrow
  \underline{\Aut}_{S\text{-}\mathrm{gr}}(G,T)
  \longrightarrow
  \underline{\Autext}(G)
  \longrightarrow1,
\]
\[
  1\longrightarrow B_{\mathrm{ad}}
  \longrightarrow
  \underline{\Aut}_{S\text{-}\mathrm{gr}}(G,B)
  \longrightarrow
  \underline{\Autext}(G)
  \longrightarrow1,
\]
\[
  1\longrightarrow T_{\mathrm{ad}}
  \longrightarrow
  \underline{\Aut}_{S\text{-}\mathrm{gr}}(G,B,T)
  \longrightarrow
  \underline{\Autext}(G)
  \longrightarrow1.
\]}

By descent of closed subschemes,\footnote{Editor's note: Cf.\ SGA~1,
VIII, 1.9.} we reduce immediately to the case in which \(G\) is
pinned and \(B\supset T\) is its canonical Killing couple
(cf.\ \SGARefLink{sga3:xxii:cor:5.5.5:item:iv}{Exp.~XXII, 5.5.5(iv)}). Since the group \(E\) of \SGARefLink{sga3:xxiv:1.2}{1.2} normalizes
\(B\) and \(T\), the result follows at once from the normalization
theorems in \(\operatorname{ad}(G)\) (\SGARefLink{sga3:xxii:cor:5.3.12}{Exp.~XXII, 5.3.12} and 5.6.1).

% Authority local p. 11 / combined-reader p. 1256.

Using now the conjugacy theorems (cf.\ \SGARefLink{sga3:XXIII:5.12}{Exp.~XXIII, 5.12}), and arguing
as in \S1, one obtains:

\SGARefTarget{sga3:xxiv:2.2}
\medskip\noindent\textbf{Corollary 2.2.} ---
{\itshape Let \(S\) be a scheme, let \(G\) and \(G'\) be two reductive
\(S\)-groups having the same type at every point, and let
\(B\supset T\) (respectively, \(B'\supset T'\)) be a Killing couple
of \(G\) (respectively, \(G'\)).

\textup{\SGARefLink{sga3:xxiv:1.19:item:i}{(i)}} The \(S\)-functor
\[
  \underline{\operatorname{Isom}}_{S\text{-}\mathrm{gr}}
    (G,T;G',T')
\]
is represented by a closed smooth subscheme of
\(\underline{\operatorname{Isom}}_{S\text{-}\mathrm{gr}}(G,G')\)
which is a principal homogeneous space under
\(\underline{\Aut}_{S\text{-}\mathrm{gr}}(G,T)\).
Moreover,
\(\operatorname{Norm}_{\operatorname{ad}(G)}(T_{\mathrm{ad}})\)
acts freely on this scheme, and there is a canonical isomorphism
\[
  \underline{\operatorname{Isom}}_{S\text{-}\mathrm{gr}}
    (G,T;G',T')
  \big/
  \operatorname{Norm}_{\operatorname{ad}(G)}(T_{\mathrm{ad}})
  \simeq
  \underline{\operatorname{Isomext}}(G,G').
\]

\textup{\SGARefLink{sga3:xxiv:1.19:item:ii}{(ii)}} The \(S\)-functor
\[
  \underline{\operatorname{Isom}}_{S\text{-}\mathrm{gr}}
    (G,B;G',B')
\]
is represented by a closed smooth subscheme of
\(\underline{\operatorname{Isom}}_{S\text{-}\mathrm{gr}}(G,G')\)
which is a principal homogeneous space under
\(\underline{\Aut}_{S\text{-}\mathrm{gr}}(G,B)\).
Moreover, \(B_{\mathrm{ad}}\) acts freely on this scheme, and there
is a canonical isomorphism
\[
  \underline{\operatorname{Isom}}_{S\text{-}\mathrm{gr}}
    (G,B;G',B')
  \big/B_{\mathrm{ad}}
  \simeq
  \underline{\operatorname{Isomext}}(G,G').
\]

\textup{\SGARefLink{sga3:xxiv:1.17:item:iii}{(iii)}} The \(S\)-functor
\[
  \underline{\operatorname{Isom}}_{S\text{-}\mathrm{gr}}
    (G,B,T;G',B',T')
\]
is represented by a closed smooth subscheme of
\(\underline{\operatorname{Isom}}_{S\text{-}\mathrm{gr}}(G,G')\),
which is a principal homogeneous space under
\(\underline{\Aut}_{S\text{-}\mathrm{gr}}(G,B,T)\).
Moreover, \(T_{\mathrm{ad}}\) acts freely on this scheme, and there
is a canonical isomorphism
\[
  \underline{\operatorname{Isom}}_{S\text{-}\mathrm{gr}}
    (G,B,T;G',B',T')
  \big/T_{\mathrm{ad}}
  \simeq
  \underline{\operatorname{Isomext}}(G,G').
\]}

Arguing once again as in \S1, one obtains:

\SGARefTarget{sga3:xxiv:2.3}
\medskip\noindent\textbf{Corollary 2.3.} ---
{\itshape Let \(S\) be a scheme, let \(G\) be a reductive \(S\)-group,
and let \(B\supset T\) be a Killing couple of \(G\). The functor}
\[
  (G',T')
  \longmapsto
  \underline{\operatorname{Isom}}_{S\text{-}\mathrm{gr}}
    (G,T;G',T'),
\]
{\itshape respectively}
\[
  (G',B')
  \longmapsto
  \underline{\operatorname{Isom}}_{S\text{-}\mathrm{gr}}
    (G,B;G',B'),
\]
{\itshape respectively}
\[
  (G',B',T')
  \longmapsto
  \underline{\operatorname{Isom}}_{S\text{-}\mathrm{gr}}
    (G,B,T;G',B',T'),
\]
{\itshape is an equivalence between the category of pairs
\((G',T')\) (respectively, pairs \((G',B')\), respectively, triples
\((G',B',T')\)), where \(G'\) is a form of \(G\) and \(T'\) is a
maximal torus of \(G'\) (respectively, \(B'\) is a Borel subgroup of
\(G'\), respectively, \(B'\supset T'\) is a Killing couple of
\(G'\)), and the category of principal homogeneous bundles under
the \(S\)-group \(H\), where
\(H=\underline{\Aut}_{S\text{-}\mathrm{gr}}(G,T)\)
(respectively,
\(H=\underline{\Aut}_{S\text{-}\mathrm{gr}}(G,B)\),
respectively,
\(H=\underline{\Aut}_{S\text{-}\mathrm{gr}}(G,B,T)\)).
Moreover, every principal homogeneous sheaf under \(H\) is
representable and quasi-isotrivial, so that}
\[
  H^1(S,H)=H^1_{\mathrm{\acute et}}(S,H)=\Fib(S,H).
\]

% Authority local p. 12 / combined-reader p. 1257.

\SGARefTarget{sga3:xxiv:2.4}
\medskip\noindent\textbf{Remark 2.4.} ---
Under the hypotheses of \SGARefLink{sga3:xxiv:2.2}{2.2}, the morphism denoted set-theoretically
by \(u\mapsto u(T)\) (respectively, \(u\mapsto u(B)\), respectively,
\(u\mapsto(u(B),u(T))\)) induces an isomorphism
\[
  \underline{\operatorname{Isom}}_{S\text{-}\mathrm{gr}}(G,G')
  \big/
  \underline{\Aut}_{S\text{-}\mathrm{gr}}(G,T)
  \simeq \operatorname{Tor}(G'),
\]
respectively,
\[
  \underline{\operatorname{Isom}}_{S\text{-}\mathrm{gr}}(G,G')
  \big/
  \underline{\Aut}_{S\text{-}\mathrm{gr}}(G,B)
  \simeq \operatorname{Bor}(G'),
\]
respectively,
\[
  \underline{\operatorname{Isom}}_{S\text{-}\mathrm{gr}}(G,G')
  \big/
  \underline{\Aut}_{S\text{-}\mathrm{gr}}(G,B,T)
  \simeq \operatorname{Kil}(G').
\]
The proof is immediate: it suffices to prove the assertion locally
for \textup{(fpqc)}, so one may assume \(G\simeq G'\), and the result
then reduces to \SGARefLink{sga3:xxii:cor:5.8.3:item:iii}{Exp.~XXII, 5.8.3(iii)}.

\SGARefTarget{sga3:xxiv:2.5}
\medskip\noindent\textbf{Remark 2.5.} ---
The preceding results are immediately interpreted in terms of
reduction of the structure group. If \(G'\) is a form of \(G\),
corresponding to the principal bundle
\(\underline{\operatorname{Isom}}_{S\text{-}\mathrm{gr}}(G,G')\),
to give a reduction of the structure group of this bundle to
\(\underline{\Aut}_{S\text{-}\mathrm{gr}}(G,T)\) amounts to giving
a maximal torus \(T'\) of \(G'\). The indicated bijections are,
on the one hand, that of \SGARefLink{sga3:xxiv:2.4}{2.4} and, on the other hand, the map
\[
  T'\longmapsto
  \underline{\operatorname{Isom}}_{S\text{-}\mathrm{gr}}
    (G,T;G',T').
\]
The same applies to Borel subgroups and Killing couples.

\SGARefTarget{sga3:xxiv:2.6}
\medskip\noindent\textbf{Proposition 2.6.} ---
{\itshape Let \(S\) be a scheme, let \(G\) and \(G'\) be two reductive
\(S\)-groups having the same type at every point, and let \(T\)
(respectively, \(T'\)) be a maximal torus of \(G\) (respectively,
\(G'\)). Then \(T_{\mathrm{ad}}\) acts freely on
\(\underline{\operatorname{Isom}}_{S\text{-}\mathrm{gr}}
(G,T;G',T')\), and the quotient}
\[
  P=
  \underline{\operatorname{Isom}}_{S\text{-}\mathrm{gr}}
    (G,T;G',T')
  \big/T_{\mathrm{ad}}
\]
{\itshape is representable; it is a principal homogeneous bundle under}
\[
  A=
  \underline{\Aut}_{S\text{-}\mathrm{gr}}(G,T)
  \big/T_{\mathrm{ad}},
\]
{\itshape where \(A\) is represented by a twisted constant \(S\)-scheme
and is an extension of \(\underline{\Autext}(G)\) by}
\[
  W_{\operatorname{ad}(G)}(T_{\mathrm{ad}})
  =
  \operatorname{Norm}_{\operatorname{ad}(G)}(T_{\mathrm{ad}})
  \big/T_{\mathrm{ad}}.
\]
{\itshape Moreover, if \(A\) is made to act on \(T\) in the evident
way, the bundle associated with \(P\) is precisely \(T'\).}\footnote{Editor's
note: Consequently, if \(P\) has a section, then \(T\) and \(T'\)
are isomorphic; in particular, \(T_P\) and \(T'_P\) are isomorphic.
See the addition \SGARefLink{sga3:xxiv:proposition:4-5-1}{4.5.1}, where this observation is developed in the
case \(G=\operatorname{\acute{E}p}_S(\mathscr R)\).}

The first part of the proposition follows immediately from the
preceding results. To prove the second, observe that there is an
evident morphism
\[
  P\times_S T\longrightarrow T',
  \qquad (u,t)\longmapsto u(t).
\]
To show that, after passage to the quotient by \(A\), it induces an
isomorphism, one may once again assume
\((G,T)\simeq(G',T')\), in which case the assertion is immediate.

In an entirely analogous way, one has:

\SGARefTarget{sga3:xxiv:2.7}
\medskip\noindent\textbf{Proposition 2.7.} ---
{\itshape Let \(S\) be a scheme, let \(G\) and \(G'\) be two reductive
\(S\)-groups having the same type at every point, and let
\(B\supset T\) (respectively, \(B'\supset T'\)) be a Killing couple
of \(G\) (respectively, \(G'\)). If
\[
  \underline{\Aut}_{S\text{-}\mathrm{gr}}(G,B,T)
  \big/T_{\mathrm{ad}}
  \simeq \underline{\Autext}(G)
\]
is made to act on \(T\) in the evident way, the bundle associated
with \(\underline{\operatorname{Isomext}}(G,G')\) is precisely
\(T'\).}

\SGARefTarget{sga3:xxiv:2.8}
\medskip\noindent\textbf{Corollary 2.8.} ---
{\itshape Let \(G\) and \(G'\) be two reductive \(S\)-groups which are
inner twisted forms of one another, and let \(B\supset T\)
(respectively, \(B'\supset T'\)) be a Killing couple of \(G\)
(respectively, \(G'\)). Then \(T\) and \(T'\) are
isomorphic.}\footnote{Editor's note: This follows from \SGARefLink{sga3:xxiv:2.7}{2.7}, since
\(\underline{\operatorname{Isomext}}(G,G')\) has a section
(cf.\ the preceding editor's note).}


% END inputs/XXIV/tex/components/06_expose_XXIV_corollary115_through_proposition28_en.tex

% BEGIN inputs/XXIV/tex/components/07_expose_XXIV_section2_conclusion_and_section3_through_39_en.tex
% Exposé XXIV, authority local pages 13--16 / combined-reader pages 1258--1261.
% Authority local p. 13 / combined-reader p. 1258.

\SGARefTarget{sga3:xxiv:2.9}
\medskip\noindent\textbf{Remark 2.9.} ---
It is not true in general that \(B\) and \(B'\) are isomorphic;
they are nevertheless inner twisted forms of one another
(cf.\ \S5).

One may develop ``\(\operatorname{Isomint}\)'' variants\footnote{Editor's
note: Recall that
\(\underline{\operatorname{Isomint}}_u(G,G')\) is defined in \SGARefLink{sga3:xxiv:1.11}{1.11}.}
of the preceding results. We mention one:

\SGARefTarget{sga3:xxiv:2.10}
\medskip\noindent\textbf{Proposition 2.10.} ---
{\itshape Let \(S\) be a scheme, let \(G\) and \(G'\) be two reductive
\(S\)-groups having the same type at every point, let \(B\supset T\)
(respectively, \(B'\supset T'\)) be a Killing couple of \(G\)
(respectively, \(G'\)), and let
\(u\in\underline{\operatorname{Isomext}}(G,G')(S)\). Consider}
\[
\begin{aligned}
  &\underline{\operatorname{Isomint}}_u
     (G,B,T;G',B',T')\\
  &\qquad =
  \underline{\operatorname{Isom}}_{S\text{-}\mathrm{gr}}
     (G,B,T;G',B',T')
  \cap
  \underline{\operatorname{Isomint}}_u(G,G').
\end{aligned}
\]
{\itshape It is a smooth affine \(S\)-scheme which is a principal
homogeneous bundle under \(T_{\mathrm{ad}}\). In particular,}
\[
  \underline{\operatorname{Isomint}}_u
    (G,B,T;G',B',T')(S)\ne\varnothing
\]
{\itshape if and only if the corresponding element of
\(H^1(S,T_{\mathrm{ad}})\) is zero.}

To conclude this section, we prove two results that will be useful
later.

\SGARefTarget{sga3:xxiv:2.11}
\medskip\noindent\textbf{Proposition 2.11.} ---
{\itshape Let \(S\) be a scheme, let \(G\) be a reductive \(S\)-group,
and let \(T\) be a maximal torus of \(G\). The evident morphism}
\[
  T_{\mathrm{ad}}
  \xrightarrow{\ \sim\ }
  \underline{\Aut}_{S\text{-}\mathrm{gr}}(G,\operatorname{id}_T)
\]
{\itshape is an isomorphism.}

This follows from the preceding statements. A direct proof was also
given in the course of the proof of \SGARefLink{sga3:xxiv:1.5.2}{1.5.2}.

\SGARefTarget{sga3:xxiv:2.12}
\medskip\noindent\textbf{Corollary 2.12.} ---
{\itshape Under the preceding hypotheses, there is an equivalence
between the category of pairs \((G',f)\), where \(G'\) is a form of
\(G\) and \(f\) is an isomorphism from \(T\) onto a maximal torus of
\(G'\), and the category of principal homogeneous bundles under
\(T_{\mathrm{ad}}\).}

\SGARefTarget{sga3:xxiv:2.13}
\medskip\noindent\textbf{Corollary 2.13.} ---
{\itshape If \(H^1(S,T_{\mathrm{ad}})=0\), and if \(G'\) is a form of
\(G\) having a maximal torus isomorphic to \(T\), then \(G'\) is
isomorphic to \(G\).}

\SGARefTarget{sga3:xxiv:2.14}
\medskip\noindent\textbf{Corollary 2.14.} ---
{\itshape Let \(S\) be a scheme such that \(\operatorname{Pic}(S)=0\),
and let \(G\) be a reductive \(S\)-group of constant type. Then \(G\)
is splittable if and only if \(G\) has a split maximal torus.}

Let \(G\) be a reductive group and let
\(\operatorname{rad}(G)\) be its radical (\SGARefLink{sga3:xxii:rem:4.3.9}{Exp.~XXII, 4.3.9}).
Since \(\operatorname{rad}(G)\) is central and characteristic in
\(G\), there is a canonical morphism
\[
  q:\underline{\Autext}(G)
  \longrightarrow
  \underline{\Aut}_{S\text{-}\mathrm{gr}}
    \bigl(\operatorname{rad}(G)\bigr).
\]

% Authority local p. 14 / combined-reader p. 1259.

\SGARefTarget{sga3:xxiv:2.15}
\medskip\noindent\textbf{Proposition 2.15.} ---
{\itshape Let \(G\) be a reductive \(S\)-group. The following sequence
is exact:}
\[
  1\longrightarrow
  \operatorname{ad}(G)
  \longrightarrow
  \underline{\Aut}_{S\text{-}\mathrm{gr}}
    \bigl(G,\operatorname{id}_{\operatorname{rad}(G)}\bigr)
  \xrightarrow{p}
  \underline{\Autext}(G)
  \xrightarrow{q}
  \underline{\Aut}_{S\text{-}\mathrm{gr}}
    \bigl(\operatorname{rad}(G)\bigr),
\]
{\itshape and \(\ker(q)=\operatorname{im}(p)\) is an open and closed
subscheme of \(\underline{\Autext}(G)\), finite over \(S\).}

One may assume that \(G\) is split. The first assertion is immediate,
and the second follows from \SGARefLink{sga3:XXI:6.7.5}{Exp.~XXI, 6.7.5} and 6.7.7.

Writing
\[
  H=
  \underline{\Aut}_{S\text{-}\mathrm{gr}}
    \bigl(G,\operatorname{id}_{\operatorname{rad}(G)}\bigr),
\]
one obtains:

\SGARefTarget{sga3:xxiv:2.16}
\medskip\noindent\textbf{Corollary 2.16.} ---
{\itshape There is an equivalence between the category of pairs
\((G',f)\), where \(G'\) is a form of \(G\) and \(f\) is an
isomorphism from \(\operatorname{rad}(G)\) onto the radical of \(G'\),
and the category of principal bundles under a certain \(S\)-group
scheme \(H\), where \(H\) is such that there is an exact sequence}
\[
  1\longrightarrow \operatorname{ad}(G)
  \longrightarrow H
  \longrightarrow F
  \longrightarrow 1,
\]
{\itshape in which the \(S\)-group \(F\) is étale and finite over
\(S\).}

\SGARefTarget{sga3:xxiv:section:3}
\section*{3. Dynkin Scheme of a Reductive Group. Quasi-Split Groups}

\SGARefTarget{sga3:xxiv:3.1}
\medskip\noindent\textbf{3.1.} ---
Recall (\SGARefLink{sga3:XXI:7.4.1}{Exp.~XXI, 7.4.1}) that a Dynkin diagram is a finite set endowed
with the structure defined by a set of pairs of distinct elements
(bonds) and a map to \(\{1,2,3\}\) (lengths). To every pinned reduced
root datum \(\mathscr R\) there is associated a Dynkin diagram
\(\Delta(\mathscr R)\), whose underlying set is the set of simple
roots.

\SGARefTarget{sga3:xxiv:3.2}
\medskip\noindent\textbf{3.2.} ---
Let \(S\) be a scheme. A Dynkin \(S\)-scheme is a finite twisted
constant \(S\)-scheme \(\Delta\), endowed with the structure defined
by a subscheme \(L\) of \(\Delta\times_S\Delta\) having empty
intersection with the diagonal, and by a morphism
\[
  \Delta\longrightarrow\{1,2,3\}_S.
\]
If \(S'\) is a connected \(S\)-scheme,\footnote{Editor's note: The
hypothesis that \(S'\) be connected has been added.}
\(\Delta(S')\) is naturally endowed with the structure of a Dynkin
diagram.

The following notions are defined immediately: an isomorphism of two
Dynkin schemes, the base extension of a Dynkin scheme, and the
constant Dynkin scheme associated with a Dynkin diagram.

Every étale descent datum on a Dynkin scheme is effective.

\SGARefTarget{sga3:xxiv:3.3}
\medskip\noindent\textbf{3.3.} ---
We shall associate a Dynkin \(S\)-scheme with every reductive
\(S\)-group \(G\). First suppose that \(G\) is splittable over \(S\).
For every pinning \(E\) of \(G\), let \(\Delta(E)\) be the constant
Dynkin scheme associated with the pinned root datum defined by \(E\).
If \(E\) and \(E'\) are two pinnings of \(G\), then by \SGARefLink{sga3:xxiv:1.5}{1.5} there is
a unique inner automorphism of \(G\) over \(S\) carrying \(E\) to
\(E'\). This automorphism of \(G\) defines an isomorphism
\[
  a_{EE'}:\Delta(E)\xrightarrow{\ \sim\ }\Delta(E').
\]
The \(a_{EE'}\) plainly form a transitive system, so the
\(\Delta(E)\) may be identified (that is, one takes their inductive
limit); the result is a constant Dynkin scheme denoted by
\(\operatorname{Dyn}(G)\). If now \(G\) is an arbitrary reductive
\(S\)-group, there is an étale covering family \(\{S_i\to S\}\)
such that

% Authority local p. 15 / combined-reader p. 1260.

\(G_{S_i}\) is splittable. Arguing as above, one therefore obtains a
canonical descent datum on the \(\operatorname{Dyn}(G_{S_i})\), and
by descent this constructs a Dynkin \(S\)-scheme
\(\operatorname{Dyn}(G)\).

\SGARefTarget{sga3:xxiv:3.4}
\medskip\noindent\textbf{3.4.} ---
This construction has the following properties (which, moreover,
essentially characterize it):
\begin{enumerate}
\item[\SGAInlineRefTarget{sga3:xxiv:3.4:item:i}(i)] With every reductive \(S\)-group \(G\) is associated a
Dynkin scheme \(\operatorname{Dyn}(G)\); with every isomorphism
\[
  u:G\xrightarrow{\ \sim\ }G'
\]
is functorially associated an isomorphism
\[
  \operatorname{Dyn}(u):
  \operatorname{Dyn}(G)\xrightarrow{\ \sim\ }\operatorname{Dyn}(G').
\]

\item[\SGAInlineRefTarget{sga3:xxiv:3.4:item:ii}(ii)] If \(S'\) is an \(S\)-scheme and \(G\) is a reductive
\(S\)-group, then
\[
  \operatorname{Dyn}(G\times_S S')
  \simeq
  \operatorname{Dyn}(G)\times_S S'.
\]

\item[(iii)] If \(E\) is a pinning of \(G\) defining the pinned root
datum with Dynkin diagram \(\Delta\), then
\[
  \operatorname{Dyn}(G)\simeq\Delta_S.
\]

\item[\SGAInlineRefTarget{sga3:xxiv:3.4:item:iv}(iv)] If \(u\) is an inner automorphism of \(G\), then
\(\operatorname{Dyn}(u)\) is the identity automorphism of
\(\operatorname{Dyn}(G)\).
\end{enumerate}

\SGARefTarget{sga3:xxiv:3.5}
\medskip\noindent\textbf{3.5.} ---
Let \(S\) be a scheme and \(G\) a reductive \(S\)-group. It is clear
that the functor
\[
  \underline{\Aut}_{\mathrm{Dyn}}
    \bigl(\operatorname{Dyn}(G)\bigr)
\]
of automorphisms of \(\operatorname{Dyn}(G)\) as a Dynkin scheme is
represented by a finite twisted constant \(S\)-scheme. By \SGARefLink{sga3:xxiv:3.4:item:i}{3.4(i)}
and \SGARefLink{sga3:xxiv:3.4:item:ii}{(ii)}, there is a canonical morphism
\[
  \underline{\Aut}_{S\text{-}\mathrm{gr}}(G)
  \longrightarrow
  \underline{\Aut}_{\mathrm{Dyn}}
    \bigl(\operatorname{Dyn}(G)\bigr),
\]
which, by \SGARefLink{sga3:xxiv:3.4:item:iv}{(iv)}, factors through a morphism
\[
  \underline{\Autext}(G)
  \longrightarrow
  \underline{\Aut}_{\mathrm{Dyn}}
    \bigl(\operatorname{Dyn}(G)\bigr).
\]
More generally, if \(G\) and \(G'\) are two reductive \(S\)-groups,
there is a canonical morphism
\[
  \underline{\operatorname{Isomext}}(G,G')
  \longrightarrow
  \underline{\operatorname{Isom}}_{\mathrm{Dyn}}
    \bigl(\operatorname{Dyn}(G),\operatorname{Dyn}(G')\bigr).
\]
In particular, if \(G'\) is an inner twisted form of \(G\) (\SGARefLink{sga3:xxiv:1.11}{1.11}),
the Dynkin schemes \(\operatorname{Dyn}(G)\) and
\(\operatorname{Dyn}(G')\) are isomorphic.

\SGARefTarget{sga3:xxiv:3.6}
\medskip\noindent\textbf{3.6.} ---
If \(G\) is semisimple (respectively, adjoint or simply connected),
the morphism
\[
  \underline{\Autext}(G)
  \longrightarrow
  \underline{\Aut}_{\mathrm{Dyn}}
    \bigl(\operatorname{Dyn}(G)\bigr)
\]
is a monomorphism (respectively, an isomorphism). Indeed, one may
assume that \(G\) is pinned, and the assertion reduces to the
corresponding result for pinned reduced root data
(cf.\ \SGARefLink{sga3:XXI:7.4.5}{Exp.~XXI, 7.4.5}).

There is an analogous result for the \(\operatorname{Isom}\) functors.
It follows in particular that two semisimple adjoint (respectively,
simply connected) \(S\)-groups are inner twisted forms of one another
if and only if their Dynkin schemes are isomorphic.

% Authority local p. 16 / combined-reader p. 1261.

\SGARefTarget{sga3:xxiv:3.7}
\medskip\noindent\textbf{3.7.} ---
One can give a different construction of the Dynkin scheme associated
with a reductive group. Let \(\mathscr R\) be a pinned reduced root
datum and \(G\) a reductive \(S\)-group of type \(\mathscr R\), and
write \(\Delta(\mathscr R)\) for the Dynkin diagram defined by
\(\mathscr R\). There is, by \SGARefLink{sga3:xxiv:3.5}{3.5}, a canonical morphism
\[
  A_S(\mathscr R)
  =
  \underline{\Aut}_{S\text{-}\mathrm{gr}}
    \bigl(\operatorname{\acute{E}p}_S(\mathscr R)\bigr)
  \longrightarrow
  \underline{\Aut}_{\mathrm{Dyn}}
    \bigl(\Delta(\mathscr R)_S\bigr).
\]
The reductive \(S\)-group \(G\) corresponds (\SGARefLink{sga3:xxiv:1.17}{1.17}) to a bundle
\[
  \underline{\operatorname{Isom}}_{S\text{-}\mathrm{gr}}
    \bigl(\operatorname{\acute{E}p}_S(\mathscr R),G\bigr),
\]
principal homogeneous under \(A_S(\mathscr R)\). The associated
bundle under
\(\underline{\Aut}_{\mathrm{Dyn}}(\Delta(\mathscr R)_S)\)
corresponds to a form over \(S\) of \(\Delta(\mathscr R)_S\): it is
\(\operatorname{Dyn}(G)\). In other words, this associated bundle is
precisely
\[
  \underline{\operatorname{Isom}}_{\mathrm{Dyn}}
    \bigl(\Delta(\mathscr R)_S,\operatorname{Dyn}(G)\bigr).
\]
In this form the proof is immediate.

\SGARefTarget{sga3:xxiv:3.8}
\medskip\noindent\textbf{3.8. Dynkin Scheme and Killing Couples.} ---
Let \(S\) be a scheme, \(G\) a reductive \(S\)-group, and
\(B\supset T\) a Killing couple of \(G\). There is a canonical
morphism
\[
  i:\operatorname{Dyn}(G)
  \longrightarrow
  \underline{\operatorname{Hom}}_{S\text{-}\mathrm{gr}}
    (T,\mathbb G_{m,S})
\]
which identifies \(\operatorname{Dyn}(G)\) with the ``scheme of
simple roots of \(B\) relative to \(T\).'' This morphism is defined
immediately by descent from the pinned case. Observe also that the
data of \(T\) and \(i\) allow one to reconstruct \(B\) (the
``one-to-one correspondence between systems of simple roots and
systems of positive roots'').

It follows from the preceding description of
\(D=\operatorname{Dyn}(G)\) that there is a canonical root of \(B_D\)
relative to \(T_D\): this root \(\alpha_D\) is the image, under
\(i(D)\), of the identity morphism of \(D\). One obtains from it a
canonical invertible \(\mathcal O_D\)-module \(\mathfrak g_D\):
\[
  \mathfrak g_D
  =
  \bigl(\mathfrak g\otimes_{\mathcal O_S}\mathcal O_D\bigr)^{\alpha_D}.
\]
In the pinned case,
\[
  D=\coprod_{\alpha\in\Delta}S_\alpha,
\]
where every \(S_\alpha\) is a copy of \(S\), and
\(\mathfrak g_D\) is then the \(\mathcal O_D\)-module inducing
\(\mathfrak g_\alpha\) on \(S_\alpha\), for every
\(\alpha\in\Delta\).

\SGARefTarget{sga3:xxiv:3.9}
\medskip\noindent\textbf{3.9. Quasi-Pinnings. Quasi-Pinned Groups.} ---
If \(G\) is a reductive \(S\)-group, a \emph{quasi-pinning of \(G\)}
means the data
\begin{enumerate}
\item[(i)] of a Killing couple \(B\supset T\) of \(G\);
\item[(ii)] of a section
\[
  X\in\Gamma\bigl(\operatorname{Dyn}(G),\mathfrak g_D\bigr)^\times.
\]
\end{enumerate}
A reductive \(S\)-group is called \emph{quasi-splittable} if it has a
quasi-pinning. A \emph{quasi-pinned group} is a reductive group
equipped with a quasi-pinning.

Let \(B\supset T\) be a Killing couple of the reductive \(S\)-group
\(G\). Then \(G\) is quasi-pinnable relative to this Killing couple
if and only if \(\mathfrak g_D\) has a section nonzero at every point,
that is, if the element of \(\operatorname{Pic}(\operatorname{Dyn}(G))\)
defined by \(\mathfrak g_D\) is zero. Suppose in particular that
\(S\) is semilocal. Then \(\operatorname{Dyn}(G)\) is also semilocal,
and hence \(\operatorname{Pic}(\operatorname{Dyn}(G))=0\). Thus:

\SGARefTarget{sga3:xxiv:3.9.1}
\medskip\noindent\textbf{Proposition 3.9.1.} ---
{\itshape Let \(S\) be a semilocal scheme and \(G\) a reductive
\(S\)-group. Then \(G\) is quasi-splittable if and only if it has a
Borel subgroup.}


% END inputs/XXIV/tex/components/07_expose_XXIV_section2_conclusion_and_section3_through_39_en.tex

% BEGIN inputs/XXIV/tex/components/08_expose_XXIV_39_proof_through_corollary312_en.tex
% Exposé XXIV, authority local pages 17--20 / combined-reader pages 1262--1265.
% Authority local p. 17 / combined-reader p. 1262.

\footnote{Editor's note: The reference to Exp.~XXII, 5.9.7 has been
given in detail.}
Indeed, \(S\) is affine, so by the first assertion of
\SGARefLink{sga3:xxii:cor:5.9.7}{Exp.~XXII, 5.9.7}, if \(G\) has a Borel subgroup \(B\), it also has a
Killing couple \(B\supset T\). Then, because
\(\operatorname{Pic}(D)=0\), the module \(\mathfrak g_D\) has a
nowhere-zero section \(X\).

Continue to let \(A\) be a semilocal ring and
\(S=\operatorname{Spec}(A)\). For every reductive \(S\)-group \(G\),
the morphism
\[
  \operatorname{Bor}(G)\longrightarrow S
\]
is surjective (because \(G_s\) has Borel subgroups for every
\(s\in S\)), smooth, and projective (\SGARefLink{sga3:xxii:cor:5.8.3}{Exp.~XXII, 5.8.3}). It therefore
has sections after a finite, surjective, étale extension of the
base.\footnote{Editor's note: The original referred to EGA~IV,
\S24, which did not appear. The point is to use ``Bertini's
theorem.'' We give the details of the argument, which were indicated
to us by O.\ Gabber. Let \(X\to S\) be a surjective smooth projective
morphism. Replacing \(S\) by a connected component, one may assume
that \(X/S\) has relative dimension \(d\) at every point
(cf.\ EGA~IV\textsubscript{4}, 17.10.2). One may also assume that
\(X\) is a closed subscheme of a projective space
\(\mathbb P^n_S=\mathbb P(E)\), where \(E\) is a free \(A\)-module
of rank \(n+1\) (cf.\ EGA~II, 5.3.3). Let
\(s_1,\ldots,s_r\) be the closed points of \(S\). By Bertini's
theorem (see, for example, \textup{[Jou83]}, I, 6.10), there is an
open subset \(U\) of
\[
  P=\mathbb P(E^\ast)^d
\]
with nonempty fibers such that, for every point
\(u=(f_1,\ldots,f_d)\) of \(U_{s_i}\), the intersection of
\(X_{\kappa(u)}\) with the \(d\) hyperplanes of
\(\mathbb P^n_{\kappa(u)}\) defined by \(u\) is étale over
\(\kappa(u)\). After shrinking \(U\), one may assume that \(U\) is
the complement in affine space \(\mathbb A^{nd}_S\) of the zero
locus \(V(Q)\) of a polynomial \(Q\) of degree \(m>0\). An easy
induction on the number of variables then shows that \(V(Q)\) cannot
contain every rational point of
\(\mathbb A^{nd}_{\kappa(s_i)}\) if
\(\lvert\kappa(s_i)\rvert>m\). Since the morphism
\[
  A\longrightarrow\prod_i\kappa(s_i)
\]
is surjective, one can find a monic polynomial \(R\in A[X]\) of
degree \(m\) whose image in \(\kappa(s_i)[X]\) is irreducible when
\(\kappa(s_i)\) is finite and has \(m\) distinct roots when
\(\kappa(s_i)\) is infinite. Put
\[
  A'=A[X]/(R),
  \qquad
  S'=\operatorname{Spec}(A').
\]
Then \(S'\to S\) is étale, finite, and surjective, and the residue
field at each of its closed points \(s'_1,\ldots,s'_t\) has
cardinality at least \(2m\).
Thus \(U\) has a rational point \(u_i\) above every closed point of
\(S'\), and, since
\[
  A'\longrightarrow\prod_i\kappa(s'_i)
\]
is surjective \emph{(Translator's note: the French text prints
\(\kappa(s_i)\); the intended closed points of \(S'\) are
\(s'_i\))}, these points lift to a section \(u\) of \(P_{S'}\).
Let \(Z\) be the intersection of \(X_{S'}\) with the \(d\)
hyperplanes of \(\mathbb P^n_{S'}\) defined by \(u\), and let \(V\)
be the open subset of \(Z\) consisting of the points at which \(Z\)
is étale over \(S'\). By EGA~IV\textsubscript{3}, 11.3.8, \(V\)
contains the fibers \(Z_{s'_i}\) for every \(i\). Since
\(\pi:Z\to S'\) is proper, it follows that the closed set
\(\pi(Z-V)\) is empty, whence \(V=Z\). Thus \(Z\to S'\) is
surjective, étale, and proper, hence finite, as is the composite
\(Z\to S\); this gives the desired section of \(X\to S\).}
It follows that:

\SGARefTarget{sga3:xxiv:3.9.2}
\medskip\noindent\textbf{Corollary 3.9.2.} ---
{\itshape Let \(S\) be a semilocal scheme and \(G\) a reductive
\(S\)-group. There is an étale, finite, and surjective morphism
\(S'\to S\) such that \(G_{S'}\) is quasi-splittable.}

\SGARefTarget{sga3:xxiv:3.9.3}
\medskip\noindent\textbf{Remark 3.9.3.} ---
Under the preceding hypotheses, let \(T\) be a maximal torus of
\(G\) (cf.\ \SGARefLink{sga3:xiv:3.20}{Exp.~XIV, 3.20}). One may moreover assume that \(G_{S'}\)
is quasi-splittable relative to \(T_{S'}\): it suffices to apply the
preceding argument to the ``scheme of Borel subgroups containing
\(T\),'' which is finite and étale over \(S\)
(\SGARefLink{sga3:xxii:corollary:5-5-5:item:ii}{Exp.~XXII, 5.5.5(ii)}).

\SGARefTarget{sga3:xxiv:3.10}
\medskip\noindent\textbf{3.10.} ---
Let \(E\) and \(E'\) be two quasi-pinnings of the reductive
\(S\)-group \(G\). There is a unique inner automorphism of \(G\)
carrying \(E\) to \(E'\). Indeed, one reduces immediately to the
split case, where the assertion has already been proved (\SGARefLink{sga3:xxiv:1.5}{1.5}; it
suffices to observe that for an inner automorphism of \(G\), preserving
a pinning is equivalent to preserving the underlying quasi-pinning).
One concludes, as in \S1, that a

% Authority local p. 18 / combined-reader p. 1263.

quasi-pinning of the reductive \(S\)-group \(G\) defines a splitting
\(h\) of the exact sequence:

\[
  % BEGIN inputs/XXIV/native_diagrams/exp24/001_split_autext_exact_sequence.tex
\begin{tikzcd}[column sep=small]
1 \arrow[r]
&
\underline{\operatorname{ad}}(G) \arrow[r]
&
\underline{\Aut}_{S\text{-}\mathrm{gr.}}(G)
  \arrow[r,shift right=.7ex,"p"']
&
\underline{\Autext}(G)
  \arrow[l,shift right=.7ex,dotted,"h"']
  \arrow[r]
&
1
\end{tikzcd}

% END inputs/XXIV/native_diagrams/exp24/001_split_autext_exact_sequence.tex

\]

the image of \(h\) being the subgroup of
\(\underline{\Aut}_{S\text{-}\mathrm{gr}}(G)\) that preserves the
quasi-pinning.

Likewise, if \(G\) and \(G'\) are two quasi-pinned \(S\)-groups, one
defines naturally the subfunctor
\[
  \underline{\operatorname{Isom}}_{S\text{-}\mathrm{gr},\,q\text{-}\acute{e}p}
    (G,G')
\]
of
\(\underline{\operatorname{Isom}}_{S\text{-}\mathrm{gr}}(G,G')\).
The projection from
\(\underline{\operatorname{Isom}}_{S\text{-}\mathrm{gr}}(G,G')\)
to \(\underline{\operatorname{Isomext}}(G,G')\) induces an
isomorphism
\[
  (*)\qquad
  \underline{\operatorname{Isom}}_{S\text{-}\mathrm{gr},\,q\text{-}\acute{e}p}
    (G,G')
  \xrightarrow{\ \sim\ }
  \underline{\operatorname{Isomext}}(G,G').
\]

\SGARefTarget{sga3:xxiv:3.11}
\medskip\noindent\textbf{Theorem 3.11.} ---
{\itshape Let \(S\) be a scheme, let \(\mathscr R\) be a pinned reduced
root datum such that
\(\operatorname{\acute{E}p}_S(\mathscr R)\) exists
(cf.\ Exp.~XXV), and let \(E\) be its group of automorphisms.
Consider the following three categories:
\begin{enumerate}
\item[\SGAInlineRefTarget{sga3:xxiv:3.11:item:i}(i)] the category \(\mathsf{Rev}\) of principal Galois coverings
of \(S\) with group \(E\) (the morphisms are isomorphisms of
\(E\)-bundles);
\item[\SGAInlineRefTarget{sga3:xxiv:3.11:item:ii}(ii)] the category \(\mathsf{Redext}\), whose objects are the
reductive \(S\)-groups of type \(\mathscr R\) (\SGARefLink{sga3:xxii:def:2.7}{Exp.~XXII, 2.7}), the
morphisms from \(G\) to \(G'\) being the elements of
\(\underline{\operatorname{Isomext}}(G,G')(S)\);
\item[(iii)] the category \(\mathsf{Q\acute{e}p}\) of quasi-pinned
reductive \(S\)-groups of type \(\mathscr R\) (the morphisms are the
isomorphisms preserving the quasi-pinnings).
\end{enumerate}
These three categories are equivalent. More precisely, there is a
diagram of functors, commutative up to functorial
isomorphisms:}

\[
  % BEGIN inputs/XXIV/native_diagrams/exp24/002_rev_qep_redext_triangle.tex
\begin{tikzcd}[column sep=large,row sep=large]
\mathsf{Rev}
  \arrow[rr,"q\acute{e}p"]
&&
\mathsf{Q\acute{e}p}
  \arrow[dl,"i"]
\\
&
\mathsf{Redext}
  \arrow[ul,"\operatorname{rev}"]
&
\end{tikzcd}

% END inputs/XXIV/native_diagrams/exp24/002_rev_qep_redext_triangle.tex

\]

We describe these three functors below, leaving to the reader the
verification that the diagram commutes.\footnote{Editor's note: In
\SGARefLink{sga3:xxiv:3.11.4}{3.11.4} we have sketched the verification that the functors \(i\),
\(\operatorname{rev}\), and \(q\acute{e}p\) are fully faithful and
that the composite
\(\operatorname{rev}\circ i\circ q\acute{e}p\) is isomorphic to the
identity functor of \(\mathsf{Rev}\).}

\SGARefTarget{sga3:xxiv:3.11.1}
\medskip\noindent\textbf{3.11.1. The functor \(i\).} ---
This is the evident functor:
\[
  i(G)=G,
  \qquad
  i(f)=\text{the image of \(f\) under the isomorphism 3.10 \((*)\)}.
\]

\SGARefTarget{sga3:xxiv:3.11.2}
\medskip\noindent\textbf{3.11.2. The functor
\(q\acute{e}p\).} ---
Let \(S'\) be a principal Galois covering of \(S\) with group \(E\).
Let
\[
  \bigl(G,T,M,\mathscr R,\Delta,
    (X_\alpha)_{\alpha\in\Delta}\bigr)
\]
be a pinning of
\(G=\operatorname{\acute{E}p}_S(\mathscr R)\), and let
\((X'_\alpha)_{\alpha\in\Delta}\) be the corresponding pinning of
\[
  G'=G_{S'}=\operatorname{\acute{E}p}_{S'}(\mathscr R).
\]
By \SGARefLink{sga3:XXIII:5.5bis}{Exp.~XXIII, 5.5 bis}, there is a unique group homomorphism
\[
  \theta:E
  \longrightarrow
  \underline{\Aut}_{S\text{-}\mathrm{gr}}
    \bigl(\operatorname{\acute{E}p}_S(\mathscr R)\bigr)
  =A(\mathscr R)(S)
\]

% Authority local p. 19 / combined-reader p. 1264.

such that, for every \(h\in E\), the automorphism \(\theta(h)\)
induces \(D_S(h)\) on \(T\), and
\[
  \Lie(\theta(h))(X_\alpha)=X_{h(\alpha)}
  \qquad(\alpha\in\Delta).
\]
Taking into account the action of \(E\) on \(S'\), this yields an
action of \(E\) on \(G'\), compatible with the action of \(E\) on
\(S'\) and preserving the canonical quasi-pinning of
\(G'\).\footnote{Editor's note: The preceding discussion expands the
original in order to show that the action of \(E\) on \(G'\) is
obtained by combining the natural action of \(E\) on
\(\operatorname{\acute{E}p}_{S'}(\mathscr R)\) with the given action
on \(S'\). The canonical Killing couple of \(G'\) is preserved by
this action, as is the quasi-pinning given by the element
\(\widetilde X'\) of
\[
  \Gamma\!\left(
    \Delta_{S'},
    \Lie(G'/S')_{\Delta_{S'}}
  \right),
\]
equal to \(X'_\alpha\) on the copy of \(S'\) indexed by \(\alpha\).
Indeed, since \(E\) acts on \(\Delta_{S'}\) by permuting the copies
of \(S'\), one has \(h(\widetilde X')=\widetilde X'\) for every
\(h\in E\).}
Since \(S'\to S\) is a covering for the \textup{(fpqc)} topology, it
is an effective descent morphism for the fibered category of affine
morphisms. Write
\[
  q\acute{e}p(S')
  =
  Q\text{-}\operatorname{\acute{E}p}_{S'/S}(\mathscr R)
\]
for the quasi-pinned \(S\)-group obtained by Galois
descent.\footnote{Editor's note: See, for example, TDTE~I, p.~22,
\SGARefLink{sga3:xxiv:section:1}{Example~1}. One may also describe \(q\acute{e}p(S')\) as the twist of
\(G=\operatorname{\acute{E}p}_S(\mathscr R)\) by the \(E\)-torsor
\(S'/S\), that is, the \textup{(fpqc)} sheaf quotient of
\(S'\times_SG\) by the right action of \(E\) defined by
\[
  (s',g)\cdot h=(s'h,h^{-1}(g)).
\]
Since \(G\) is affine over \(S\), and since \(E\) acts on \(G\) by
group automorphisms, this sheaf is represented by an \(S\)-group
\(G^\sharp\), which is a ``twisted'' form of \(G\), and
\(D^\sharp=\operatorname{Dyn}(G^\sharp)\) is the twist of the
constant Dynkin scheme \(\Delta_S\) by the torsor \(S'/S\). Since
\(E\) normalizes \(B\) and \(T\), one obtains likewise a Killing
couple \((B^\sharp,T^\sharp)\) of \(G^\sharp\). On the other hand,
put
\(\mathfrak g=\Lie(G/S)\) and
\(\mathfrak g^\sharp=\Lie(G^\sharp/S)\). For every \(U\to S\), the
sections of \(\mathfrak g^\sharp\) over \(U\) are the
\(E\)-equivariant \(S\)-morphisms
\[
  U\times_SS'\longrightarrow W(\mathfrak g).
\]
Since \(D^\sharp\times_SS'\) is \(E\)-isomorphic to
\(\Delta\times S'\), endowed with the action
\[
  (\alpha,s')\cdot h=(h^{-1}(\alpha),s'h),
\]
the morphism \((\alpha,s')\mapsto X_\alpha\) gives a section of
\(\mathfrak g^\sharp\) over \(D^\sharp\) which is a quasi-pinning,
that is, a nowhere-zero section of
\[
  \bigl(
    \mathfrak g^\sharp\otimes_{\mathcal O_S}\mathcal O_{D^\sharp}
  \bigr)^{D^\sharp}
\]
(cf.\ \SGARefLink{sga3:xxiv:3.8}{3.8}).
}

\SGARefTarget{sga3:xxiv:3.11.3}
\medskip\noindent\textbf{3.11.3. The functor
\(\operatorname{rev}\).} ---
Let \(G\) be a reductive \(S\)-group of type \(\mathscr R\). Write
\[
  \operatorname{rev}(G)
  =
  \underline{\operatorname{Isomext}}
    \bigl(\operatorname{\acute{E}p}_S(\mathscr R),G\bigr).
\]
This is a principal homogeneous bundle under \(E_S\)
(cf.\ \SGARefLink{sga3:xxiv:1.10}{1.10} and \SGARefLink{sga3:xxiv:1.19}{1.19}), that is, an object of \(\mathsf{Rev}\).

\SGARefTarget{sga3:xxiv:3.11.4}
\medskip\noindent\textbf{3.11.4.} ---
\footnote{Editor's note: This subsection number has been added.}
It is clear from the isomorphism \SGARefLink{sga3:xxiv:3.10}{3.10} \((*)\) that the functor \(i\)
is fully faithful. On the other hand, if \(G_0,G,G'\) are reductive
\(S\)-groups, the natural morphism
\[
  \underline{\operatorname{Isom}}(G_0,G)
  \times_S
  \underline{\operatorname{Isom}}(G,G')
  \longrightarrow
  \underline{\operatorname{Isom}}(G_0,G')
\]
induces a morphism
\[
  (\star)\qquad
  \underline{\operatorname{Isomext}}(G_0,G)
  \times_S
  \underline{\operatorname{Isomext}}(G,G')
  \longrightarrow
  \underline{\operatorname{Isomext}}(G_0,G'),
\]
since, for every \(S'\to S\), if
\[
  g_0\in G_0(S'),\quad g\in G(S'),\quad
  u\in\underline{\operatorname{Isom}}(G_0,G)(S'),\quad
  v\in\underline{\operatorname{Isom}}(G,G')(S'),
\]
one has
\[
  v\circ\operatorname{int}(g)\circ u\circ\operatorname{int}(g_0)
  =
  v\circ u\circ
  \operatorname{int}\bigl(u^{-1}(g)g_0\bigr).
\]
Taking \(G_0=\operatorname{\acute{E}p}_S(\mathscr R)\), one obtains
a map
\[
\begin{aligned}
  \underline{\operatorname{Isomext}}(G,G')(S)
  \longrightarrow
  \operatorname{Hom}_{\mathsf{Rev}}\bigl(
  &\underline{\operatorname{Isomext}}
    (\operatorname{\acute{E}p}_S(\mathscr R),G),\\
  &\underline{\operatorname{Isomext}}
    (\operatorname{\acute{E}p}_S(\mathscr R),G')
  \bigr),
\end{aligned}
\]
and, to show that this map is a bijection, one may assume by
\textup{(fpqc)} descent that \(G\) and \(G'\) are pinned, in which
case it is evident.

% Authority local p. 20 / combined-reader p. 1265.

Similarly, if \(S_1,S_2\) are two objects of \(\mathsf{Rev}\), and
if one writes
\[
  S_i\times^E\operatorname{\acute{E}p}_S(\mathscr R)
\]
for the \(S\)-group \(q\acute{e}p(S_i)\) obtained by twisting
\(\operatorname{\acute{E}p}_S(\mathscr R)\) by the \(E\)-torsor
\(S_i\) (cf.\ editor's note~(21)), there is a natural map
\[
\begin{aligned}
  \operatorname{Hom}_{\mathsf{Rev}}(S_1,S_2)
  \longrightarrow
  \operatorname{Hom}_{\mathsf{Q\acute{e}p}}\bigl(
    &S_1\times^E\operatorname{\acute{E}p}_S(\mathscr R),\\
    &S_2\times^E\operatorname{\acute{E}p}_S(\mathscr R)
  \bigr).
\end{aligned}
\]
To show that this map is a bijection, one may assume by
\textup{(fpqc)} descent that \(S_1\) and \(S_2\) are trivial
\(E\)-bundles, in which case the assertion is evident.

Thus the three functors in the diagram are fully faithful. Moreover,
for every object \(S'\) of \(\mathsf{Rev}\), there is a natural
morphism
\[
  S'
  \longrightarrow
  \underline{\operatorname{Isomext}}\!\left(
    \operatorname{\acute{E}p}_S(\mathscr R),
    S'\times^E\operatorname{\acute{E}p}_S(\mathscr R)
  \right),
\]
and to see that it is an isomorphism, one may assume by
\textup{(fpqc)} descent that \(S'\) is a trivial \(E\)-bundle, in
which case the assertion is evident.

There is therefore an isomorphism of functors
\[
  \operatorname{Id}_{\mathsf{Rev}}
  \longrightarrow
  \operatorname{rev}\circ i\circ q\acute{e}p.
\]
Since \(\operatorname{rev}\) is fully faithful, for every reductive
\(S\)-group \(H\) (not necessarily quasi-split) one obtains a
bijection
\[
  \underline{\operatorname{Isomext}}\!\left(
    Q\text{-}\operatorname{\acute{E}p}_{\operatorname{rev}(H)/S}
      (\mathscr R),
    H
  \right)(S)
  \simeq
  \operatorname{Aut}_{\mathsf{Rev}}\bigl(\operatorname{rev}(H)\bigr),
\]
functorial in \(H\). In particular, the identity morphism of
\(\operatorname{rev}(H)\) corresponds to a ``canonical'' element
\[
  u_0\in
  \underline{\operatorname{Isomext}}\!\left(
    Q\text{-}\operatorname{\acute{E}p}_{\operatorname{rev}(H)/S}
      (\mathscr R),
    H
  \right)(S).
\]
Moreover, every element \(u\) of this same set corresponds to an
automorphism \(\varphi_u\) of \(\operatorname{rev}(H)\), and
\(q\acute{e}p(\varphi_u)\) is the unique automorphism of the pinned
\(S\)-group
\[
  Q\text{-}\operatorname{\acute{E}p}_{\operatorname{rev}(H)/S}
    (\mathscr R)
\]
such that
\[
  u_0\circ q\acute{e}p(\varphi_u)=u.
\]
This proves \SGARefLink{sga3:xxiv:3.11}{Theorem 3.11}.

We develop one of the corollaries of \SGARefLink{sga3:xxiv:3.11}{3.11}.

\SGARefTarget{sga3:xxiv:3.12}
\medskip\noindent\textbf{Corollary 3.12.} ---
{\itshape For every reductive \(S\)-group \(G\), there is a quasi-pinned
\(S\)-group \(G_{q\text{-}\acute{e}p}\) and an ``outer isomorphism''}
\[
  u\in
  \underline{\operatorname{Isomext}}
    (G_{q\text{-}\acute{e}p},G)(S).
\]
{\itshape The pair \((G_{q\text{-}\acute{e}p},u)\) is unique up to a
unique isomorphism.}

Indeed, one may assume that \(G\) has constant type \(\mathscr R\),
and one takes
\[
  G_{q\text{-}\acute{e}p}
  =
  Q\text{-}\operatorname{\acute{E}p}_{\operatorname{rev}(G)/S}
    (\mathscr R).
\]

\SGARefTarget{sga3:xxiv:3.12.1}
\medskip\noindent\textbf{3.12.1.} ---
Observe that the data of the pair
\((G_{q\text{-}\acute{e}p},u)\) canonically define the \(S\)-scheme
(cf.\ \SGARefLink{sga3:xxiv:1.11}{1.11})
\[
  Q=
  \underline{\operatorname{Isomint}}
    (G_{q\text{-}\acute{e}p},G),\footnote{Editor's note: We write
    \(\underline{\operatorname{Isomint}}
    (G_{q\text{-}\acute{e}p},G)\) instead of
    \(\underline{\operatorname{Isomint}}_u
    (G_{q\text{-}\acute{e}p},G)\) (cf.\ 1.11), since the pair
    \((G_{q\text{-}\acute{e}p},u)\) is unique up to a unique
    isomorphism.}
\]
which is a principal homogeneous bundle under
\(\operatorname{ad}(G_{q\text{-}\acute{e}p})\), and whose data are
``equivalent'' to those of the inner twisted form \(G\) of
\(G_{q\text{-}\acute{e}p}\). Moreover, \(Q\) is precisely the
``scheme of quasi-pinnings of \(G\),'' a definition that explains its
structure as a principal homogeneous bundle (on the left) under
\(\operatorname{ad}(G)\) (\SGARefLink{sga3:xxiv:3.10}{3.10}); compare \SGARefLink{sga3:xxiv:1.20}{1.20}.

% END inputs/XXIV/tex/components/08_expose_XXIV_39_proof_through_corollary312_en.tex

% BEGIN inputs/XXIV/tex/components/09_expose_XXIV_proposition313_through_remark3168_en.tex
% Authority coverage: Polo--Gille Exposé XXIV, local pages 21--24.
% Lead-reviewed against the authority PDF; the Reinhold Markdown lineage was comparison-only.

\textbf{\SGAInlineRefTarget{sga3:xxiv:proposition:3-13}Proposition 3.13.} \emph{Let \(S\) be a scheme, \(G\) a semisimple adjoint (resp. simply connected) \(S\)-group, \(B \supset T\) a
Killing couple of \(G\), \(\operatorname{Dyn}(G)\) the Dynkin \(S\)-scheme}

\emph{of \(G\). There exists a canonical isomorphism of \(S\)-group schemes}

\[
  T\xrightarrow{\sim}
  \prod_{\operatorname{Dyn}(G)/S}
  \mathbb G_{m,\operatorname{Dyn}(G)}.
\]

\emph{(Recall (cf.~Exp. II, § 1) that the right-hand side is by definition the \(S\)-functor which to \(S' \to S\) associates
\(G_{m}(\operatorname{Dyn}(G) \times_{S} S')\), or, what amounts to the same, \(G_{m}(Dyn(G_{S'}))\).)}

\emph{First proof.} Let us do it for simplicity in the adjoint case. Consider the composite morphism

\[
  T\longrightarrow\prod_{D/S}T_D
  \longrightarrow\prod_{D/S}\mathbb G_{m,D}.
\]

where the first morphism is the canonical morphism, the second is \(\prod_{D/S} \alpha_{D}\) (we have written \(D =
\operatorname{Dyn}(G)\)). To verify that this morphism is an isomorphism, one may assume \(G\) split; now, in this case, this is none
other than the morphism \(T \to (G_{m, S})^{\Delta}\) with components \(\alpha\), for \(\alpha \in \Delta\), and the latter is
an isomorphism (\SGARefLink{sga3:xxii:rem:4.3.8}{Exp. XXII, 4.3.8}).

\emph{Second proof.} By \SGARefLink{sga3:xxiv:2.8}{2.8}, \SGARefLink{sga3:xxiv:3.5}{3.5} and \SGARefLink{sga3:xxiv:3.11}{3.11}, one may assume that \(G = Q-\acute{E}p_{S'/S}(R)\), \(T\) being the canonical maximal
torus. On \(S'\), one has by \SGARefLink{sga3:xxii:rem:4.3.8}{Exp. XXII 4.3.8}, an isomorphism \(T_{S'} \xrightarrow{\sim} (G_{m, S'})^{\Delta}\), defined by
the simple roots (resp. the simple coroots). The group \(E\) operates on the right-hand side by permutation of \(\Delta\).
Now the bundle associated to \(S'/S\) by \(E \to \operatorname{Aut}(\Delta)\) is \(\operatorname{Dyn}(G)\) (\SGARefLink{sga3:xxiv:3.7}{3.7}), and one concludes at once.

Using Prop. \SGARefLink{sga3:xxiv:proposition:8-4}{8.4} of the appendix, one deduces:

\textbf{\SGAInlineRefTarget{sga3:xxiv:corollary:3-14}Corollary 3.14.} \emph{Under the preceding conditions, one has}

\[
  H^1(S,T)\xrightarrow{\sim}
  H^1\bigl(\operatorname{Dyn}(G),\mathbb G_m\bigr)
  \xrightarrow{\sim}\operatorname{Pic}\bigl(\operatorname{Dyn}(G)\bigr).
\]

\emph{In particular, \(H^{1}(S, T) = 0\) when \(S\) is semi-local.}

\textbf{\SGAInlineRefTarget{sga3:xxiv:remark:3-15}Remark 3.15.} Let \(S\) be a scheme, \(G\) (resp. \(G'\)) an \(S\)-reductive group, \(B \supset T\) (resp. \(B' \supset T'\)) a
Killing couple of \(G\) (resp. \(G'\)), \(u \in Isomext(G, G')(S)\). Set (cf.~\SGARefLink{sga3:xxiv:2.10}{2.10})

\[
  P=\underline{\operatorname{Isomint}}_u(G,B,T;G',B',T').
\]

this is a principal homogeneous bundle under \(T_{ad}\) (by \((f, t) \mapsto f \circ int(t)\)).

Let on the other hand \(D = \operatorname{Dyn}(G) = \operatorname{Dyn}(G')\) (identified via \(u\) (\SGARefLink{sga3:xxiv:3.5}{3.5})), and let \(L = Isom_{D}(g_{D}, g'_{D})\) be the
principal homogeneous bundle under \(G_{m,D}\) defined by the invertible \(\mathcal O_D\)-module

\[
  \operatorname{Hom}_{\mathcal O_D}(\mathfrak g_D,\mathfrak g'_D)
  =\mathfrak g'_D\otimes(\mathfrak g_D)^\vee .
\]

Each \(f \in P(S')\) defines an isomorphism of \(g_{D} \otimes_{O_{S}} O_{S'}\) onto \(g'_{D} \otimes_{O_{S}} O_{S'}\), whence
a canonical morphism

\[ P \to \prod_{D/S} L. \]

This morphism is an isomorphism, compatible with the isomorphism of operator groups

\[
  T_{\mathrm{ad}}\xrightarrow{\sim}
  \prod_{D/S}\mathbb G_{m,D}.
\]

defined above. Indeed, it suffices to verify it in the case where the two groups are pinned, where it is easy.

It follows in particular that in the isomorphism

\[ H^{1}(S, T_{ad}) \xrightarrow{\sim} \operatorname{Pic}(\operatorname{Dyn}(G)) \]

of \SGARefLink{sga3:xxiv:corollary:3-14}{3.14}, the class of the bundle \(P\) is transformed into \(c\ell(g'^{D}) - c\ell(g_{D})\). The bundle \(P\) is therefore
trivial if and only if \(g'^{D}\) and \(g_{D}\) are isomorphic.

If \((G, B, T)\) is quasi-splittable, for example if one takes for \(G\) the group \(G'_{q-\acute{e}p.}\), with its canonical
Killing couple, it follows that the image of the class of \(P\) is none other than the obstruction to the quasi-splitting
of \(G'\) defined in \SGARefLink{sga3:xxiv:3.9}{3.9}.

\subsection*{\SGAInlineRefTarget{sga3:xxiv:subsection:3-16}3.16. Symmetries}

\SGAInlineRefTarget{sga3:xxiv:subsubsection:3-16-1}\subsubsection*{3.16.1.}

Let \(S\) be a scheme, \(G\) an \(S\)-reductive group, \(B \supset T\) a Killing couple of \(G\). Write \(D = \operatorname{Dyn}(G)\). Recall (\SGARefLink{sga3:xxii:prop:5.9.2}{Exp.
XXII, 5.9.2}) that there exists a unique Borel subgroup \(B^{-}\) of \(G\) such that \(B \cap B^{-} = T\). If \(X \in \Gamma(D,
g_{D})^{\times}\) defines a quasi-pinning of \(G\) relative to \((B, T)\) (\SGARefLink{sga3:xxiv:3.9}{3.9}), then \(Y = -X^{-1} \in \Gamma(D,
(g_{D})^{\vee})^{\times}\) defines a quasi-pinning of \(G\) relative to \((B^{-}, T)\); one says that this is the \emph{opposite
quasi-pinning}.

If \(R\) is a pinned reduced root datum and if \(E\) is an \(R\)-pinning of the \(S\)-reductive group \(G\), one defines an
\(R\)-pinning \(E^{-}\) said to be \emph{opposite to \(E\)} in the following way: one keeps the same maximal torus \(T\), one takes
the opposite of the isomorphism \(D_{S}(M) \xrightarrow{\sim} T\), and one ``pins'' by \(Y_{\alpha} = -X^{-1}_{\alpha} \in
\Gamma(S, g^{-\alpha})^{\times}\), for \(\alpha \in \Delta\). The quasi-pinning underlying \(E^{-}\) is the quasi-pinning
opposite to the quasi-pinning underlying \(E\).

\textbf{Remark.} In the notation of \SGARefLink{sga3:xx:3.1}{Exp. XX, 3.1}, if one sets

\[
  w_\alpha(X_\alpha)
  =\exp(X_\alpha)\exp(-X_\alpha^{-1})\exp(X_\alpha).
\]

one has \(w_{\alpha}(X_{\alpha}) = w_{-\alpha}(Y_{\alpha})\) (loc. cit., 3.1 (vi)).

\textbf{\SGAInlineRefTarget{sga3:xxiv:proposition:3-16-2}Proposition 3.16.2.} \emph{Let \(S\) be a scheme, \(G\) an \(S\)-reductive group.}

\emph{\SGAInlineRefTarget{sga3:xxiv:proposition:3-16-2:item:i}(i) Let \(T\) be a maximal torus of \(G\); there exists a unique}

\[
  i_T\in
  \left(
    \underline{\operatorname{Aut}}_{S\text{-}\mathrm{gr.}}(G,T)/T_{\mathrm{ad}}
  \right)(S)
  \subset
  \underline{\operatorname{Aut}}_{S\text{-}\mathrm{gr.}}(T).
\]

\emph{such that \(i_{T}(t) = t^{-1}\) for every section \(t\) of \(T\).}

\emph{\SGAInlineRefTarget{sga3:xxiv:proposition:3-16-2:item:ii}(ii) Let \((B, T)\) be a Killing couple of \(G\); there exists a unique section}

\[
  w_{B,T}\in
  \bigl(\operatorname{Norm}_G(T)/T\bigr)(S)=W_G(T)(S).
\]

\emph{such that \(int(w_{B, T})(B) = B^{-}\) (with the obvious abuse of language).}

\emph{\SGAInlineRefTarget{sga3:xxiv:proposition:3-16-2:item:iii}(iii) Let \(Q = (B, T, X)\) be a quasi-pinning of \(G\), \(Q^{-} = (B^{-}, T, Y)\) the opposite quasi-pinning; there exists a
unique inner automorphism of \(G\)}

\[ n_{Q} \in Norm_{ad(G)}(T_{ad})(S) \subset \underline{\operatorname{Aut}}_{S\text{-}\mathrm{gr.}}(G) \]

\emph{such that \(n_{Q}(Q) = Q^{-}\), that is, \(n_{Q}(T) = T\), \(n_{Q}(B) = B^{-}\), \(n_{Q}(X) = Y\).}

\emph{\SGAInlineRefTarget{sga3:xxiv:proposition:3-16-2:item:iv}(iv) Let \((R, E)\) be a pinning of \(G\), \((R, E^{-})\) the opposite pinning; there exists a unique automorphism of \(G\)}

\[
  u_E\in
  \underline{\operatorname{Aut}}_{S\text{-}\mathrm{gr.}}(G,T)
  \subset
  \underline{\operatorname{Aut}}_{S\text{-}\mathrm{gr.}}(G).
\]

\emph{such that \(u_{E}(E) = E^{-}\), i.e.~\(u_{E}(t) = t^{-1}\) for every section \(t\) of \(T\), and \(Ad(u_{E}) X_{\alpha} =
Y_{\alpha}\) for every \(\alpha \in \Delta\).}

\emph{Proof.} \SGARefLink{sga3:xxiv:proposition:3-16-2:item:ii}{(ii)} follows from \SGARefLink{sga3:xxii:corollary:5-5-5:item:ii}{Exp. XXII, 5.5.5 (ii)}, then \SGARefLink{sga3:xxiv:proposition:3-16-2:item:iii}{(iii)} follows from \SGARefLink{sga3:xxiv:3.10}{3.10}, and \SGARefLink{sga3:xxiv:proposition:3-16-2:item:iv}{(iv)} follows from \SGARefLink{sga3:xxiii:theorem:4-1}{Exp. XXIII, 4.1}.
Finally to prove \SGARefLink{sga3:xxiv:proposition:3-16-2:item:i}{(i)}, one may assume \(G\) pinned. Existence follows from \SGARefLink{sga3:xxiv:proposition:3-16-2:item:iv}{(iv)} for example, uniqueness from the fact that
an automorphism of \(G\) inducing the identity on \(T\) is given by a section of \(T_{ad}\) (\SGARefLink{sga3:xxiv:2.11}{2.11}).

\textbf{\SGAInlineRefTarget{sga3:xxiv:corollary:3-16-3}Corollary 3.16.3.} \emph{One has}

\[
  i_T^2=w_{B,T}^2=n_Q^2=u_E^2=e.
\]

\emph{Moreover, \(i_{T}\) (resp. \(u_{E}\)) is \(\neq e\) if \(G \neq e\), and \(w_{B, T}\) (resp. \(n_{Q}\)) is \(\neq e\) if \(G\) is not a
torus.}

\textbf{\SGAInlineRefTarget{sga3:xxiv:corollary:3-16-4}Corollary 3.16.4.} \emph{In the situation of \SGARefLink{sga3:xxiv:proposition:3-16-2:item:iii}{(iii)} (resp. \SGARefLink{sga3:xxiv:proposition:3-16-2:item:iv}{(iv)}), \(n_{Q}\) projects onto \(w_{B, T}\) (resp. \(u_{E}\) projects
onto \(i_{T}\)) by the canonical morphism}

\[ Norm_{ad(G)}(T_{ad}) \to W_{ad(G)}(T_{ad}) \simeq W_{G}(T), \]

\emph{resp.}

\[
  \underline{\operatorname{Aut}}_{S\text{-}\mathrm{gr.}}(G,T)
  \longrightarrow
  \underline{\operatorname{Aut}}_{S\text{-}\mathrm{gr.}}(G,T)/T_{\mathrm{ad}}.
\]

\textbf{\SGAInlineRefTarget{sga3:xxiv:corollary:3-16-5}Corollary 3.16.5.} \emph{The preceding definitions are compatible with base change, and are functorial under isomorphism
(in an obvious sense).}

\textbf{\SGAInlineRefTarget{sga3:xxiv:proposition:3-16-6}Proposition 3.16.6.} \emph{\SGAInlineRefTarget{sga3:xxiv:proposition:3-16-6:item:i}(i) One can define uniquely for each reductive group \(G\) over a scheme \(S\) an element}

\[ s_{G} \in Autext(G)(S) \]

\emph{in such a way that this construction is functorial in \(G\) under isomorphism, is compatible with base change, and is
such that whenever \(T\) is a maximal torus of the \(S\)-reductive group \(G\), \(s_{G}\) is the image of \(i_{T}\) under the
canonical morphism}

\[
  \underline{\operatorname{Aut}}_{S\text{-}\mathrm{gr.}}(G,T)/T_{\mathrm{ad}}
  \longrightarrow\underline{\underline{\operatorname{Autext}}}(G).
\]

\emph{\SGAInlineRefTarget{sga3:xxiv:proposition:3-16-6:item:ii}(ii) One has \(s^{2}_{G} = e\), and \(s_{G}\) is a central element of \(Autext(G)\).}

\emph{\SGAInlineRefTarget{sga3:xxiv:proposition:3-16-6:item:iii}(iii) Under the conditions of \SGARefLink{sga3:xxiv:proposition:3-16-2:item:ii}{3.16.2 (ii)}, if one identifies \(\underline{\operatorname{Aut}}_{S\text{-}\mathrm{gr.}}(G, B, T) / T_{ad}\) with
\(Autext(G)\) (\SGARefLink{sga3:xxiv:2.2}{2.2}), one has}

\[
  w_{B,T}i_T=i_Tw_{B,T}=s_G.
\]

\emph{\SGAInlineRefTarget{sga3:xxiv:proposition:3-16-6:item:iv}(iv) Under the conditions of \SGARefLink{sga3:xxiv:proposition:3-16-2:item:iv}{3.16.2 (iv)}, if one identifies \(\operatorname{Aut}_{S-gr.-\acute{e}p}(G, R, E)\) with
\(Autext(G)\) (\SGARefLink{sga3:xxiv:1.3:item:iii}{1.3 (iii)}), one has}

\[
  n_Eu_E=u_En_E=s_G.
\]

\emph{Proof.} \SGARefLink{sga3:xxiv:proposition:3-16-6:item:i}{(i)} is proved without difficulty by descent. On the other hand, since \(i_{T}\) is evidently a central section of
square \(e\) in \(\underline{\operatorname{Aut}}_{S\text{-}\mathrm{gr.}}(T)\), \SGARefLink{sga3:xxiv:proposition:3-16-6:item:ii}{(ii)} follows immediately; \SGARefLink{sga3:xxiv:proposition:3-16-6:item:iii}{(iii)} is a consequence of \SGARefLink{sga3:xxiv:proposition:3-16-6:item:iv}{(iv)} by descent.
Finally, under the conditions of \SGARefLink{sga3:xxiv:proposition:3-16-6:item:iv}{(iv)}, it is clear that \(n_{E} u_{E} = u_{E} n_{E}\) and that this automorphism of \(G\)
respects the pinning; modulo the identification made, it is therefore equal to its image in \(Autext(G)\); but \(n_{E}\) is
inner and \(u_{E}\) projects onto \(s_{G}\).

\textbf{\SGAInlineRefTarget{sga3:xxiv:remark:3-16-7}Remark 3.16.7.} \SGARefLink{sga3:xxiv:proposition:3-16-6:item:i}{(i)} One determines \(s_{G}\) explicitly in each case of the classification thanks to \SGARefLink{sga3:xxiv:proposition:3-16-6:item:iii}{(iii)}: for each
irreducible pinned root datum it suffices to compose the symmetry through the origin with the symmetry in the Weyl group
(i.e.~the element \(w_{0}\) of the Weyl group such that \(w_{0}(\Delta) = -\Delta\)). One finds the following results: one
has \(s_{G} = e\) except for \(A_{n}\) (\(n \geqslant 2\)), \(D_{n}\) (\(n\) odd) and \(E_6\), in which case \(s_{G}\) is the unique
non-trivial ``outer automorphism''.

\begin{enumerate}
\def\labelenumi{(\roman{enumi})}
\setcounter{enumi}{1}
\item
  The automorphism \(u_{E}\) is the one used to construct ``the compact real forms'' in the theory of semisimple Lie
  algebras.
\end{enumerate}

\textbf{\SGAInlineRefTarget{sga3:xxiv:remark:3-16-8}Remark 3.16.8.} We defined in \SGARefLink{sga3:xxiv:subsubsection:3-16-1}{3.16.1} an involution in the \(S\)-scheme
\[
Q = Isomint(G_{q-\acute{e}p.}, G)
\]
of quasi-pinnings of \(G\) (cf.~\SGARefLink{sga3:xxiv:3.12.1}{3.12.1}); by transport of structure from \(G_{q-\acute{e}p.}\) to \(G\), one sees at once that
this involution is given by the action of an element of \(ad(G_{q-\acute{e}p.})(S)\): the element \(n_{0}\) defined (\SGARefLink{sga3:xxiv:proposition:3-16-2:item:iii}{3.16.2
(iii)}) by the canonical quasi-pinning of \(G_{q-\acute{e}p.}\).

In the same way, one has defined an involution in the \(S\)-scheme

\[ P = Isom_{S-gr.}(\acute{E}p_{S}(R), G) \]

of \(R\)-pinnings of \(G\) (cf.~\SGARefLink{sga3:xxiv:1.20}{1.20}). Arguing as before, one sees that this involution is given by the action of the
automorphism \(u_{0}\) of \(\acute{E}p_{S}(R)\) defined (\SGARefLink{sga3:xxiv:proposition:3-16-2:item:iv}{3.16.2 (iv)}) by the canonical pinning of \(\acute{E}p_{S}(R)\).

% END inputs/XXIV/tex/components/09_expose_XXIV_proposition313_through_remark3168_en.tex

% BEGIN inputs/XXIV/tex/components/10_expose_XXIV_section4_through_theorem441_setup_en.tex
% Authority coverage: Polo--Gille Exposé XXIV, local pages 25--29.
% Lead-reviewed against the authority PDF; authority images remain explicit native-diagram debt.

\section*{\SGAInlineRefTarget{sga3:xxiv:section:4}4. Isotriviality of reductive groups and of principal bundles under reductive groups}

\subsection*{\SGAInlineRefTarget{sga3:xxiv:subsection:4-1}4.1. Definitions. Isotriviality theorem}

\textbf{\SGAInlineRefTarget{sga3:xxiv:definition:4-1-1}Definition 4.1.1.} \emph{Let \(S\) be a scheme, \(G\) an \(S\)-group scheme, \(P\) a principal homogeneous bundle under \(G\). One
says that \(P\) is} locally isotrivial \emph{(resp.} semi-locally isotrivial) if for every point \(s \in S\) (resp. every
finite set \(F\) of points of \(S\) contained in an affine open set) there exist an open set \(U\) of \(S\) containing \(s\)
(resp. \(F\)) and a finite surjective étale morphism \(S' \to U\) such that \(P_{S'}\) is trivial.

\textbf{\SGAInlineRefTarget{sga3:xxiv:definition:4-1-2}Definition 4.1.2.} \emph{Let \(S\) be a scheme, \(G\) an \(S\)-reductive group. One says that \(G\) is} locally isotrivial
\emph{(resp.} semi-locally isotrivial) if for every point \(s \in S\) (resp. every finite set \(F\) of points of \(S\) contained
in an affine open set) there exist an open set \(U\) of \(S\) containing \(s\) (resp. \(F\)) and a finite surjective étale
morphism \(S' \to U\) such that \(G_{S'}\) is splittable.*

\textbf{\SGAInlineRefTarget{sga3:xxiv:remark:4-1-3}Remark 4.1.3.} (i) The equivalence of categories \SGARefLink{sga3:xxiv:definition:4-1-1}{4.1.1} of \SGARefLink{sga3:xxiv:1.17}{1.17} respects by definition local (resp. semi-local)
isotriviality.

\begin{enumerate}
\def\labelenumi{(\roman{enumi})}
\setcounter{enumi}{1}
\item
  Add to the conditions of \SGARefLink{sga3:xxiv:definition:4-1-1}{4.1.1}: \(G\) locally of finite presentation over \(S\). Then the principal bundle \(P\) (or the
  reductive group \(G\)) is locally isotrivial (resp. semi-locally isotrivial) if and only if for every \(S' \to S\), \(S'\)
  local (resp. semi-local), \(P_{S'}\) is isotrivial (or \(G_{S'}\) isotrivial), that is to say, if there exists \(S'' \to S'\)
  finite surjective étale such that \(P_{S''}\) is trivial (or \(G_{S''}\) split).
\item
  In the case of tori, \SGARefLink{sga3:xxiv:definition:4-1-2}{definition 4.1.2} coincides with that of \SGARefLink{sga3:ix:definition:1:1}{Exp. IX, 1.1}.
\end{enumerate}

\subsubsection*{4.1.4.}

Recall (\SGARefLink{sga3:xxii:section:4.3}{Exp. XXII, 4.3} and \SGARefLink{sga3:xxiv:lemma:6-2}{6.2}) that for every reductive group \(G\), we have introduced the groups \(rad(G)\), \(corad(G)\)
and \(d\acute{e}r(G)\). The groups \(rad(G)\) and \(corad(G)\) are tori, which are split when \(G\) is split; moreover, there
exists an isogeny \(rad(G) \to corad(G)\). The \(S\)-group \(d\acute{e}r(G)\) is semisimple, one has \(G / d\acute{e}r(G) =
corad(G)\); it follows that for every principal homogeneous bundle \(P\) under \(G\), \(P / d\acute{e}r(G)\) is a principal
homogeneous bundle under \(corad(G)\). This said, one has:

\textbf{\SGAInlineRefTarget{sga3:xxiv:theorem:4-1-5}Theorem 4.1.5.} \emph{Let \(S\) be a scheme, \(G\) an \(S\)-reductive group of constant type.}

\emph{\SGAInlineRefTarget{sga3:xxiv:theorem:4-1-5:item:i}(i) The following conditions are equivalent:}

\emph{\SGAInlineRefTarget{sga3:xxiv:theorem:4-1-5:item:a}(a) \(G\) is locally (resp. semi-locally) isotrivial.}

\emph{\SGAInlineRefTarget{sga3:xxiv:theorem:4-1-5:item:b}(b) The torus \(rad(G)\) is.}

\emph{(b') The torus \(corad(G)\) is.}

\emph{\SGAInlineRefTarget{sga3:xxiv:theorem:4-1-5:item:c}(c) The Galois covering of \(S\) associated to \(G\) (\SGARefLink{sga3:xxiv:3.11}{3.11}) is.}

\emph{If \(T\) is a maximal torus of \(G\), these conditions are also equivalent to}

\emph{\SGAInlineRefTarget{sga3:xxiv:theorem:4-1-5:item:d}(d) The torus \(T\) is locally (resp. semi-locally) isotrivial.}

\emph{\SGAInlineRefTarget{sga3:xxiv:theorem:4-1-5:item:ii}(ii) Let \(P\) be a principal homogeneous bundle under \(G\); for \(P\) to be locally (resp. semi-locally) isotrivial, it is
necessary and sufficient that the \(corad(G)\)-principal bundle \(P / d\acute{e}r(G)\) be so.}

\textbf{\SGAInlineRefTarget{sga3:xxiv:corollary:4-1-6}Corollary 4.1.6.} \emph{The conditions of \SGARefLink{sga3:xxiv:theorem:4-1-5:item:i}{(i)} are satisfied when \(G\) is semisimple or when \(S\) is locally noetherian and
normal (or more generally geometrically unibranch). The conditions of \SGARefLink{sga3:xxiv:theorem:4-1-5:item:ii}{(ii)} are satisfied when \(G\) is locally (resp.
semi-locally) isotrivial.}

For \SGARefLink{sga3:xxiv:theorem:4-1-5:item:i}{(i)}, the first assertion is trivial from \SGARefLink{sga3:xxiv:theorem:4-1-5:item:b}{(b)}, the second follows from \SGARefLink{sga3:xxiv:theorem:4-1-5:item:c}{(c)} and \SGARefLink{sga3:x:corollary:5-14}{Exp. X, 5.14} and 5.15. For \SGARefLink{sga3:xxiv:theorem:4-1-5:item:ii}{(ii)}, it
suffices to note that by virtue of \SGARefLink{sga3:viii:corollary:4:5}{theorem 90}, a principal bundle under a split torus is semi-locally isotrivial.

\subsection*{\SGAInlineRefTarget{sga3:xxiv:subsection:4-2}4.2. Proof: the semisimple case}

Let us first prove, for later reference:

\textbf{\SGAInlineRefTarget{sga3:xxiv:proposition:4-2-1}Proposition 4.2.1.} \emph{Let \(S\) be a scheme, \(G\) an \(S\)-reductive group, \(T\) a maximal torus of \(G\) (resp. \(B\) a Borel
subgroup, resp. \(B \supset T\) a Killing couple of \(G\)), \(P\) a principal homogeneous bundle under \(G\), \(G'\) the twisted
form of \(G\) associated to \(P\) (via the morphism \(int : G \to \operatorname{Aut}_{S-gr.}(G)\)). One has canonical
isomorphisms}

\[
  P/\operatorname{Norm}_G(T)\xrightarrow{\sim}\operatorname{Tor}(G'),
  \qquad
  P/B\xrightarrow{\sim}\operatorname{Bor}(G'),
  \qquad
  P/T\xrightarrow{\sim}\operatorname{Kil}(G').
\]

By construction, \(G'\) is the quotient of \(P \times_{S} G\) by a certain action of \(G\) (\(((p, g')) g = (pg, g^{-1} g'
g)\)); one therefore has a morphism \(P \times_{S} G \to G'\), that is, a morphism

\[ P \to \operatorname{Hom}_{S}(G, G'), \]

which, as one sees at once, factors through a morphism

\[
  f:P\longrightarrow\operatorname{Isom}_{S\text{-}\mathrm{gr.}}(G,G').
\]

(to verify this assertion, one may assume \(G = P\), in which case one has \(G' = G\), \(f = int\)). The set-theoretic map \(p
\mapsto f(p)(T)\) defines a morphism

\[ P \to Tor(G'). \]

To verify that this morphism induces an isomorphism \(P / Norm_{G}(T) \xrightarrow{\sim} Tor(G')\) as announced, one may
again assume \(P = G\) in which case one is reduced to \SGARefLink{sga3:xxii:cor:5.8.3:item:iii}{Exp. XXII, 5.8.3 (iii)}. (In fact, \emph{loc. cit.} should be replaced by
the statement above.) One argues similarly for \(\operatorname{Bor}\) and \(\operatorname{Kil}\).

\textbf{\SGAInlineRefTarget{sga3:xxiv:proposition:4-2-2}Proposition 4.2.2.} \emph{Let \(S\) be a semi-local scheme, \(G\) an \(S\)-semisimple group of constant type.}

\emph{\SGAInlineRefTarget{sga3:xxiv:proposition:4-2-2:item:i}(i) \(G\) is isotrivial.}

\emph{(ii) Every principal bundle under \(G\) is isotrivial.}

Let us prove \SGARefLink{sga3:xxiv:proposition:4-2-2:item:i}{(i)}. Up to making a finite surjective étale extension of the base, one may, by \SGARefLink{sga3:xxiv:3.9.2}{3.9.2}, assume \(G\)
quasi-split. But then \(G\) is isotrivial by construction (\SGARefLink{sga3:xxiv:3.11}{3.11}, the group \(E\) being finite). To prove \SGARefLink{sga3:xxiv:3.11:item:ii}{(ii)}, one may, by
\SGARefLink{sga3:xxiv:3.11:item:i}{(i)}, assume \(G\) split; one is then reduced to:

\textbf{\SGAInlineRefTarget{sga3:xxiv:lemma:4-2-3}Lemma 4.2.3.} \emph{Let \(S\) be a semi-local scheme. Every principal bundle under a split reductive group is isotrivial.}

Indeed, with the notation of \SGARefLink{sga3:xxiv:proposition:4-2-1}{4.2.1}, where \(B \supset T\) denotes the canonical Killing couple of the split \(G\), the
\(S\)-scheme \(Kil(G')\) possesses a section, after finite surjective étale extension of the base, by \SGARefLink{sga3:xxiv:3.9.2}{3.9.2}. One may
therefore reduce the structure group of \(G\) to \(T\); but \(T\) is split, so \(H^{1}(S, T) = 0\) (theorem 90).

\textbf{\SGAInlineRefTarget{sga3:xxiv:corollary:4-2-4}Corollary 4.2.4.} \emph{Let \(S\) be a scheme and}

\[ 1 \to G \to G' \to G'' \to 1 \]

\emph{an exact sequence of \(S\)-group schemes (for (fpqc)), \(G\) being semisimple of constant type. Let \(P\) be a principal
homogeneous bundle under \(G'\), suppose the sheaf\footnote{N.D.E.: We have replaced ``bundle'' by ``sheaf'', and then \(G'\) by \(G''\).} \(P / G\) associated representable (for example \(G''\)
affine over \(S\)). For \(P\) to be locally isotrivial (resp. semi-locally isotrivial), it is necessary and sufficient that
\(P / G\) be so (as bundle under \(G''\)).}

If \(P\) is trivial, \(P / G\) is too, which shows that the condition is necessary. Conversely, suppose \(S\) local (resp.
semi-local) and \(P / G\) isotrivial, so trivialized by a finite surjective étale extension \(S'\) of \(S\). Extending the
base to \(S'\), one may reduce the structure group of \(P_{S'}\) to \(G_{S'}\). But \(S'\) is still semi-local and \(G_{S'}\)
semisimple of constant type, so the obtained bundle is isotrivial (\SGARefLink{sga3:xxiv:proposition:4-2-2}{4.2.2}).

\subsection*{\SGAInlineRefTarget{sga3:xxiv:subsection:4-3}4.3. Proof: general case}

Let us first note that \SGARefLink{sga3:xxiv:theorem:4-1-5:item:ii}{4.1.5 (ii)} follows at once from the application of \SGARefLink{sga3:xxiv:corollary:4-2-4}{4.2.4} to the exact sequence

\[ 1 \to d\acute{e}r(G) \to G \to corad(G) \to 1. \]

Let us therefore prove \SGARefLink{sga3:xxiv:proposition:4-2-2:item:i}{(i)}. If \(G\) is split, \(rad(G)\) and \(corad(G)\) are split, as is \(rev(G)\); so \SGARefLink{sga3:xxiv:theorem:4-1-5:item:a}{(a)} implies \SGARefLink{sga3:xxiv:theorem:4-1-5:item:b}{(b)},
(b'), and \SGARefLink{sga3:xxiv:theorem:4-1-5:item:c}{(c)}.

\subsubsection*{4.3.1.}

One has \SGARefLink{sga3:xxiv:theorem:4-1-5:item:c}{(c)} \(\Rightarrow\) \SGARefLink{sga3:xxiv:theorem:4-1-5:item:a}{(a)}. Let \(R\) be the type of \(G\); one has an exact sequence

\[ 1 \to ad(\acute{E}p_{S}(R)) \to A_{S}(R) \to E_{S} \to 1. \]

Applying \SGARefLink{sga3:xxiv:corollary:4-2-4}{4.2.4} to the bundle \(P = Isom_{S-gr.}(\acute{E}p_{S}(R), G)\) and to the associated bundle \(rev(G) = P /
ad(\acute{E}p_{S}(R))\), one has \SGARefLink{sga3:xxiv:theorem:4-1-5:item:c}{(c)} \(\Rightarrow\) \SGARefLink{sga3:xxiv:theorem:4-1-5:item:a}{(a)}.

\subsubsection*{4.3.2.}

One has \SGARefLink{sga3:xxiv:theorem:4-1-5:item:b}{(b)} \(\Rightarrow\) \SGARefLink{sga3:xxiv:theorem:4-1-5:item:a}{(a)}. It suffices to prove that if \(rad(G)\) is split, \(G\) is semi-locally isotrivial. Now let \(G_{0} =
\acute{E}p_{S}(R)\); consider the category of pairs \((G',f)\), where \(G'\) is a form of \(G_0\) and \(f\) is an isomorphism of
\(rad(G_{0})\) onto \(rad(G)\).

One knows (\SGARefLink{sga3:xxiv:2.16}{2.16}) that this category is equivalent to the category of principal homogeneous bundles under a certain
\(S\)-group \(H\) extension of a finite twisted constant group by a semisimple group. It suffices to prove that every
principal bundle under \(H\) is semi-locally isotrivial; this follows at once from \SGARefLink{sga3:xxiv:corollary:4-2-4}{4.2.4}.

\subsubsection*{4.3.3.}

One has (b') \(\Rightarrow\) \SGARefLink{sga3:xxiv:theorem:4-1-5:item:a}{(a)}. One can argue as previously (it will moreover be the same group \(H\) that arises). One can also see
that \SGARefLink{sga3:xxiv:theorem:4-1-5:item:b}{(b)} and (b') are equivalent: a torus isogenous to a locally split torus is also locally split (cf.~\SGARefLink{sga3:ix:proposition:2:11:item:iii:01}{Exp. IX, 2.11
(iii)}).

\subsubsection*{4.3.4.}

One has \SGARefLink{sga3:xxiv:theorem:4-1-5:item:d}{(d)} \(\Rightarrow\) \SGARefLink{sga3:xxiv:theorem:4-1-5:item:a}{(a)}. It suffices to prove that a group of constant type possessing a split maximal torus is semi-locally
isotrivial; this follows at once from \SGARefLink{sga3:xxiv:2.14}{2.14}.

\subsubsection*{4.3.5.}

One has \SGARefLink{sga3:xxiv:theorem:4-1-5:item:a}{(a)} \(\Rightarrow\) \SGARefLink{sga3:xxiv:theorem:4-1-5:item:d}{(d)}. It suffices to prove that a maximal torus of a split group is semi-locally isotrivial. More
precisely:

\textbf{\SGAInlineRefTarget{sga3:xxiv:lemma:4-3-6}Lemma 4.3.6.} \emph{Let \(S\) be a semi-local scheme, \(G\) an \(S\)-reductive group, \(T_0\) a split maximal torus of \(G\), \(W_{0}
= Norm_{G}(T_{0}) / T_{0}\) (this is a locally constant \(S\)-group, constant if \(G\) is of constant type, by \SGARefLink{sga3:xxiv:2.14}{2.14}), \(T\) a
maximal torus of \(G\).}

\emph{There exists a morphism \(S' \to S\) which is principal homogeneous under \(W_0\) (hence étale finite and surjective, and
even principal Galois if \(G\) is of constant type) such that \(T_{S'}\) is conjugate to \((T_{0})_{S'}\) by an element of
\(G(S')\) (and so in particular split).}

Indeed, one knows that \(Transp_{G}(T_{0}, T)\) is a principal homogeneous bundle under \(Norm_{G}(T_{0})\) (cf.~for example
\SGARefLink{sga3:xi:corollary:5:4:bis}{Exp. XI, 5.4 bis}). Set \(S' = Transp_{G}(T_{0}, T) / T_{0}\); this is a principal homogeneous bundle under \(W_0\).
Extending the base from \(S\) to \(S'\), one can reduce the structure group of \(Transp_{G}(T_{0}, T)\) to \(T_0\). Now \(S'\) is
semi-local and \(T_0\) split, so \(Transp_{G}(T_{0}, T)\) possesses a section over \(S'\).\footnote{N.D.E.: by ``theorem 90''.} QED

\subsection*{\SGAInlineRefTarget{sga3:xxiv:subsection:4-4}4.4. Use of the existence of maximal tori}

Using Grothendieck's existence theorem for maximal tori (\SGARefLink{sga3:xiv:3.20}{Exp. XIV, 3.20}), one can considerably sharpen the preceding
results. Let us state at once:

\textbf{\SGAInlineRefTarget{sga3:xxiv:theorem:4-4-1}Theorem 4.4.1.} \emph{Let \(S\) be a semi-local scheme, \(R\) a pinned reduced root datum, \(W\) its Weyl group, \(E\) the group
of its automorphisms. (Recall that \(E\) operates naturally on \(W\) and that the semi-direct product \(A = W \cdot E\) is
identified with the group of automorphisms of \(R\) non-pinned, cf.~\SGARefLink{sga3:XXI:6.7.2}{Exp. XXI, 6.7.2}).}

\emph{(i) Every principal homogeneous bundle under \(\acute{E}p_{S}(R)\) is trivialized by a principal Galois covering \(S' \to
S\) of group \(W\).}

\emph{(ii) Let \(G\) be an \(S\)-reductive group of type \(R\); let \(rev(G) = Isomext(\acute{E}p_{S}(R), G)\) be the associated
Galois covering of \(S\) with group \(E\). Let \(W_0\) be the form of \(W_S\) associated to \(rev(G)\). There exists a morphism
\(S' \to S\), which is a principal homogeneous bundle under \(W_0\), such that \(G_{S'}\) is quasi-splittable (i.e.~possesses
a Borel subgroup).}

\emph{\SGAInlineRefTarget{sga3:xxiv:theorem:4-4-1:item:iii}(iii) Every \(S\)-reductive group of type \(R\) is split by a principal Galois covering \(\tilde{S} \to S\) with group \(A\).}

Let us first state:

\textbf{\SGAInlineRefTarget{sga3:xxiv:proposition:4-4-2}Proposition 4.4.2.} \emph{Let \(S\) be a scheme, \(G\) an \(S\)-reductive group, \(T\) a maximal torus of \(G\), \(P\) a principal
homogeneous bundle under \(G\), \(G'\) the twisted form of \(G\) associated to \(P\) (one then has, cf.~\SGARefLink{sga3:xxiv:proposition:4-2-1}{4.2.1}, a canonical
isomorphism \(P / Norm_{G}(T) \xrightarrow{\sim} Tor(G')\)). Let \(T'\) be a maximal torus of \(G'\), defining a composite
morphism}

\[
  S\longrightarrow\operatorname{Tor}(G')
  \xrightarrow{\sim}P/\operatorname{Norm}_G(T).
\]

\emph{Consider the canonical morphisms}

\[
  P\longrightarrow P/T\longrightarrow P/\operatorname{Norm}_G(T).
\]

\emph{and take their inverse images by the preceding morphism:}

\[
  % BEGIN inputs/XXIV/native_diagrams/exp24/003_inverse_image_grid.tex
\begin{tikzcd}[column sep=large,row sep=large]
P \arrow[r]
&
P/T \arrow[r]
&
P/\operatorname{Norm}_G(T)
\\
H \arrow[u] \arrow[r]
&
S' \arrow[u] \arrow[r]
&
S \arrow[u]
\end{tikzcd}

% END inputs/XXIV/native_diagrams/exp24/003_inverse_image_grid.tex

\]

\emph{Then \(S'\) (resp. \(H\)) is a principal homogeneous bundle over \(S\) (resp. \(S'\)) under \(W_{G}(T)\) (resp. \(T_{S'}\)).
Moreover, if one makes \(W_{G}(T)\) operate on \(T\) in the obvious way, the bundle associated to \(S'\) is isomorphic to
\(T'\).}

The first two assertions are trivial, the last is proved like the corresponding assertion of \SGARefLink{sga3:xxiv:2.6}{2.6}.

\textbf{\SGAInlineRefTarget{sga3:xxiv:corollary:4-4-3}Corollary 4.4.3.} \emph{Let \(S\) be a semi-local scheme, \(G\) an \(S\)-reductive group, \(T\) a maximal torus of \(G\). Suppose
one of the two following conditions is satisfied:}

\emph{\SGAInlineRefTarget{sga3:xxiv:corollary:4-4-3:item:i}(i) \(T\) is split.}

\emph{\SGAInlineRefTarget{sga3:xxiv:corollary:4-4-3:item:ii}(ii) \(T\) is contained in a Borel subgroup of \(G\), and \(G\) is either adjoint, or simply connected.}

\emph{Let moreover \(P\) be a principal homogeneous bundle under \(G\). There exists an \(S\)-scheme \(S'\), which is a principal
homogeneous bundle under \(W_{G}(T)\), such that \(P_{S'}\) is trivial.}

Indeed, if \(G'\) is the form of \(G\) associated to \(P\), then \(G'\) possesses a maximal torus \(T'\) (\SGARefLink{sga3:xiv:3.20}{Exp. XIV, 3.20}).
Resuming the notation of the preceding proposition, one sees that \(H^{1}(S', T_{S'}) = 0\) (by theorem 90 for \SGARefLink{sga3:xxiv:corollary:4-4-3:item:i}{(i)}, by
\SGARefLink{sga3:xxiv:corollary:3-14}{3.14} for \SGARefLink{sga3:xxiv:corollary:4-4-3:item:ii}{(ii)}). The morphism \(H \to S'\) possesses a section, so \(P_{S'}\) also possesses a section over \(S'\). QED

Let us now prove the theorem. Assertion \SGARefLink{sga3:xxiv:corollary:4-4-3:item:i}{(i)} is a particular case of the preceding corollary (take \(G =
\acute{E}p_{S}(R)\), endowed with its canonical split torus). Let us prove \SGARefLink{sga3:xxiv:corollary:4-4-3:item:ii}{(ii)}. One knows (\SGARefLink{sga3:xxiv:3.12}{3.12}) that \(G\) is an inner
twisted form of

\[ G_{0} = Q-\acute{E}p_{rev(G)/S}(R). \]

If \(T_0\) is the canonical maximal torus of \(G_0\), \(W_{G_{0}}(T_{0})\) is the group \(W_0\) described in the statement. The
form \(G\) of \(G_0\) corresponds to a principal homogeneous bundle \(P\) under \(ad(G_{0})\) (\(P = Isomint(G_{0}, G)\)). The
group \(W_{ad(G_{0})}(T^{ad}_{0})\) is canonically isomorphic to \(W_0\), and one obtains the wanted result by applying
\SGARefLink{sga3:xxiv:corollary:4-4-3}{4.4.3} to the situation \((ad(G_{0}), T^{ad}_{0}, P)\), hypothesis \SGARefLink{sga3:xxiv:corollary:4-4-3:item:ii}{(ii)} of \SGARefLink{sga3:xxiv:corollary:4-4-3}{4.4.3} being well verified. Let us finally prove
\SGARefLink{sga3:xxiv:theorem:4-4-1:item:iii}{(iii)}. Resume the notation of \SGARefLink{sga3:xxiv:corollary:4-4-3:item:ii}{(ii)}; one has a diagram

\[
  % BEGIN inputs/XXIV/native_diagrams/exp24/004_rev_base_triangle.tex
\begin{tikzcd}[column sep=large,row sep=large]
&
\operatorname{rev}(G) \arrow[d,"E_S"']
\\
S'
  \arrow[r,"W_0"']
&
S
\end{tikzcd}

% END inputs/XXIV/native_diagrams/exp24/004_rev_base_triangle.tex

\]

One knows that \(G_{S'}\) is isomorphic to \((G_{0})_{S'}\) and that \((G_{0})_{rev(G)}\) is splittable. If one sets
\(\tilde{S} = S' \times_{S} rev(G)\), \(G_{\tilde{S}}\) is splittable, and it only remains to verify that \(\tilde{S}\) is
indeed a principal Galois covering of \(S\) with group \(A\), which follows from the following more general lemma (naturally
valid in any site):

% END inputs/XXIV/tex/components/10_expose_XXIV_section4_through_theorem441_setup_en.tex

% BEGIN inputs/XXIV/tex/components/11_expose_XXIV_lemma444_through_section54_en.tex
% Authority coverage: Polo--Gille Exposé XXIV, local pages 29--32.
% Lead-reviewed against the authority PDF; the next source cursor is Proposition 5.5 on local page 32.


\textbf{\SGAInlineRefTarget{sga3:xxiv:lemma:4-4-4}Lemma 4.4.4.} \emph{Let \(S\) be a scheme, \(G\) and \(H\) two \(S\)-group schemes, \(G \to \operatorname{Aut}_{S-gr.}(H)\) an
action of \(G\) on \(H\), \(E\) a \(G\)-principal homogeneous bundle, \(F\) an \(H_0\)-principal homogeneous bundle, where \(H_0\) is
the form of \(H\) associated to \(E\). Then \(E \times_{S} F\) is endowed naturally with a structure of principal homogeneous
bundle under the semi-direct product \(H \cdot G\).}

Write \((e, g) \mapsto eg\) (resp. \((f, u) \mapsto f u\)) for the action (on the right) of \(G\) on \(E\) (resp. of \(H_0\) on
\(F\)). Write

\[
  p:E\times_S H\longrightarrow H_0.
\]

for the canonical projection (\(H_0\) is by definition the quotient of \(E \times_{S} H\) by \(G\) acting on it by the formula
\((e, h) g = (eg, g^{-1} h g)\)). Consider the morphism

\[
  r:E\times_S F\times_S H\times_S G\longrightarrow E\times_S F.
\]

defined set-theoretically by

\[
  r(e,f,h,g)=\bigl(eg,f\cdot p(e,h)\bigr).
\]

The morphism \(r\) indeed defines an action of the semi-direct product \(H \cdot G\) on \(E \times_{S} F\). Indeed, one has
set-theoretically

\[
  r\bigl(r(e,f,h,g),h',g'\bigr)
  =\bigl(egg',f\,p(e,h)\,p(eg,h')\bigr).
\]

but

\[
  p(e,h)p(eg,h')
  =p(e,h)p(e,gh'g^{-1})
  =p(e,hgh'g^{-1}).
\]

whence

\[
  r\bigl(r(e,f,h,g),h',g'\bigr)
  =r(e,f,hgh'g^{-1},gg').
\]

which had to be proved.

To prove now that this law is indeed a law of principal homogeneous bundle, one may assume that \(E\) and \(F\) are trivial,
in which case one sees at once that \(E \times_{S} F\) is also a trivial bundle.

\subsection*{\SGAInlineRefTarget{sga3:xxiv:subsection:4-5}4.5. Isotriviality of maximal tori of semisimple groups}
\footnote{\textit{Editors' note.} This no.~\SGARefLink{sga3:xxiv:subsection:4-5}{4.5} has been added.}

Let us bring out here the following result, contained implicitly in \SGARefLink{sga3:xxiv:2.6}{2.6} (cf.~the N.D.E. (13)) and which was mentioned at
the end of \SGARefLink{sga3:ix:comments:1:2}{Exp. IX 1.2}. Recall (\SGARefLink{sga3:XXIII:5.11}{XXIII 5.11}) that, for every scheme \(S\) and every reduced root datum \(R\), one writes
\(\acute{E}p_{S}(R)\) for the unique pinned \(S\)-group of type \(R\) (which exists by Exp. XXV) and \(T_{S}(R)\) its canonical
maximal torus.

\textbf{\SGAInlineRefTarget{sga3:xxiv:proposition:4-5-1}Proposition 4.5.1.} \emph{Let \(G\) be an \(S\)-semisimple group of constant type \(R\); assume that \(G\) possesses a maximal
torus \(T\) (which is the case if \(S\) is semi-local, by \SGARefLink{sga3:xiv:3.20}{XIV 3.20}).}

\emph{\SGAInlineRefTarget{sga3:xxiv:proposition:4-5-1:item:i}(i) Then \(P = Isom_{S-gr.}(\acute{E}p_{S}(R), T_{S}(R); G, T) / T_{S}(R)_{ad}\) is a principal \(S\)-bundle under the
finite group \(\operatorname{Aut}(R)\), and \(T_P\) is a split maximal torus of \(G_P\). Consequently, \(T\) is isotrivial.}

\emph{\SGAInlineRefTarget{sga3:xxiv:proposition:4-5-1:item:ii}(ii) Moreover, for every finite set \(F\) of points of \(P\) contained in an affine open set, there exists an open set \(U\)
of \(P\) containing \(F\) such that \(G_U\) is split.}

\emph{Proof.} The exact sequence

\[
  1\longrightarrow T_S(\mathscr R)_{\mathrm{ad}}
  \longrightarrow
  \underline{\operatorname{Aut}}_{S\text{-}\mathrm{gr.}}
  \bigl(\operatorname{\acute{E}p}_S(\mathscr R),T_S(\mathscr R)\bigr)
  \longrightarrow
  \underline{\operatorname{Aut}}_{S\text{-}\mathrm{gr.}}\bigl(T_S(\mathscr R)\bigr).
\]

induces for every \(S' \to S\) a morphism of bundles \(P_{S'} \to Isom_{S'-gr.}(T_{S'}(R), T_{S'})\), so if \(P_{S'}\)
possesses a section, then \(T_{S'}(R)\) and \(T_{S'}\) are isomorphic. This is in particular the case for \(S' = P\).

Note moreover that \(\operatorname{Aut}_{S-gr.}(\acute{E}p_{S}(R), T_{S}(R)) / T_{S}(R)_{ad}\) is none other than the
constant \(S\)-group \(\operatorname{Aut}(R)_{S}\), and that \(P\) is identified with the \(\operatorname{Aut}(R)_{S}\)-torsor
\(Isom(R, \tilde{R}(G, T))\), where \(\tilde{R}(G, T)\) denotes the twisted root datum associated to \((G, T)\) (cf.~\SGARefLink{sga3:xxii:prop:1.10}{XXII
1.10}).

Then \SGARefLink{sga3:xxiv:proposition:4-5-1:item:ii}{(ii)} follows from \SGARefLink{sga3:xxiv:proposition:4-5-1:item:i}{(i)} and from \SGARefLink{sga3:xxiv:2.14}{2.14}, taking into account the fact that the \(g_{\alpha}\) (\(\alpha\) ranging over the
roots of \(R\)) are free \(\mathcal O_U\)-modules of rank 1 on a sufficiently small affine open set \(U\) containing \(F\).

\section*{\SGAInlineRefTarget{sga3:xxiv:section:5}5. Canonical decomposition of an adjoint or simply connected group}

In this section, we shall use the results of no 1 to generalize to the case of schemes a classical decomposition of
adjoint (resp. simply connected) groups. So as not to overburden the exposition indefinitely, the proofs are sketched
and the details are left to the reader; in fact it is always a matter of absolutely standard proofs in the theory of
principal bundles: reduction of the structure group, twisting, etc.

\subsection*{\SGAInlineRefTarget{sga3:xxiv:subsection:5-1}5.1.}

Recall (\SGARefLink{sga3:XXI:7.4}{Exp. XXI, 7.4}) that a Dynkin diagram is a disjoint union of its connected components, which are Dynkin diagrams.
Moreover, every non-empty connected Dynkin diagram corresponding to a root datum is isomorphic to one of the standard
diagrams (\(A_{n}\), \(B_{n}\), \ldots, \(G_2\)) which were exhibited in \SGARefLink{sga3:XXI:7.4.6}{Exp. XXI, 7.4.6}. In the sequel, we shall be interested
only in Dynkin diagrams whose connected components are of one of the preceding types. Let \(T\) be the set of these
standard diagrams. For every Dynkin diagram \(D\), let \(n(t) = n_{D}(t)\) be the number of connected components of \(D\)
isomorphic to \(t\), where \(t \in T\). The \emph{type} of \(D\) is by definition \(\sum_{t \in T} n_{D}(t) \cdot t\).

A Dynkin diagram of type \(t\) is said to be \emph{simple of type \(t\)}, a Dynkin diagram of type \(n \cdot t\) is said to be
\emph{isotypic of type \(t\)}. Let \(D_0\) be the set of connected components of \(D\) and let

\[ a : D_{0} \to T \]

be the obvious map. The type of \(D\) is none other than \(\sum_{x \in D_{0}} a(x)\).

\subsection*{\SGAInlineRefTarget{sga3:xxiv:subsection:5-2}5.2.}

Let \(S\) be a scheme, \(D\) an \(S\)-Dynkin scheme (verifying the restrictive condition stated above). The cokernel of the
pair of morphisms \(L \Rightarrow D\) (\(L =\)\footnote{N.D.E.: We have replaced \(R\) by \(L\), as in \SGARefLink{sga3:xxiv:3.2}{3.2}.} scheme of bonds of \(D\)) is written \(D_0\). This is the ``scheme
of connected components'' of \(D\) (it

exists trivially when \(D\) is constant; the general case is deduced from this by descent); it is a finite twisted
constant \(S\)-scheme. One has a canonical morphism

\[ a : D_{0} \to T_{S}. \]

For \(t \in T\), set \(a^{-1}(t) = D_{0, t}\); this is a subscheme of \(D_0\), whose inverse image in \(D\), written \(D_{t}\), is
the \emph{isotypic component of type \(t\)} of the Dynkin scheme \(D\). Each \(D_{t}\) is a subscheme of \(D\), and one has

\[
  D=\coprod_{t\in T}D_t.
\]

Note that the degree of \(D_{0, t}\) at \(s \in S\) is \(n(s, t)\), if the type of \(D\) at \(s\) is \(\sum_{t} n(s, t) \cdot t\).

\subsection*{\SGAInlineRefTarget{sga3:xxiv:subsection:5-3}5.3.}

In what follows, we shall consider only semisimple adjoint (resp. simply connected) groups. To simplify the language, we
shall state the results for adjoint groups; all statements will remain valid if one substitutes everywhere ``simply
connected'' for ``adjoint''.

Recall that an adjoint reduced root datum is determined up to isomorphism by the type of its Dynkin diagrams. We shall
therefore say that an adjoint root datum \(R\) (resp. a semisimple adjoint group \(G\)) is \emph{of type \(\sum n(t) \cdot t\)} if
its Dynkin diagrams are (resp. if its type is given by an adjoint root datum of type \(\sum n(t) \cdot t\)). We shall say
that \(R\) or \(G\) is \emph{simple of type \(t\)} (resp. \emph{isotypic of type \(t\)}) if its type is \(t\) (resp. \(n \cdot t\), \(n > 0\)).

If \(G\) is a semisimple adjoint group, we shall use the symbols \(Dyn_{0}(G)\) and \(Dyn_{0, t}(G)\) in the sense defined in
\SGARefLink{sga3:xxiv:subsection:5-2}{5.2}.

\subsection*{\SGAInlineRefTarget{sga3:xxiv:subsection:5-4}5.4.}

Let \(t_i\), \(i=1,2,\ldots,n\), be distinct elements of \(T\), and let \(R_i\) be a pinned adjoint root datum isotypic of
type \(t_{i}\). Consider the product \(R = R_{1} \times \cdot\cdot\cdot \times R_{n}\) (\SGARefLink{sga3:XXI:6.4.1}{Exp. XXI, 6.4.1}). Let \(S\) be a
scheme such that the different \(\acute{E}p_{S}(R_{i})\) exist (cf.~Exp. XXV). One verifies at once that there exists a
canonical isomorphism

\[
  \tag{*}\label{sga3:xxiv:subsection:5-4:equation:star}
  \operatorname{\acute{E}p}_S(\mathscr R)
  =
  \operatorname{\acute{E}p}_S(\mathscr R_1)
  \times_S\cdots\times_S
  \operatorname{\acute{E}p}_S(\mathscr R_n).
\]



% END inputs/XXIV/tex/components/11_expose_XXIV_lemma444_through_section54_en.tex

% BEGIN inputs/XXIV/tex/components/12_expose_XXIV_proposition55_through_remark631_en.tex
% Authority coverage: Polo--Gille Exposé XXIV, local pages 32--37.
% Reconstructed manually from the authority PDF.  The Reinhold Markdown
% lineage was used only as a declared comparison translation.

\par\medskip\noindent\SGAInlineRefTarget{sga3:xxiv:proposition:5-5}\textbf{Proposition 5.5.---}\enspace
\emph{Let \(S\) be a scheme and \(G\) an adjoint semisimple
\(S\)-group.  There is a unique decomposition}
\[
  G\simeq\prod_{t\in\mathrm T}G_t,
\]
\emph{where \(G_t\) is an adjoint semisimple \(S\)-group, isotypic of
type \(t\).  Moreover, the preceding decomposition induces
isomorphisms}
\[
  \operatorname{Dyn}_t(G)\simeq\operatorname{Dyn}(G_t),
  \qquad
  \underline{\operatorname{Aut}}_{S\text{-}\mathrm{gr.}}(G)
  \simeq
  \prod_{t\in\mathrm T}
  \underline{\operatorname{Aut}}_{S\text{-}\mathrm{gr.}}(G_t).
\]

This was proved above when \(G\) is split.  In the general case one may
suppose that \(G\) has constant type \(\mathscr R\).  Using the
preceding decomposition of \(A_S(\mathscr R)\) and \SGARefLink{sga3:xxiv:1.17}{1.17}, one obtains
the desired decomposition of \(G\).  The other assertions then follow
by descent.

\par\medskip\noindent\SGAInlineRefTarget{sga3:xxiv:remark:5-6}\textbf{Remark 5.6.---}\enspace
More generally, if \(G\) and \(H\) are two adjoint semisimple groups,
there are canonical isomorphisms in the following commutative diagram:
\[
  % BEGIN inputs/XXIV/native_diagrams/exp24/012_isomorphism_dynkin_comparison.tex
\begin{tikzcd}[column sep=huge,row sep=large]
\underline{\operatorname{Isom}}_{S\text{-}\mathrm{gr.}}(G,H)
  \arrow[r,"\sim"] \arrow[d]
&
\displaystyle\prod_{t\in\mathrm T}
\underline{\operatorname{Isom}}_{S\text{-}\mathrm{gr.}}(G_t,H_t)
  \arrow[d]
\\
\underline{\operatorname{Isomext}}(G,H)
  \arrow[r,"\sim"] \arrow[d,"\wr"']
&
\displaystyle\prod_{t\in\mathrm T}
\underline{\operatorname{Isomext}}(G_t,H_t)
  \arrow[d,"\wr"]
\\
\underline{\operatorname{Isom}}_{\mathrm{Dyn}}
  \bigl(\operatorname{Dyn}(G),\operatorname{Dyn}(H)\bigr)
  \arrow[r,"\sim"]
&
\displaystyle\prod_{t\in\mathrm T}
\underline{\operatorname{Isom}}_{\mathrm{Dyn}}
  \bigl(\operatorname{Dyn}(G_t),\operatorname{Dyn}(H_t)\bigr).
\end{tikzcd}

% END inputs/XXIV/native_diagrams/exp24/012_isomorphism_dynkin_comparison.tex

\]

\par\medskip\noindent\SGAInlineRefTarget{sga3:xxiv:remark:5-7}\textbf{Remark 5.7.---}\enspace
One can give the following intrinsic characterization of \(G_t\),
which we state without proof: it is the largest reductive subgroup of
\(G\) that is invariant (and indeed characteristic) and isotypic of
type \(t\).

The preceding proposition reduces the study of adjoint semisimple
groups to that of isotypic adjoint semisimple groups.  We now undertake
this study.

\subsection*{\SGAInlineRefTarget{sga3:xxiv:subsection:5-8}5.8.}

Let \(\mathscr R\) be a pinned reduced adjoint simple root datum of type
\(t\), let \(I\) be a nonempty finite set, and let
\(\mathscr R^I\) be the product root datum of copies
\(\mathscr R_i\) of \(\mathscr R\), for \(i\in I\).  Denote by \(E\)
the group of automorphisms of \(\mathscr R\), identified with the group
of automorphisms of its Dynkin diagram \(D\).  The Dynkin diagram of
\(\mathscr R^I\) is identified with \(D^I\), and hence there is an
exact sequence
\[
  1\longrightarrow\operatorname{Aut}(D)^I
  \longrightarrow\operatorname{Aut}(D^I)
  \longrightarrow\operatorname{Aut}(I)
  \longrightarrow1,
\]
where \(\operatorname{Aut}(I)\) is the group of permutations of \(I\).
It follows that the canonical isomorphism
\[
  \operatorname{\acute Ep}_S(\mathscr R)^I
  \simeq
  \operatorname{\acute Ep}_S(\mathscr R^I)
\]
induces an exact sequence
\[
  1\longrightarrow A_S(\mathscr R)^I
  \longrightarrow A_S(\mathscr R^I)
  \longrightarrow\underline{\operatorname{Aut}}(I_S)
  \longrightarrow1.
\]
The last group is the constant \(S\)-group associated with
\(\operatorname{Aut}(I)\).  Notice also that \(I\) is canonically
identified with the set of connected components of \(D^I\).

Let \(G\) be an \(S\)-semisimple group of type \(\mathscr R^I\), thus
defining (cf.\ \SGARefLink{sga3:xxiv:1.17}{1.17}) a principal homogeneous bundle \(P\) under
\(A_S(\mathscr R^I)\).  The definitions by descent of
\(\operatorname{Dyn}(G)\) (\SGARefLink{sga3:xxiv:3.7}{3.7}) and of
\(\operatorname{Dyn}_0(G)\) (\SGARefLink{sga3:xxiv:subsection:5-2}{5.2}) show that the bundle associated with
\(P\) by
\[
  A_S(\mathscr R^I)\longrightarrow
  \underline{\operatorname{Aut}}(I_S)
\]
is precisely
\(\underline{\operatorname{Isom}}_S(I_S,\operatorname{Dyn}_0(G))\),
corresponding to the form \(\operatorname{Dyn}_0(G)\) of \(I_S\).
Using again the equivalence of categories \SGARefLink{sga3:xxiv:1.17}{1.17} and the preceding exact
sequence, one deduces formally that there is a reductive
\(\operatorname{Dyn}_0(G)\)-group \(G_0\) of type \(\mathscr R\), and
an \(S\)-isomorphism
\[
  G\simeq\prod_{\operatorname{Dyn}_0(G)/S}G_0.
\]

\par\medskip\noindent\SGAInlineRefTarget{sga3:xxiv:proposition:5-9}\textbf{Proposition 5.9.---}\enspace
\emph{Let \(S\) be a scheme and \(G\) an adjoint semisimple
\(S\)-group isotypic of type \(t\).  There exist an adjoint semisimple
\(\operatorname{Dyn}_0(G)\)-group \(G_0\), simple of type \(t\), and a
unique \(S\)-isomorphism}
\[
  G\simeq\prod_{\operatorname{Dyn}_0(G)/S}G_0.
\]
\emph{Moreover, this isomorphism induces an exact sequence}
\[
  1\longrightarrow
  \prod_{\operatorname{Dyn}_0(G)/S}
    \underline{\operatorname{Aut}}_{\mathrm{gr.}}(G_0)
  \longrightarrow
  \underline{\operatorname{Aut}}_{S\text{-}\mathrm{gr.}}(G)
  \longrightarrow
  \underline{\operatorname{Aut}}_S\bigl(\operatorname{Dyn}_0(G)\bigr)
  \longrightarrow1.
\]

One may plainly suppose that \(G\) has constant type \(n\,t\).  The
first assertion has already been proved (and uniqueness is clear).
The second then follows from the split case by descent.

Combining \SGARefLink{sga3:xxiv:proposition:5-5}{5.5} and \SGARefLink{sga3:xxiv:proposition:5-9}{5.9} gives:

\par\medskip\noindent\SGAInlineRefTarget{sga3:xxiv:proposition:5-10}\textbf{Proposition 5.10.---}\enspace
\emph{Let \(S\) be a scheme, let \(G\) be an adjoint semisimple
\(S\)-group, and put \(\mathfrak D=\operatorname{Dyn}(G)\).}

\emph{(i) There is a canonical decomposition}
\[
  G\simeq
  \prod_{t\in\mathrm T}
  \prod_{\mathfrak D_{0,t}/S}G_{0,t}
  \simeq
  \prod_{\operatorname{Dyn}_0(G)/S}G_0,
\]
\emph{where each \(G_{0,t}\) is a simple adjoint
\(\operatorname{Dyn}_{0,t}(G)\)-group of type \(t\); equivalently,
\(G_0\) is an adjoint semisimple
\(\operatorname{Dyn}_0(G)\)-group whose type at
\(x\in\operatorname{Dyn}_0(G)\) is \(a(x)\in\mathrm T\), where
\(a:\operatorname{Dyn}_0(G)\to\mathrm T_S\) was defined in \SGARefLink{sga3:xxiv:subsection:5-2}{5.2}.}

\emph{(ii) The preceding decompositions induce isomorphic exact
sequences, written vertically below.  Here
\(\underline{\operatorname{Aut}}_{S,a}(\mathfrak D_0)\) denotes the
scheme of automorphisms of \(\operatorname{Dyn}_0(G)\) commuting with
\(a:\operatorname{Dyn}_0(G)\to\mathrm T_S\):}
\[
  % BEGIN inputs/XXIV/native_diagrams/exp24/013_automorphism_exact_columns.tex
\begin{tikzcd}[column sep=large,row sep=large]
1 \arrow[d] & 1 \arrow[d] \\
\displaystyle
\prod_{t\in\mathrm T}\prod_{\mathfrak D_{0,t}/S}
  \underline{\operatorname{Aut}}_{\mathrm{gr.}}(G_{0,t})
  \arrow[r,"\sim"] \arrow[d]
&
\displaystyle
\prod_{\operatorname{Dyn}_0(G)/S}
  \underline{\operatorname{Aut}}_{\mathrm{gr.}}(G_0)
  \arrow[d]
\\
\underline{\operatorname{Aut}}_{S\text{-}\mathrm{gr.}}(G)
  \arrow[r,equal] \arrow[d]
&
\underline{\operatorname{Aut}}_{S\text{-}\mathrm{gr.}}(G)
  \arrow[d]
\\
\displaystyle\prod_{t\in\mathrm T}
  \underline{\operatorname{Aut}}_S
    \bigl(\operatorname{Dyn}_{0,t}(G)\bigr)
  \arrow[r,"\sim"] \arrow[d]
&
\underline{\operatorname{Aut}}_{S,a}
  \bigl(\operatorname{Dyn}_0(G)\bigr)
  \arrow[d]
\\
1 & 1
\end{tikzcd}

% END inputs/XXIV/native_diagrams/exp24/013_automorphism_exact_columns.tex

\]

\par\medskip\noindent\SGAInlineRefTarget{sga3:xxiv:corollary:5-11}\textbf{Corollary 5.11.---}\enspace
\emph{Under the preceding hypotheses, the following three categories
are equivalent:}
\begin{enumerate}
\item[(i)] the category of principal homogeneous bundles under \(G\);
\item[(ii)] the category of principal homogeneous bundles under \(G_0\);
\item[(iii)] the product, for \(t\in\mathrm T\), of the categories of
principal homogeneous bundles under the \(G_{0,t}\).
\end{enumerate}
This follows formally from the preceding decompositions and \SGARefLink{sga3:xxiv:proposition:8-4}{8.4}.

\par\medskip\noindent\SGAInlineRefTarget{sga3:xxiv:corollary:5-12}\textbf{Corollary 5.12.---}\enspace
\emph{There are canonical isomorphisms}
\[
  \underline{\operatorname{Tor}}(G)
  \simeq
  \prod_{t\in\mathrm T}
  \prod_{\mathfrak D_{0,t}/S}
    \underline{\operatorname{Tor}}(G_{0,t})
  \simeq
  \prod_{\operatorname{Dyn}_0(G)/S}
    \underline{\operatorname{Tor}}(G_0),
\]
\emph{and likewise with \(\underline{\operatorname{Tor}}\) replaced
by \(\underline{\operatorname{Bor}}\), respectively by
\(\underline{\operatorname{Kil}}\).}

This is immediate from the split case.

\par\medskip\noindent\SGAInlineRefTarget{sga3:xxiv:remark:5-13}\textbf{Remark 5.13.---}\enspace
The canonical morphism
\[
  \operatorname{Dyn}(G)\longrightarrow\operatorname{Dyn}_0(G)
\]
allows one to regard \(\operatorname{Dyn}(G)\) as a Dynkin scheme over
\(\operatorname{Dyn}_0(G)\); in fact, it is the Dynkin scheme
\(\operatorname{Dyn}(G_0)\) of \(G_0\).

Likewise, if \(T\subset B\) is a Killing couple of \(G\), corresponding
canonically to a Killing couple \(T_0\subset B_0\) of \(G_0\), then the
obstructions to the quasi-splitting of \(G\) and \(G_0\), which lie by
\SGARefLink{sga3:xxiv:3.9}{3.9} in
\[
  \operatorname{Pic}\bigl(\operatorname{Dyn}(G)\bigr)
  =
  \operatorname{Pic}\bigl(\operatorname{Dyn}(G_0)\bigr),
\]
coincide.  Hence:

\par\medskip\noindent\SGAInlineRefTarget{sga3:xxiv:corollary:5-14}\textbf{Corollary 5.14.---}\enspace
\emph{The following conditions are equivalent:}
\begin{enumerate}
\item[(i)] \(G\) is quasi-splittable;
\item[(ii)] \(G_0\) is quasi-splittable;
\item[(iii)] each \(G_{0,t}\), \(t\in\mathrm T\), is quasi-splittable.
\end{enumerate}

\section*{\SGAInlineRefTarget{sga3:xxiv:section:6}6. Automorphisms of Borel Subgroups of Reductive Groups}

\par\medskip\noindent\SGAInlineRefTarget{sga3:xxiv:lemma:6-1}\textbf{Lemma 6.1.---}\enspace
\emph{Let \(S\) be a scheme, let \((G,T,M,R)\) be a split \(S\)-group,
let \(\Delta\) be a system of simple roots of \(R\), and let \(B\) be
the corresponding Borel subgroup.  Then \(B^u\) is generated, as an
\textup{(fppf)} sheaf of groups, by the \(U_\alpha\),
\(\alpha\in\Delta\).}

Let \(H\) be the subgroup sheaf of \(B^u\) generated by the
\(U_\alpha\), \(\alpha\in\Delta\).  We prove
\(H\supset U_\beta\), for \(\beta\in R_+\), by induction on
\(\operatorname{ord}(\beta)=\operatorname{ord}_\Delta(\beta)>0\)
(cf.\ \SGARefLink{sga3:XXI:3.2.15}{Exp.~XXI, 3.2.15}).  The assertion holds by hypothesis when
\(\operatorname{ord}(\beta)=1\).  Suppose
\(\operatorname{ord}(\beta)>1\), and suppose
\(U_\gamma\subset H\) whenever
\(\operatorname{ord}(\gamma)<\operatorname{ord}(\beta)\).  There is an
\(\alpha\in\Delta\) such that \(\beta-\alpha\in R_+\)
(\SGARefLink{sga3:XXI:3.1.1}{Exp.~XXI, 3.1.1}).  Let \(p\) be the largest integer such that
\(\beta-p\alpha=\beta'\in R_+\).  Then
\(U_\alpha,U_{\beta'}\subset H\) and \(\beta'-\alpha\notin R\).
Thus it remains to prove:

\par\medskip\noindent\SGAInlineRefTarget{sga3:xxiv:lemma:6-2}\textbf{Lemma 6.2.---}\enspace
\emph{Keep the notation of \SGARefLink{sga3:XXIII:6.4}{Exp.~XXIII, 6.4}.  Suppose \(p=1\), that is,
\(\beta-\alpha\) is not a root.  If \(H\) is a subgroup sheaf of
\(B^u\) such that \(U_\alpha,U_\beta\subset H\), then
\(U_{i\alpha+j\beta}\subset H\) whenever
\(i\alpha+j\beta\in R\), \(i>0\), \(j>0\).}

Distinguish four cases according as \(q=0,1,2,3\).  In what follows,
\(x\) and \(y\) denote arbitrary sections of
\(\mathbb G_{a,S'}\), where \(S'\to S\).

For \(q=0\), there is nothing to prove.

For \(q=1\), one has
\[
 p_{\alpha+\beta}(\pm x)
 =p_\beta(-y)p_\alpha(-x)p_\beta(y)p_\alpha(x)\in H(S'),
\]
and hence \(U_{\alpha+\beta}\subset H\).

For \(q=2\), one similarly has
\[
 p_{\alpha+\beta}(\pm xy)
 p_{2\alpha+\beta}(\pm x^2y)
 =p_\beta(-y)p_\alpha(-x)p_\beta(y)p_\alpha(x)\in H(S').
\]
Thus, after changing some signs,
\[
 F(x,y)=p_{\alpha+\beta}(xy)p_{2\alpha+\beta}(x^2y)\in H(S').
\]
Since \(p_{\alpha+\beta}\) and \(p_{2\alpha+\beta}\) commute,
\[
 F(x,1)F(1,-x)=p_{2\alpha+\beta}(x^2-x).
\]
\setcounter{footnote}{27}
\footnote{\textit{Editors' note.} We corrected \(F(-1,x)\) to
\(F(1,-x)\).}
If \(a\in\mathbb G_a(S)\), the equation \(X^2-X=a\) defines a free
surjective extension of \(S\), namely
\(\operatorname{Spec}\mathcal O_S[X]/(X^2-X-a)\).  Consequently
\(p_{2\alpha+\beta}(a)\in H(S')\), so
\(U_{2\alpha+\beta}\subset H\), and then
\(U_{\alpha+\beta}\subset H\), since \(F(1,a)\in H(S')\).

For \(q=3\), one similarly has
\[
 F(x,y)=
 p_{\alpha+\beta}(xy)
 p_{2\alpha+\beta}(x^2y)
 p_{3\alpha+\beta}(x^3y)
 p_{3\alpha+2\beta}(x^3y^2)\in H(S'),
\]
and, since
\[
 p_{3\alpha+2\beta}(\pm x)
 =F(1,1)^{-1}p_\beta(-x)F(1,1)p_\beta(x)\in H(S'),
\]
one obtains \(U_{3\alpha+2\beta}\subset H\), and therefore
\[
 K(x,y)=
 p_{\alpha+\beta}(xy)
 p_{2\alpha+\beta}(x^2y)
 p_{3\alpha+\beta}(x^3y)\in H(S').
\]
Computing
\[
 K(x+y,1)K(-x,1)K(1,y)^{-1}
\]
modulo \(U_{3\alpha+2\beta}\), one finds
\[
 p_{2\alpha+\beta}(2x^2+y^2+2xy-y)
 p_{3\alpha+\beta}(y^3+3xy^2+3x^2y-y)\in H(S').
\]
If \(a\in\mathbb G_a(S)\), the ``equations''
\[
  x^2=-xy-y+1-a,\qquad
  y^2=3y-2+3a
\]
define a free surjective extension of \(S\).  One then has
\(y^3+3xy^2+3x^2y-y=0\), and the preceding expression gives
\(p_{2\alpha+\beta}(a)\in H(S')\).
\footnote{\textit{Editors' note.} We corrected
\(p_{3\alpha+\beta}(a)\) to \(p_{2\alpha+\beta}(a)\) and supplied the
end of the argument.}
Thus \(H\) contains \(U_{2\alpha+\beta}\) and
\(U_{3\alpha+2\beta}\), and
\[
 E(x,y)=p_{\alpha+\beta}(xy)p_{3\alpha+\beta}(x^3y)\in H(S').
\]
Since \(p_{\alpha+\beta}(xy)\) and \(p_{3\alpha+\beta}\) commute, the
preceding computation applies:
\[
 E(1,x)E(1,-x)=p_{3\alpha+\beta}(x^3-x).
\]
Hence \(U_{3\alpha+\beta}\subset H\), and then
\(U_{\alpha+\beta}\subset H\).

\par\medskip\noindent\SGAInlineRefTarget{sga3:xxiv:remark:6-2-1}\textbf{Remark 6.2.1.---}\enspace
The preceding proof shows that the hypothesis
``\(H\) contains \(U_\alpha\) and \(U_\beta\)'' could have been
replaced by: ``\(H\) contains \(U_\alpha\) or \(U_\beta\), and \(H\)
is invariant under \(\operatorname{int}(U_\alpha)\) and
\(\operatorname{int}(U_\beta)\).''

\par\medskip\noindent\SGAInlineRefTarget{sga3:xxiv:theorem:6-3}\textbf{Theorem 6.3.---}\enspace
\emph{Let \(S\) be a scheme, let \(G\) and \(G'\) be two semisimple
\(S\)-groups, and let \(B\) (respectively \(B'\)) be a Borel subgroup
of \(G\) (respectively \(G'\)).  Every isomorphism
\(B\xrightarrow{\sim}B'\) is induced by a unique isomorphism
\(G\xrightarrow{\sim}G'\).}

The assertion is local for the étale topology, so we may suppose that
\(G\) is split: \(G=(G,T,M,R)\), with \(B\) defined by the system
\(R_+\) of positive roots of \(R\).  The given isomorphism
\(u:B\xrightarrow{\sim}B'\) induces an isomorphism from \(T\) onto a
maximal torus \(T'\) of \(B'\), hence of \(G'\).  The given
isomorphism \(T\simeq D_S(M)\) gives an isomorphism
\(T'\simeq D_S(M)\) in which the elements of \(R_+\) become the
constant roots of \(B'\) relative to \(T'\), and hence the elements of
\(R\) become the constant roots of \(G'\) relative to \(T'\).  Since
\(G\) and \(G'\) are semisimple, the coroots are determined by duality
(\SGARefLink{sga3:XXI:1.2.5}{Exp.~XXI, 1.2.5}); this proves that \((T',M,R)\) is a splitting of
\(G'\) such that \(\mathscr R(G)=\mathscr R(G')\).

Applying \SGARefLink{sga3:XXIII:5.1}{Exp.~XXIII, 5.1} (the uniqueness theorem), one obtains a unique
isomorphism \(G\xrightarrow{\sim}G'\) agreeing with \(u\) on \(T\)
and on the \(U_\alpha\), \(\alpha\in\Delta\).  By \SGARefLink{sga3:xxiv:lemma:6-1}{6.1}, its restriction
to \(B\) is \(u\).

\par\medskip\noindent\SGAInlineRefTarget{sga3:xxiv:remark:6-3-1}\textbf{Remark 6.3.1.---}\enspace
(i) Using \SGARefLink{sga3:xxii:prop:4.1.9}{Exp.~XXII, 4.1.9} and arguing as in \emph{loc.\ cit.}, 4.2.12,
one may replace ``isomorphism'' in the theorem by ``isogeny''
(respectively, ``central isogeny'').

(ii) The theorem is false for reductive groups.  For example, take
\(G=G'=\mathrm{SL}_2\times\mathbb G_m\), identified with the group
\[
 \left\{
 \begin{pmatrix}
   a&b&0\\
   c&d&0\\
   0&0&h
 \end{pmatrix}
 \ \middle|\ ad-bc=1,\ h\text{ invertible}
 \right\}.
\]

% END inputs/XXIV/tex/components/12_expose_XXIV_proposition55_through_remark631_en.tex

% BEGIN inputs/XXIV/tex/components/13_expose_XXIV_remark631_continuation_through_bibliography_en.tex
% Exposé XXIV, authority local pp. 38--54 / combined-reader pp. 1283--1299.
% Lead-audited against the Polo--Gille authority at 2400 dpi.
% Reinhold English Markdown was comparison-only and is credited in STATUS.
% Provisional authority images preserve diagram content pending native reconstruction.

\setcounter{footnote}{29}

take \(B=B'\) to be the Borel subgroup defined by \(c=0\), and let \(u\)
be the following automorphism of \(B\):
\[
  u(a,b,d,h)=(a,b,d,ah).
\]

\par\medskip\noindent\SGAInlineRefTarget{sga3:xxiv:corollary:6-4}\textbf{Corollary 6.4.---}\enspace
\emph{The functor
\[
  (G,B)\longmapsto B
\]
from the category of pairs \((G,B)\), where \(G\) is an
\(S\)-semisimple group and \(B\) is a Borel subgroup of \(G\), to the
category of \(S\)-group schemes, with isomorphisms as morphisms, is
fully faithful.}

\par\medskip\noindent\SGAInlineRefTarget{sga3:xxiv:corollary:6-5}\textbf{Corollary 6.5.---}\enspace
\emph{Let \(S\) be a scheme, \(G\) an \(S\)-semisimple group, \(B\) a
Borel subgroup of \(G\), and \(B'\) an \(S\)-group locally isomorphic to
\(B\) for the \textup{(fpqc)} topology. Then \(B'\) is locally
isomorphic to \(B\) for the local finite étale
topology,\footnote{\textit{Editors' note.} This topology is denoted
\textup{(étf)} in \SGARefLink{sga3:iv:subsection:6:3}{Exposé IV, 6.3}. In other words, for every \(s\in S\)
there exist an open neighborhood \(U\) of \(s\) and a finite surjective
étale morphism \(V\to U\) such that
\(B'\times_S V\simeq B\times_S V\).} and there exists an
\(S\)-semisimple group \(G'\) of which \(B'\) is a Borel subgroup.
Moreover, \(G'\) is unique up to a unique isomorphism inducing the
identity on \(B'\).}

\par\medskip\noindent\SGAInlineRefTarget{sga3:xxiv:corollary:6-6}\textbf{Corollary 6.6.---}\enspace
\emph{Let \(S\) be a scheme, \(G\) an \(S\)-semisimple group, and \(B\)
a Borel subgroup of \(G\). Then
\(\underline{\operatorname{Aut}}_{S\text{-}\mathrm{gr.}}(B)\) is
representable by an affine smooth \(S\)-scheme,
\(\underline{\operatorname{Autext}}(B)\) is representable by an étale
finite \(S\)-scheme, and the evident morphisms induce an isomorphism of
exact sequences.}

\[
  % BEGIN inputs/XXIV/native_diagrams/exp24/005_borel_automorphism_exact_sequences.tex
\begin{tikzcd}[column sep=small,row sep=large]
1 \arrow[r]
&
\underline{B}_{\mathrm{ad}}(G)
  \arrow[r]
  \arrow[d,"\wr" description]
&
\underline{\Aut}_{S\text{-}\mathrm{gr.}}(G,B)
  \arrow[r]
  \arrow[d,"\wr" description]
&
\underline{\Autext}(G)
  \arrow[r]
  \arrow[d,"\wr" description]
&
1
\\
1 \arrow[r]
&
\underline{B}_{\mathrm{ad}}
  \arrow[r]
&
\underline{\Aut}_{S\text{-}\mathrm{gr.}}(G)
  \arrow[r]
&
\underline{\Autext}(B)
  \arrow[r]
&
1 .
\end{tikzcd}

% END inputs/XXIV/native_diagrams/exp24/005_borel_automorphism_exact_sequences.tex

\]

This follows immediately from \SGARefLink{sga3:xxiv:2.1}{2.1} and the preceding results. We leave
to the reader the task of developing the analogues of the results of
nos.~1, 2, 3, and 4 in the foregoing setting.

\par\medskip\noindent\SGAInlineRefTarget{sga3:xxiv:remark:6-7}\textbf{Remark 6.7.---}\enspace
If \(S\) is a scheme and \(B\) an \(S\)-group, then \(B\) can be a
Borel subgroup of only one semisimple group, uniquely determined by
\(B\). It remains to characterize those \(S\)-groups \(B\) which are
Borel subgroups of semisimple
groups.\footnote{\textit{Editors' note.} One may ask whether every
smooth affine \(S\)-group whose geometric fibers are Borel subgroups of
semisimple groups is itself a Borel subgroup of a semisimple
\(S\)-group. We have suppressed the counterexample in the original,
with \(S\) the scheme of dual numbers over a field \(k\), because it
was erroneous, as M.~Demazure pointed out.}

\section*{\SGAInlineRefTarget{sga3:xxiv:section:7}7. Representability of the functors
\(\underline{\operatorname{Hom}}_{S\text{-}\mathrm{gr.}}(G,H)\), for
\(G\) reductive}

\subsection*{\SGAInlineRefTarget{sga3:xxiv:subsection:7-1}7.1. The split case}

\subsubsection*{7.1.1.}

Let \(S\) be a scheme, \(G\) a pinned \(S\)-group, \(T\) its maximal
torus, and \(\Delta\) the set of simple roots. For
\(\alpha\in\Delta\), let
\[
  u_\alpha\in U_\alpha^\times(S),
  \qquad
  w_\alpha\in\operatorname{Norm}_G(T)(S)
\]
be the elements defined by the pinning.

Let \(H\) be an \(S\)-group scheme. We are interested in the functor
\(\underline{\operatorname{Hom}}_{S\text{-}\mathrm{gr.}}(G,H)\), and
more precisely in the morphism
\[
  q:\underline{\operatorname{Hom}}_{S\text{-}\mathrm{gr.}}(G,H)
  \longrightarrow
  \underline{\operatorname{Hom}}_{S\text{-}\mathrm{gr.}}(T,H).
\]
Let
\[
  f_T:T\longrightarrow H
\]
be an \(S\)-group homomorphism and put \(q^{-1}(f_T)=\mathscr F\).
Thus
\[
  \mathscr F(S')
  =
  \left\{
    f\in\operatorname{Hom}_{S'\text{-}\mathrm{gr.}}
      (G_{S'},H_{S'})
    \,\middle|\,
    f|_{T_{S'}}=(f_T)_{S'}
  \right\}.
\]
There is a morphism of \(S\)-functors
\[
  i:\mathscr F\longrightarrow H^{2n},
\]
where \(n=\operatorname{Card}(\Delta)\), defined on points by
\[
  i(f)=\bigl(f(u_\alpha),f(w_\alpha)\bigr)_{\alpha\in\Delta}.
\]
By \SGARefLink{sga3:XXIII:1.9}{Exposé XXIII, 1.9}, \(i\) is a monomorphism.

\par\medskip\noindent\SGAInlineRefTarget{sga3:xxiv:proposition:7-1-2}\textbf{Proposition 7.1.2.---}\enspace
\emph{If \(H\) is separated over \(S\), then \(\mathscr F\) is
representable and \(i\) is a closed immersion.}

The usual technique of relative
representability\footnote{\textit{Editors' note.} See, for example, the
proof of \SGARefLink{sga3:xxii:thm:5.8.1}{Exposé XXII, 5.8.1}.} shows that it is enough to prove the
following. Given sections
\[
  v_\alpha,h_\alpha\in H(S),
  \qquad \alpha\in\Delta,
\]
the \(S\)-schemes \(S'\) for which there exists an \(S'\)-homomorphism
\[
  f:G_{S'}\longrightarrow H_{S'}
\]
extending \((f_T)_{S'}\) and satisfying
\[
  f(u_\alpha)=v_\alpha,
  \qquad
  f(w_\alpha)=h_\alpha
\]
are exactly those which factor through a certain closed subscheme of
\(S\). We may plainly suppose \(S\) affine.

For convenience, say that a morphism \(Y\to
X\)\footnote{\textit{Editors' note.} We have corrected \(X\to Y\) to
\(Y\to X\).} of affine schemes satisfies condition \textup{(L)} if
\(Y\) is the spectrum of an \(\mathcal O(X)\)-algebra which is a free
\(\mathcal O(X)\)-module. Within the category of affine schemes,
products and composites of morphisms satisfying \textup{(L)} again
satisfy \textup{(L)}, and \textup{(L)} is stable under base change.

\par\medskip\noindent\SGAInlineRefTarget{sga3:xxiv:lemma:7-1-3}\textbf{Lemma 7.1.3.---}\enspace
\emph{Suppose that \(S\) is affine, and let \(\alpha\in\Delta\).
Consider the morphism
\[
  a:T\times_S T\longrightarrow G_{a,S},
  \qquad
  a(t,t')=\alpha(t)+\alpha(t').
\]
Then:
\begin{enumerate}
\item[(i)] \(a\) is faithfully flat and of finite presentation;
\item[(ii)] if \(R\) is the fiber square of \(a\), the structure
  morphism \(R\to S\) satisfies \textup{(L)}.
\end{enumerate}}

The morphism \(a\) is surjective, since
\(\alpha:T\to G_{m,S}\) is surjective, and it is plainly of finite
presentation. Since \(\alpha\) satisfies \textup{(L)}, it is enough,
for both assertions, to show that
\[
  u:G_{m,S}^2\longrightarrow G_{a,S},
  \qquad
  u(x,y)=x+y
\]
satisfies \textup{(L)}. Equivalently, for every ring \(A\), the ring
\[
  A[X,Y,X^{-1},Y^{-1}]
\]
must be free over its subring \(A[X+Y]\).\footnote{\textit{Editors'
note.} We have simplified the proof in the original.} Now \(A[X,Y]\)
is free with basis \(\{1,X\}\) over
\[
  B=A[X+Y,XY],
\]
because \(X^2-(X+Y)X+XY=0\). Hence
\[
  A[X^{\pm1},Y^{\pm1}]
  =A[X,Y]_{XY}
\]
is free over
\[
  B_{XY}=A[X+Y,(XY)^{\pm1}]
\]
with basis \(\{1,X\}\), and free over \(A[X+Y]\) with basis
\[
  (XY)^p,\quad X(XY)^p,
  \qquad p\in\mathbb Z.
\]

\par\medskip\noindent\SGAInlineRefTarget{sga3:xxiv:lemma:7-1-4}\textbf{Lemma 7.1.4.---}\enspace
\emph{Let \(\alpha\in\Delta\), and let
\[
  b:T\times_S T\longrightarrow H
\]
be the morphism defined on points by
\[
  b(t,t')
  =
  \operatorname{int}(f_T(t))(v_\alpha)\,
  \operatorname{int}(f_T(t'))(v_\alpha).
\]
Let \(f_\alpha:U_\alpha\to H\) be an \(S\)-morphism. The following
conditions are equivalent:\footnote{\textit{Editors' note.} In what
follows the group law of \(U_\alpha\) is written additively. If
\(p_\alpha:G_{a,S}\xrightarrow{\sim}U_\alpha\) is the isomorphism with
\(p_\alpha(1)=u_\alpha\), and \(x=p_\alpha(z)\), then
\(\alpha(t)x\) denotes \(p_\alpha(\alpha(t)z)\), while
\(a(t,t')u_\alpha\) denotes
\(p_\alpha(a(t,t'))=\alpha(t)u_\alpha+\alpha(t')u_\alpha\).}
\begin{enumerate}
\item[\SGAInlineRefTarget{sga3:xxiv:lemma:7-1-4:item:i}(i)] \(f_\alpha\) is a group homomorphism,
  \(f_\alpha(u_\alpha)=v_\alpha\), and
  \[
    f_T(t)f_\alpha(x)f_T(t)^{-1}
    =
    f_\alpha(\alpha(t)x)
  \]
  for all \(t\in T(S')\), \(x\in U_\alpha(S')\), and \(S'\to S\);
\item[\SGAInlineRefTarget{sga3:xxiv:lemma:7-1-4:item:ii}(ii)] \(f_\alpha(u_\alpha)=v_\alpha\), and
  \[
    f_\alpha(a(t,t')u_\alpha)=b(t,t')
  \]
  for all \(t,t'\in T(S')\) and \(S'\to S\).
\end{enumerate}}

If \(f_\alpha\) satisfies \SGARefLink{sga3:xxiv:lemma:7-1-4:item:i}{(i)}, then
\[
  f_\alpha(\alpha(t)u_\alpha)
  =
  \operatorname{int}(f_T(t))(v_\alpha),
\]
which immediately gives \SGARefLink{sga3:xxiv:lemma:7-1-4:item:ii}{(ii)}. Conversely, suppose \SGARefLink{sga3:xxiv:lemma:7-1-4:item:ii}{(ii)}. Since \(a\)
is an \textup{(fpqc)} covering, to prove the last condition in \SGARefLink{sga3:xxiv:lemma:7-1-4:item:i}{(i)} it
suffices to show that, for \(t,t',t''\in T(S')\),
\[
  f_T(t)\,
  f_\alpha(a(t',t'')u_\alpha)\,
  f_T(t)^{-1}
  =
  f_\alpha(\alpha(t)a(t',t'')u_\alpha).
\]
This is the same as
\[
  f_T(t)b(t',t'')f_T(t)^{-1}=b(tt',tt''),
\]
which follows directly from the definition of \(b\).

It remains to prove that \(f_\alpha\) is a group homomorphism. The
preceding property gives
\[
\begin{aligned}
 f_\alpha(\alpha(t)u_\alpha)
 f_\alpha(\alpha(t')u_\alpha)
 &=
 \bigl(f_T(t)f_\alpha(u_\alpha)f_T(t)^{-1}\bigr)
 \bigl(f_T(t')f_\alpha(u_\alpha)f_T(t')^{-1}\bigr)\\
 &=b(t,t')\\
  &=f_\alpha\bigl(\alpha(t)u_\alpha+
                   \alpha(t')u_\alpha\bigr).
\end{aligned}
\]
Thus \(f_\alpha(x+x')=f_\alpha(x)f_\alpha(x')\) whenever \(x,x'\) are
sections of the schematically dense open subset
\(U_\alpha^\times\subset U_\alpha\). \SGARefLink{sga3:xviii:prop:1.4}{Exposé XVIII, 1.4} then gives the
result.

\subsubsection*{7.1.5.}

Fix provisionally an \(\alpha\in\Delta\). Since \(a\) is faithfully
flat and quasi-compact, \(G_{a,S}\) is identified with the quotient of
\(T\times_S T\) by the equivalence relation \(R\) defined by \(a\).
Write \(i_1,i_2:R\to T\times_S T\) for its two projections.

\[
  % BEGIN inputs/XXIV/native_diagrams/exp24/006_equivalence_relation_pair.tex
\begin{tikzcd}[column sep=huge]
R
  \arrow[r,shift left=.7ex,"i_1"]
  \arrow[r,shift right=.7ex,"i_2"']
&
T\times_S T
\end{tikzcd}

% END inputs/XXIV/native_diagrams/exp24/006_equivalence_relation_pair.tex

\]

For \(b\) to factor through \(a\),\footnote{\textit{Editors' note.}
That is, for condition \SGARefLink{sga3:xxiv:lemma:7-1-4:item:ii}{7.1.4(ii)} to hold.} it is necessary and
sufficient that
\[
  b\circ i_1=b\circ i_2.
\]
Let \(K\) be the equalizer of \(b\circ i_1\) and \(b\circ i_2\).

\[
  % BEGIN inputs/XXIV/native_diagrams/exp24/007_kernel_pair.tex
\begin{tikzcd}[column sep=huge]
R
  \arrow[r,shift left=.7ex,"b\circ i_1"]
  \arrow[r,shift right=.7ex,"b\circ i_2"']
&
H
\end{tikzcd}

% END inputs/XXIV/native_diagrams/exp24/007_kernel_pair.tex

\]

Then the required condition is \(K=R\). Since \(H\) is separated over
\(S\), \(K\) is a closed subscheme of \(R\). On the other hand,
\(R\) is essentially free over \(S\), by \SGARefLink{sga3:xxiv:lemma:7-1-3}{7.1.3}.

\subsubsection*{7.1.6.}

It follows that, for an \(S\)-scheme \(S'\), an \(f_\alpha\) satisfying
\SGARefLink{sga3:xxiv:lemma:7-1-4:item:i}{7.1.4(i)} exists over \(S'\), and is then unique, if and only if
\[
  K_{S'}=R_{S'}
\]
and the resulting morphism satisfies
\[
  f_\alpha(u_\alpha)=v_\alpha.
\]
The first condition is equivalent to requiring \(S'\to S\) to factor
through a certain closed subscheme of \(S\) (\SGARefLink{sga3:viii:theorem:6:4}{Exposé VIII,
6.4}\footnote{\textit{Editors' note.} See also the addition in Exposé
VI\textsubscript{B}, \SGARefLink{thm:sga3-VIB-6.2.3}{6.2.3}.}). After replacing \(S\) by this closed
subscheme, the second condition again defines a closed subscheme of
\(S\), since it is the equality of two sections of the separated
\(S\)-scheme \(H\).

After replacing \(S\) by the resulting closed subscheme, we may thus
suppose that a morphism \(f_\alpha:U_\alpha\to H\) satisfying \SGARefLink{sga3:xxiv:lemma:7-1-4:item:i}{7.1.4(i)}
exists. Intersecting the closed subschemes obtained for all
\(\alpha\in\Delta\), we may suppose this condition holds for every
\(\alpha\in\Delta\).

\subsubsection*{7.1.7.}

Similarly, for each \(\alpha\in\Delta\), consider the two
\(S\)-group homomorphisms
\[
  f_T\circ\operatorname{int}(w_\alpha),
  \qquad
  \operatorname{int}(h_\alpha)\circ f_T
  :T\longrightarrow H.
\]
Since \(H\) is separated over \(S\) and \(T\) is essentially free over
\(S\), the same reasoning shows that, after replacing \(S\) by a
closed subscheme, one may suppose that
\[
  f_T\circ\operatorname{int}(w_\alpha)
  =
  \operatorname{int}(h_\alpha)\circ f_T
\]
for every \(\alpha\in\Delta\).

\subsubsection*{7.1.8.}

The generators-and-relations theorem (\SGARefLink{sga3:xxiii:subsection:3-5}{Exposé XXIII, 3.5}) now shows that
a group homomorphism \(f:G\to H\) with the required properties exists
if and only if a certain finite set of algebraic relations among the
sections
\[
  h_\alpha,\qquad
  v_\alpha,\qquad
  f_T\bigl(\alpha^*(-1)\bigr),
  \qquad \alpha\in\Delta,
\]
is satisfied:
\[
  R_i\!\left(
    (h_\alpha)_\alpha,
    (v_\alpha)_\alpha,
    \bigl(f_T(\alpha^*(-1))\bigr)_\alpha
  \right)=e,
  \qquad i=1,\ldots,n.
\]
Because \(H\) is separated over \(S\), these relations again define a
closed subscheme of \(S\), proving \SGARefLink{sga3:xxiv:proposition:7-1-2}{Proposition 7.1.2}.

\par\medskip\noindent\SGAInlineRefTarget{sga3:xxiv:corollary:7-1-9}\textbf{Corollary 7.1.9.---}\enspace
\emph{Let \(S\) be a scheme, \(G\) a split reductive \(S\)-group, \(T\)
its maximal torus, and \(H\) a separated \(S\)-group scheme. Let
\(\mathcal P\) be a property of morphisms satisfying:
\begin{enumerate}
\item[(i)] every closed immersion has \(\mathcal P\);
\item[(ii)] a composite of morphisms having \(\mathcal P\) has
  \(\mathcal P\);
\item[(iii)] \(\mathcal P\) is stable under base change;
\item[(iv)] the structure morphism \(H\to S\) has \(\mathcal P\).
\end{enumerate}
Then the canonical morphism
\[
  \underline{\operatorname{Hom}}_{S\text{-}\mathrm{gr.}}(G,H)
  \longrightarrow
  \underline{\operatorname{Hom}}_{S\text{-}\mathrm{gr.}}(T,H)
\]
is relatively representable by a separated morphism having
\(\mathcal P\).\footnote{\textit{Editors' note.} More explicitly, for
every representable
\(S'\to\underline{\operatorname{Hom}}_{S\text{-}\mathrm{gr.}}(T,H)\),
the fiber product
\(\underline{\operatorname{Hom}}_{S\text{-}\mathrm{gr.}}(G,H)
\times S'\) is representable, and its morphism to \(S'\) is separated
and has \(\mathcal P\).}}

Indeed, we may suppose \(G\) pinned. The structure morphism
\(H^{2n}\to S\) has \(\mathcal P\), and \SGARefLink{sga3:xxiv:proposition:7-1-2}{7.1.2} gives the assertion.

\par\medskip\noindent\SGAInlineRefTarget{sga3:xxiv:corollary:7-1-10}\textbf{Corollary 7.1.10.---}\enspace
\emph{Let \(S\) be a scheme, \(G\) a split reductive \(S\)-group, and
\(H\) a smooth \(S\)-group scheme with affine
fibers.\footnote{\textit{Editors' note.} We have suppressed the
superfluous hypothesis that \(H\) be quasi-separated.} Then
\(\underline{\operatorname{Hom}}_{S\text{-}\mathrm{gr.}}(G,H)\) is
representable by an \(S\)-scheme locally of finite presentation and
separated over \(S\).}

Since \(H\) is smooth over \(S\), its neutral component \(H^0\) is
defined; it has affine fibers and is smooth, separated, and of finite
presentation over \(S\) (Exposé VI\textsubscript{B}, \SGARefLink{thm:sga3-VIB-3.10}{3.10} and \SGARefLink{sga3:xxiv:proposition:5-5}{5.5}).
Since \(G\) has connected fibers, \(H\) may be replaced by \(H^0\).
As \(G\) and \(H\) are then of finite presentation, we may suppose
\(S\) noetherian and apply \SGARefLink{sga3:xxiv:corollary:7-1-9}{7.1.9} with \(\mathcal P\) equal to the
property ``of finite type''. By \SGARefLink{sga3:xv:theorem:8:8}{Exposé XV, 8.9},
\(\underline{\operatorname{Hom}}_{S\text{-}\mathrm{gr.}}(T,H)\) is
representable by a separated \(S\)-scheme locally of finite type.

\par\medskip\noindent\SGAInlineRefTarget{sga3:xxiv:remark:7-1-11}\textbf{Remark 7.1.11.---}\enspace
If \(H\) is affine over \(S\), the reference to Exposé XV may be
replaced by \SGARefLink{sga3:xi:corollary:4:2}{Exposé XI, 4.2}.

\subsection*{\SGAInlineRefTarget{sga3:xxiv:subsection:7-2}7.2. General case}

\par\medskip\noindent\SGAInlineRefTarget{sga3:xxiv:proposition:7-2-1}\textbf{Proposition 7.2.1.---}\enspace
\emph{Let \(S\) be a scheme, \(G\) a reductive \(S\)-group, and \(T\) a
maximal torus of \(G\). Let \(H\) be an \(S\)-group scheme whose
structure morphism \(H\to S\) satisfies:
\begin{quote}
\textup{(+)} Every point \(s\in S\) has an open neighborhood \(U\)
such that \(H_U\to U\) is quasi-projective.
\end{quote}
Then the canonical morphism
\[
  \underline{\operatorname{Hom}}_{S\text{-}\mathrm{gr.}}(G,H)
  \longrightarrow
  \underline{\operatorname{Hom}}_{S\text{-}\mathrm{gr.}}(T,H)
\]
is relatively representable by a morphism satisfying \textup{(+)}.}

When \(G\) is split relative to \(T\), apply 7.1.9 with
\(\mathcal P=\textup{(+)}\). When \(G\) is locally isotrivial, for
example when \(G\) is semisimple (4.2.2), note that the assertion is
local for the local finite étale topology. Indeed,
quasi-projectivity is local for the global finite étale topology and
ensures effectiveness of descent for that topology (SGA~1, VIII, 7.7).
The general case follows from the next lemma.

\par\medskip\noindent\SGAInlineRefTarget{sga3:xxiv:lemma:7-2-2}\textbf{Lemma 7.2.2.---}\enspace
\emph{Let \(S\) be a scheme, \(G\) a reductive \(S\)-group, \(T\) a
maximal torus of \(G\), \(G'\) the derived group of \(G\) (Exposé XXII,
6.2), and \(T'=T\cap G'\) the maximal torus of \(G'\) corresponding to
\(T\) (Exposé XXII, 6.2.8). The natural diagram involving
\(G,G',T,T'\) is a pushout in the category of \(S\)-group sheaves.
Equivalently, for every \(S\)-group sheaf \(H\), the corresponding
square of Hom-functors is cartesian.}

\[
  % BEGIN inputs/XXIV/native_diagrams/exp24/008_pushout_square.tex
\begin{tikzcd}[column sep=large,row sep=large]
G & G' \arrow[l] \\
T \arrow[u] & T' \arrow[l] \arrow[u]
\end{tikzcd}

% END inputs/XXIV/native_diagrams/exp24/008_pushout_square.tex

\]

\[
  % BEGIN inputs/XXIV/native_diagrams/exp24/009_hom_cartesian_square.tex
\begin{tikzcd}[column sep=large,row sep=large]
\underline{\operatorname{Hom}}_{S\text{-}\mathrm{gr.}}(G,H)
  \arrow[r]
  \arrow[d]
&
\underline{\operatorname{Hom}}_{S\text{-}\mathrm{gr.}}(G',H)
  \arrow[d]
\\
\underline{\operatorname{Hom}}_{S\text{-}\mathrm{gr.}}(T,H)
  \arrow[r]
&
\underline{\operatorname{Hom}}_{S\text{-}\mathrm{gr.}}(T',H)
\end{tikzcd}

% END inputs/XXIV/native_diagrams/exp24/009_hom_cartesian_square.tex

\]

If \(\operatorname{rad}(G)\) is the radical of \(G\), then
\(\operatorname{rad}(G)\subset T\), so
\[
  \operatorname{rad}(G)\cap G'
  =
  \operatorname{rad}(G)\cap T'
  =
  K.
\]
Multiplication in \(G\) induces isogenies (\SGARefLink{sga3:xxii:subsection:6-2}{Exposé XXII, 6.2})
\[
  i:G'\times_S\operatorname{rad}(G)\longrightarrow G,
  \qquad
  j:T'\times_S\operatorname{rad}(G)\longrightarrow T.
\]
Let
\[
  f_{G'}:G'\to H,\qquad
  f_T:T\to H,\qquad
  f_{T'}:T'\to H
\]
be \(S\)-group homomorphisms such that
\[
  f_{G'}|_{T'}=f_T|_{T'}=f_{T'}.
\]
We show that there is a unique \(S\)-group homomorphism \(f:G\to H\)
inducing \(f_{G'}\) and \(f_T\). Put
\[
  f_{\mathrm{rad}}
  =
  f_T|_{\operatorname{rad}(G)}
\]
\footnote{\textit{Editors' note.} We have replaced \(f_r\) by
\(f_{\mathrm{rad}}\).}
and define
\[
  f_1=f_{G'}\cdot f_{\mathrm{rad}}:
  G'\times_S\operatorname{rad}(G)\longrightarrow H.
\]
For \(f\) to exist, and it is plainly then unique, it is necessary and
sufficient that \(f_1\) be trivial on the kernel of \(i\), that is,
that \(f_{G'}\) and \(f_{\mathrm{rad}}\) coincide on \(K\). The
existence of \(f_T\) similarly shows that \(f_{T'}\) and
\(f_{\mathrm{rad}}\) coincide on \(K\). Finally, \(f_{G'}\) and
\(f_{T'}\) coincide on \(K\subset T'\), proving the lemma.

Arguing as in \SGARefLink{sga3:xxiv:corollary:7-1-10}{7.1.10}, \SGARefLink{sga3:xxiv:proposition:7-2-1}{Proposition 7.2.1} gives:

\par\medskip\noindent\SGAInlineRefTarget{sga3:xxiv:corollary:7-2-3}\textbf{Corollary 7.2.3.---}\enspace
\emph{Let \(S\) be a scheme, \(G\) a reductive \(S\)-group, and \(H\)
a smooth, quasi-projective \(S\)-group scheme with affine fibers. Then
\(\underline{\operatorname{Hom}}_{S\text{-}\mathrm{gr.}}(G,H)\) is
representable by an \(S\)-scheme locally of finite presentation and
separated over \(S\).}

\subsection*{\SGAInlineRefTarget{sga3:xxiv:subsection:7-3}7.3. Phenomena specific to characteristic zero}

If \(G\) and \(H\) are smooth \(S\)-groups, with Lie algebras
\(\mathfrak g\) and \(\mathfrak h\), there is a canonical morphism
\[
  \Lie:
  \underline{\operatorname{Hom}}_{S\text{-}\mathrm{gr.}}(G,H)
  \longrightarrow
  \underline{\operatorname{Hom}}_{\mathcal O_S\text{-}\mathrm{Lie}}
    (\mathfrak g,\mathfrak h),
\]
where the functor on the right has the evident meaning.

\par\medskip\noindent\SGAInlineRefTarget{sga3:xxiv:proposition:7-3-1}\textbf{Proposition 7.3.1.---}\enspace
\emph{Let \(S\) be a scheme of characteristic zero, that is, a
\(\mathbb Q\)-scheme; let \(G\) be a reductive \(S\)-group, and \(H\) a
smooth, quasi-projective \(S\)-group scheme with affine fibers.
\begin{enumerate}
\item[\SGAInlineRefTarget{sga3:xxiv:proposition:7-3-1:item:i}(i)]
\(\underline{\operatorname{Hom}}_{S\text{-}\mathrm{gr.}}(G,H)\) is
representable by a smooth separated \(S\)-scheme.
\item[(ii)] If \(G\) is semisimple, this scheme is affine and of finite
presentation over \(S\).
\item[\SGAInlineRefTarget{sga3:xxiv:proposition:7-3-1:item:iii}(iii)] If \(G\) is simply connected (\SGARefLink{sga3:xxii:def:4.3.3}{Exposé XXII, 4.3.3}), the
canonical morphism
\[
  \Lie:
  \underline{\operatorname{Hom}}_{S\text{-}\mathrm{gr.}}(G,H)
  \longrightarrow
  \underline{\operatorname{Hom}}_{\mathcal O_S\text{-}\mathrm{Lie}}
    (\mathfrak g,\mathfrak h)
\]
is an isomorphism.
\end{enumerate}}

\par\medskip\noindent\SGAInlineRefTarget{sga3:xxiv:lemma:7-3-2}\textbf{Lemma 7.3.2.---}\enspace
\emph{Let \(k\) be a field of characteristic zero, \(G\) a reductive
\(k\)-group, \(V\) a finite-dimensional \(k\)-vector space, and
\(G\to\operatorname{GL}(V)\) a linear representation. Then
\[
  H^0(G,V)=H^0(\mathfrak g,V),
  \qquad
  H^i(G,V)=0\quad(i>0).
\]}

The first equality is valid in general for a smooth connected
group.\footnote{\textit{Editors' note.} The inclusion
\(V^G\subset V^{\mathfrak g}\) is clear. Conversely,
\(C=\operatorname{Centr}_G(V^{\mathfrak g})\) is a closed subgroup of
\(G\), smooth because \(\operatorname{char}(k)=0\). By \SGARefLink{sga3:ii:theorem:5-3-1}{Exposé II,
5.3.1},
\(\Lie(C)=\operatorname{Centr}_{\mathfrak g}(V^{\mathfrak g})
=\mathfrak g\); because \(C\) is smooth and \(G\) connected, \(C=G\).
Thus \(V^{\mathfrak g}\subset V^G\). See also \([{\rm DG70}]\),
\S II.6, Proposition 2.1(c).} For a reductive group it can also be
proved as follows. We may suppose \(k\) algebraically closed, hence
\(G\) split and generated by subgroups isomorphic to
\(G_{m,k}\),\footnote{\textit{Editors' note.} The union of the maximal
tori of \(G\) is dense in \(G\); cf.\ \emph{Bible}, \S6.5,
Theorem~5, or \([{\rm Ch05}]\), \S6.6, Theorem~6.} and the assertion
for this group is immediate.

The first equality shows that the functor \(H^0(G,V)\) is exact in
\(V\) when \(G\) is semisimple: \(\mathfrak g\) is then a semisimple
Lie algebra, and one applies \([{\rm BLie}]\), \S I.6,
Exercise~1(b). The assertion remains true when \(G\) is reductive.
Indeed, introducing the radical \(R\) of
\(G\)\footnote{\textit{Editors' note.} Observe that \(R\) is a torus.}
and the semisimple quotient \(G/R\), one has
\[
  H^0(G,V)=H^0\bigl(G/R,H^0(R,V)\bigr),
\]
so \(H^0(G,-)\) is a composite of two exact functors. \SGARefLink{sga3:i:theorem:5-3-1}{Exposé I, 5.3.1}
then gives \(H^i(G,V)=0\) for \(i>0\).

\par\medskip\noindent\SGAInlineRefTarget{sga3:xxiv:remark:7-3-3}\textbf{Remark 7.3.3.---}\enspace
Under the preceding hypotheses, if \(G\) is semisimple, then
\[
  H^1(\mathfrak g,V)=H^2(\mathfrak g,V)=0;
\]
see \([{\rm BLie}]\), \emph{loc.\ cit.}, (b) and (d).

\subsubsection*{7.3.4.}

We prove the proposition. By \SGARefLink{sga3:xxiv:corollary:7-2-3}{7.2.3},
\(\underline{\operatorname{Hom}}_{S\text{-}\mathrm{gr.}}(G,H)\) is
already representable by an \(S\)-scheme locally of finite
presentation and separated over \(S\). To show it smooth, it suffices
to show it infinitesimally smooth (\SGARefLink{sga3:xi:remark:1:8}{Exposé XI,
1.8}),\footnote{\textit{Editors' note.} This means that for every local
artinian \(S\)-scheme \(S'\), with closed point \(s'\), every
homomorphism of \(\kappa(s')\)-groups
\(G_{s'}\to H_{s'}\) lifts to an \(S'\)-group homomorphism
\(G_{S'}\to H_{S'}\). By \SGARefLink{sga3:iii:corollary:2-8}{Exposé III, 2.8}, this follows from the
vanishing of \(H^2(G_{s'},V)\), where
\(V=\Lie(H_{s'}/\kappa(s'))\).} and this follows from \SGARefLink{sga3:iii:corollary:2-8:item:i}{Exposé III,
2.8(i)} and \SGARefLink{sga3:xxiv:lemma:7-3-2}{7.3.2}. This proves \SGARefLink{sga3:xxiv:proposition:7-3-1:item:i}{(i)}.

We next show that \SGARefLink{sga3:iii:corollary:2-8:item:ii}{(ii)} follows from \SGARefLink{sga3:xxiv:proposition:7-3-1:item:iii}{(iii)}. First it is enough to prove
that
\(\underline{\operatorname{Hom}}_{S\text{-}\mathrm{gr.}}(G,H)\) is
affine over \(S\), since it is then of finite presentation, being
smooth by \SGARefLink{sga3:xxiv:proposition:7-3-1:item:i}{(i)}. In any case,
\[
  \underline{\operatorname{Hom}}_{\mathcal O_S\text{-}\mathrm{Lie}}
    (\mathfrak g,\mathfrak h)
\]
is represented by a closed subscheme of
\[
  \underline{\operatorname{Hom}}_{\mathcal O_S\text{-}\mathrm{mod}}
    (\mathfrak g,\mathfrak h)
  \simeq
  \mathrm W\!\left(
    \mathfrak g^\vee\otimes_{\mathcal O_S}\mathfrak h
  \right),
\]
which is an affine \(S\)-scheme. Thus the desired conclusion follows
when \(G\) is simply connected.

In general, we may suppose \(G\) split, so that
\(G\simeq\operatorname{\acute Ep}_S(\mathscr R)\). Introduce the
simply connected root datum \(\operatorname{sc}(\mathscr R)\)
(\SGARefLink{sga3:XXI:6.5.5}{Exposé XXI, 6.5.5}). The existence theorem (\SGARefLink{sga3:xxv:theorem:1:1}{Exposé XXV, 1.1}) gives a
simply connected \(S\)-group \(\widetilde G\) and a central isogeny
\[
  \widetilde G\longrightarrow G.
\]
Its kernel \(K\) is a finite diagonalizable \(S\)-group, hence a
twisted constant \(S\)-group because \(S\) has characteristic zero.
The functor
\(\underline{\operatorname{Hom}}_{S\text{-}\mathrm{gr.}}(K,H)\) is
plainly represented by a separated \(S\)-scheme. For example, if
\(K\simeq(\mathbb Z/n\mathbb Z)_S\), then it is isomorphic to
\[
  \ker\!\left(H\xrightarrow{\,n\,}H\right).
\]
There is an exact sequence of pointed \(S\)-schemes
\[
  1\longrightarrow
  \underline{\operatorname{Hom}}_{S\text{-}\mathrm{gr.}}(G,H)
  \xrightarrow{u}
  \underline{\operatorname{Hom}}_{S\text{-}\mathrm{gr.}}
    (\widetilde G,H)
  \longrightarrow
  \underline{\operatorname{Hom}}_{S\text{-}\mathrm{gr.}}(K,H).
\]
Thus \(u\) is a closed immersion, and
\(\underline{\operatorname{Hom}}_{S\text{-}\mathrm{gr.}}(G,H)\) is
affine over \(S\).

\SGAInlineRefTarget{sga3:xxiv:subsubsection:7-3-5}\subsubsection*{7.3.5.}

It remains to prove \SGARefLink{sga3:xxiv:proposition:7-3-1:item:iii}{(iii)}. As in the proof of \SGARefLink{sga3:xxiv:proposition:7-3-1:item:i}{(i)}, and using \SGARefLink{sga3:xxiv:remark:7-3-3}{7.3.3},
one shows that
\[
  \underline{\operatorname{Hom}}_{\mathcal O_S\text{-}\mathrm{Lie}}
    (\mathfrak g,\mathfrak h)
\]
is smooth over \(S\). We may therefore suppose
\(S=\operatorname{Spec}(k)\), where \(k\) is algebraically closed of
characteristic zero. It is enough to show that \(\Lie\) is bijective
on \(k\)-points and induces a bijection on tangent spaces at
corresponding points. First we show that
\[
  \operatorname{Hom}_{k\text{-}\mathrm{gr.}}(G,H)
  \longrightarrow
  \operatorname{Hom}_{k\text{-}\mathrm{Lie}}
    (\mathfrak g,\mathfrak h)
\]
is bijective.

If \(u:G\to H\) is a \(k\)-group homomorphism, its graph is a connected
subgroup of \(G\times_k H\), determined by its Lie algebra, which is
the graph of \(\Lie(u)\). Hence the map is injective. Conversely, if
\(v:\mathfrak g\to\mathfrak h\) is a Lie algebra homomorphism, its
graph \(\mathfrak k\) is a Lie subalgebra of
\(\mathfrak g\times\mathfrak h\), isomorphic to \(\mathfrak g\).
Because \(\mathfrak g\) is its own derived algebra, so is
\(\mathfrak k\). A theorem of Chevalley
\([{\rm Ch51}]\), \S14, Theorem~15, therefore shows that
\(\mathfrak k\) is the Lie algebra of a connected subgroup \(K\) of
\(G\times_k H\).\footnote{\textit{Editors' note.} See also
\([{\rm Bo91}]\), II 7.9. We have also added the following sentence.}

Since \(\mathfrak k\simeq\mathfrak g\) is semisimple, \(K\) is a
semisimple \(k\)-group (the French prints \(G\) here; the argument
and the following sentence require \(K\)).
Moreover,
\[
  \dim K=\dim\mathfrak k=\dim\mathfrak g=\dim G,
\]
and \(K\cap H\) is finite because its Lie algebra is zero. Thus the
canonical morphism
\[
  \operatorname{pr}_1:K\longrightarrow G
\]
is finite and dominant; because \(G\) is connected, it is surjective,
and hence an isogeny. Since \(G\) is simply connected and \(K\) is
semisimple, \SGARefLink{sga3:XXI:6.2.7}{Exposé XXI, 6.2.7} shows that \(\operatorname{pr}_1\) is an
isomorphism. Consequently \(K\) is the graph of a \(k\)-group
homomorphism \(u:G\to H\) with \(\Lie(u)=v\).

Finally, let \(u:G\to H\) be a \(k\)-group homomorphism. The tangent
space at \(u\) to
\[
  \underline{\operatorname{Hom}}_{k\text{-}\mathrm{gr.}}(G,H)
\]
identifies with \(Z^1(G,\mathfrak h)\) (\SGARefLink{sga3:ii:subsection:4-2}{Exposé II, 4.2}), where
\(G\) acts on \(\mathfrak h\) through
\(\operatorname{Ad}_H\circ u\). Similarly, the tangent space to
\[
  \underline{\operatorname{Hom}}_{k\text{-}\mathrm{Lie}}
    (\mathfrak g,\mathfrak h)
\]
at \(\Lie(u)\) identifies with \(Z^1(\mathfrak g,\mathfrak h)\). We
must show that
\[
  Z^1(G,\mathfrak h)\longrightarrow
  Z^1(\mathfrak g,\mathfrak h)
\]
is bijective.

\[
  % BEGIN inputs/XXIV/native_diagrams/exp24/010_group_lie_cohomology_rows.tex
\begin{tikzcd}[column sep=small,row sep=large]
0 \arrow[r]
&
\mathfrak h^G
  \arrow[r]
  \arrow[d]
&
\mathfrak h
  \arrow[r]
  \arrow[d,equal]
&
Z^1(G,\mathfrak h)
  \arrow[r]
  \arrow[d]
&
H^1(G,\mathfrak h)
  \arrow[d]
\\
0 \arrow[r]
&
\mathfrak h^{\mathfrak g}
  \arrow[r]
&
\mathfrak h
  \arrow[r]
&
Z^1(\mathfrak g,\mathfrak h)
  \arrow[r]
&
H^1(\mathfrak g,\mathfrak h)
\end{tikzcd}

% END inputs/XXIV/native_diagrams/exp24/010_group_lie_cohomology_rows.tex

\]

By \SGARefLink{sga3:xxiv:lemma:7-3-2}{7.3.2} and \SGARefLink{sga3:xxiv:remark:7-3-3}{7.3.3},
\[
  \mathfrak h^G=\mathfrak h^{\mathfrak g},
  \qquad
  H^1(G,\mathfrak h)=0=H^1(\mathfrak g,\mathfrak h),
\]
which gives the desired conclusion.

\par\medskip\noindent\SGAInlineRefTarget{sga3:xxiv:remark:7-3-6}\textbf{Remark 7.3.6.---}\enspace
It is presumably true that, if \(S\) is locally noetherian, \(G\) a
simply connected semisimple \(S\)-group, \(H\) a smooth \(S\)-group
scheme, and \(\widehat G,\widehat H\) their formal completions along
the identity section, then every homomorphism
\(\widehat G\to\widehat H\) comes from a unique homomorphism \(G\to H\).
This would generalize \SGARefLink{sga3:xxiv:proposition:7-3-1:item:iii}{7.3.1(iii)}. When \(S=\operatorname{Spec}(k)\)
and \(H\) is affine and of finite type, this follows from a theorem of
Dieudonné \([{\rm Di57}]\), \S15, Theorem~4.\footnote{\textit{Editors'
note.} See Exposé VII\textsubscript{B} for the definition of the formal
groups \(\widehat G\) and \(\widehat H\). Suppose that
\(S=\operatorname{Spec}(k)\), where \(k\) is a field. By loc. cit.,
2.2.1, giving a morphism of formal \(k\)-groups
\(v:\widehat G\to\widehat H\) is equivalent to giving a morphism of
Hopf \(k\)-algebras
\(\phi:\operatorname{Dist}(G)\to\operatorname{Dist}(H)\), where
\(\operatorname{Dist}(G)\) denotes the algebra of distributions (at
the origin) of \(G\); see Exposé VII\textsubscript{A}, \SGARefLink{sga3:VIIA:2.1}{2.1},
\([{\rm DG70}]\), \S II.4, \SGARefLink{sga3:xxiv:lemma:6-1}{6.1}, or \([{\rm Ja87}]\), I 7.7. It is called
the hyperalgebra of \(G\) in \([{\rm Di57}]\) and \([{\rm Ta75}]\).
Theorem~4 of \([{\rm Di57}]\) (see also \([{\rm Ta75}]\), 0.3.4 (f)
and (g)) generalizes in this context the theorem of Chevalley used in
\SGARefLink{sga3:xxiv:subsubsection:7-3-5}{7.3.5}: one obtains a closed connected \(k\)-subgroup
\(K\subset G\times H\) such that \(\widehat K\) equals the graph of
\(v\). Since \(\operatorname{Dist}(K)\to\operatorname{Dist}(G)\) is an
isomorphism, \(K\to G\) is finite étale. It follows that \(K\) is
semisimple and then, by \SGARefLink{sga3:XXI:6.2.7}{Exposé XXI, 6.2.7}, that \(K\to G\) is an
isomorphism (because \(G\) is simply connected; see also
\([{\rm Ta75}]\), \SGARefLink{sga3:xxiv:1.8}{1.8} and \SGARefLink{sga3:xxiv:2.2}{2.2}). More generally, loc. cit. studies the
\(k\)-groups \(G\) having property \textup{(SC)}: every finite étale
morphism of \(k\)-groups \(G'\to G\), with \(G'\) connected, is an
isomorphism. Finally, the preceding argument shows that every
finite-dimensional \(\operatorname{Dist}(G)\)-module is, uniquely, a
\(G\)-module; for an extension to a split reductive \(k\)-group \(G\)
(or even to a Borel subgroup of \(G\)), see \([{\rm Ja87}]\), II 1.20
(and the references given there).}

\subsection*{\SGAInlineRefTarget{sga3:xxiv:subsection:7-4}7.4. An example}

As an example, we shall determine
\[
  \underline{\operatorname{Hom}}_{S\text{-}\mathrm{gr.}}
    (\mathrm{SL}_{2,S},\mathrm{SL}_{2,S}).
\]

\subsubsection*{7.4.1.}

Recall (Exposé XX, no.~5) that \(\mathrm{SL}_{2,S}\) is the
\(S\)-group scheme of matrices
\[
  \begin{pmatrix}a&b\\c&d\end{pmatrix}
\]
over \(S\) satisfying \(ad-bc=1\). A maximal torus \(T\) is defined
by the monomorphism
\[
  \alpha^*:\mathbb G_{m,S}\longrightarrow\mathrm{SL}_{2,S},
  \qquad
  \alpha^*(z)=
  \begin{pmatrix}z&0\\0&z^{-1}\end{pmatrix}.
\]
A root of \(G\) with respect to \(T\) is defined by
\[
  \alpha\bigl(\alpha^*(z)\bigr)=z^2,
\]
and a corresponding monomorphism is
\[
  p:\mathbb G_{a,S}\longrightarrow\mathrm{SL}_{2,S},
  \qquad
  p(x)=
  \begin{pmatrix}1&x\\0&1\end{pmatrix}.
\]
Finally, putting \(u=p(1)\), the corresponding representative of the
Weyl group is
\[
  w=
  \begin{pmatrix}0&1\\-1&0\end{pmatrix}.
\]

Recall also (\SGARefLink{sga3:xx:6.2}{Exposé XX, 6.2}) that \(\mathrm{SL}_{2,S}\) is generated by
\(T\), \(p(\mathbb G_{a,S})\), and \(w\), subject to the relations
\[
  \alpha^*(z)p(x)\alpha^*(z^{-1})=p(z^2x),
\]
\[
  w\alpha^*(z)w^{-1}=\alpha^*(z^{-1}),
  \qquad
  w^2=\alpha^*(-1),
  \qquad
  (wu)^3=1.
\]

\subsubsection*{7.4.2.}

Let \(f\) be an endomorphism of \(\mathrm{SL}_{2,S}\). It first
defines a homomorphism
\[
  f\circ\alpha^*:\mathbb G_{m,S}\longrightarrow\mathrm{SL}_{2,S}.
\]
The closed subgroup
\(\ker(f\circ\alpha^*)\subset\mathbb G_{m,S}\) is, locally on \(S\),
either \(\mu_{n,S}\), with \(n\geq1\), or
\(\mathbb G_{m,S}\). After restricting \(S\), we may therefore assume
that there are an integer \(n\in\mathbb N\) and a monomorphism
\[
  f':\mathbb G_{m,S}\longrightarrow\mathrm{SL}_{2,S}
\]
such that
\[
  f\circ\alpha^*(z)=f'(z^n).
\]
By conjugacy of monomorphisms
\(\mathbb G_{m,S}\to\mathrm{SL}_{2,S}\), after an étale covering
extension of the base there is a section \(g\) of \(G\) such that
\[
  f\circ\alpha^*(z)
  =
  \operatorname{int}(g)\circ\alpha^*(z^n).
\]
Replacing \(f\) by an inner conjugate, we are thus reduced to the case
\[
  f\circ\alpha^*(z)=\alpha^*(z^n)
\]
for some \(n\in\mathbb N\).

\subsubsection*{7.4.3.}

Consider now
\[
  f\circ p:\mathbb G_{a,S}\longrightarrow\mathrm{SL}_{2,S}.
\]
It satisfies
\[
  \alpha^*(z^n)(f\circ p)(x)\alpha^*(z^n)^{-1}
  =
  (f\circ p)(z^2x).
\]
If \(n=0\), it follows at once that \(f\circ p\) is invariant under
the homotheties of \(\mathbb G_{a,S}\), and hence is constant. If
\(n\neq0\), \SGARefLink{sga3:xxii:prop:4.1.9}{Exposé XXII, 4.1.9} applies: the map
\(x\mapsto x^n\) is an endomorphism of \(\mathbb G_{a,S}\), and there
is a \(\lambda\in\mathbb G_{m,S}\) such that
\[
  f\circ p(x)=p(\lambda x^n).
\]
This relation is also valid for \(n=0\), on taking \(\lambda=1\).
After extending \(S\) once more, choose
\(z\in\mathbb G_{m,S}\) with \(z^2=\lambda\). Replacing \(f\) by
\(\operatorname{int}(\alpha^*(z))\circ f\), we are reduced to
\[
  f\circ\alpha^*(z)=\alpha^*(z^n),
  \qquad
  f\circ p(x)=p(x^n).
\]

\subsubsection*{7.4.4.}

Finally, one must have
\[
  f(w)\alpha^*(z^n)f(w)^{-1}=\alpha^*(z^n)^{-1},
  \qquad
  (f(w)u)^3=1.
\]
By \SGARefLink{sga3:xx:3.8}{Exposé XX, 3.8}, this implies \(f(w)=w^n\). Since \(G\) is generated
by \(T\), \(p(\mathbb G_{a,S})\), and \(w\), it follows that, for every
\(S'\to S\) and every
\(\begin{pmatrix}a&b\\c&d\end{pmatrix}\in G(S')\),
\[
  f\!\begin{pmatrix}a&b\\c&d\end{pmatrix}
  =
  \begin{pmatrix}a^n&b^n\\c^n&d^n\end{pmatrix}.
\]

\subsubsection*{7.4.5.}

Summarizing, locally on \(S\) for the étale topology, every
endomorphism \(f\) of \(\mathrm{SL}_{2,S}\) can be written
\[
  f=\phi\circ F_n,
\]
where \(\phi\) is an inner automorphism, \(n\geq0\), and
\[
  F_n\!\begin{pmatrix}a&b\\c&d\end{pmatrix}
  =
  \begin{pmatrix}a^n&b^n\\c^n&d^n\end{pmatrix}.
\]
If \(f=\phi\circ F_n\), the integer \(n\) is determined by any fiber
\(f_s\), for example by \(\ker(f_s)\). Consequently \(n\) is a
locally constant function on \(S\).

\subsubsection*{7.4.6.}

It follows that, for every endomorphism \(f\) of
\(\mathrm{SL}_{2,S}\), the scheme \(S\) has a canonical decomposition
into open and closed subschemes
\[
  S_0,\qquad S_1,\qquad S_{p^n},
\]
where \(p^n\) runs through the positive powers of prime numbers, such
that:
\begin{enumerate}
\item[(i)] \(f_{S_0}\) is the zero morphism;
\item[(ii)] \(f_{S_1}\) is an isomorphism, necessarily an inner
  automorphism;
\item[(iii)] \(S_{p^n}\) has characteristic \(p\), and
  \(f_{S_{p^n}}\) has a unique decomposition
  \[
    f_{S_{p^n}}=\phi\circ F_p^n,
  \]
  where \(\phi\) is an inner automorphism and \(F_p\) is the Frobenius
  endomorphism of \(\mathrm{SL}_{2,\mathbb F_p}\).
\end{enumerate}

\subsubsection*{7.4.7.}

Equivalently,
\[
  \underline{\operatorname{Hom}}_{\mathbb Z\text{-}\mathrm{gr.}}
    (\mathrm{SL}_{2,\mathbb Z},\mathrm{SL}_{2,\mathbb Z})
\]
is the disjoint-union scheme consisting of:
\begin{enumerate}
\item[(i)] one scheme isomorphic to \(\operatorname{Spec}(\mathbb Z)\);
\item[(ii)] one scheme isomorphic to
  \[
    \underline{\operatorname{Aut}}_{\mathbb Z\text{-}\mathrm{gr.}}
      (\mathrm{SL}_{2,\mathbb Z})
    \simeq
    \operatorname{ad}(\mathrm{SL}_{2,\mathbb Z});
  \]
\item[(iii)] for each prime \(p\) and each integer \(n>0\), one scheme
  isomorphic to
  \[
    \underline{\operatorname{Aut}}_{\mathbb F_p\text{-}\mathrm{gr.}}
      (\mathrm{SL}_{2,\mathbb F_p})
    \simeq
    \operatorname{ad}(\mathrm{SL}_{2,\mathbb F_p}).
  \]
\end{enumerate}

\subsubsection*{7.4.8.}

In particular:
\begin{enumerate}
\item[(i)]
\(\underline{\operatorname{Hom}}_{\mathbb F_p\text{-}\mathrm{gr.}}
(\mathrm{SL}_{2,\mathbb F_p},\mathrm{SL}_{2,\mathbb F_p})\)
has infinitely many connected components and is therefore not
quasi-compact.
\item[(ii)] If \(S\) is a scheme of unequal characteristics, then
\(\underline{\operatorname{Hom}}_{S\text{-}\mathrm{gr.}}
(\mathrm{SL}_{2,S},\mathrm{SL}_{2,S})\) is not flat over \(S\).
\end{enumerate}

\section*{\SGAInlineRefTarget{sga3:xxiv:section:8}8. Appendix: Cohomology of a smooth group\\
over a henselian ring. Cohomology and the functor \(\prod\)}

\par\medskip\noindent\SGAInlineRefTarget{sga3:xxiv:proposition:8-1}\textbf{Proposition 8.1.---}\enspace
\emph{Let \(S\) be a henselian local scheme, \(s\) its closed point,
and \(G\) a smooth \(S\)-group scheme such that every finite subset of
\(G\) is contained in an affine open subset.%
\begingroup
\renewcommand{\thefootnote}{*}%
\footnote{This last condition is in fact unnecessary; see
A.~Grothendieck, \textit{Le groupe de Brauer III}, in
\textit{Dix Exposés sur la cohomologie des schémas}, North-Holland,
1968, Theorem~11.7 and Remarks~11.8.3.}%
\endgroup
\addtocounter{footnote}{-1}
\begin{enumerate}
\item[(i)] If \(P\) is a principal homogeneous bundle under \(G\),
  there is a finite surjective étale morphism \(S'\to S\) which
  trivializes \(P\). Moreover,
  \[
    \operatorname{Fib}(S,G)\simeq H^1_{\acute et}(S,G).
  \]
\item[\SGAInlineRefTarget{sga3:xxiv:proposition:8-1:item:ii}(ii)] For every finite surjective étale morphism \(S'\to S\),
  the canonical map
  \[
    H^1(S'/S,G)
    \longrightarrow
    H^1(S'\otimes_S\kappa(s)/\kappa(s),G_s)
  \]
  is bijective.
\item[(iii)] The canonical map
  \[
    \operatorname{Fib}(S,G)
    \longrightarrow
    \operatorname{Fib}(\kappa(s),G_s)
  \]
  is bijective.
\end{enumerate}}

\subsubsection*{8.1.2.}

If \(K\) is a finite separable extension of \(\kappa(s)\), there is a
finite surjective étale morphism \(S'\to S\) such that
\[
  K\simeq S'\otimes_S\kappa(s).
\]
\footnote{\textit{Editors' note.} The scheme \(S'\) is still local and
henselian. We have also retained the numbering of the original: there
is no no.~8.1.1.}
If \(P\) is a principal homogeneous bundle under \(G\), then \(P\) is
smooth over \(S\), hence \(P_s\) is smooth over \(\kappa(s)\).
Consequently there is a finite separable extension \(K/\kappa(s)\)
such that \(P_K\) has a section (EGA IV\textsubscript{4}, 17.15.10).
Representing \(K\) as above, Hensel's lemma
(EGA IV\textsubscript{4}, 18.5.17) shows that \(P_{S'}\) has a section.
This proves the first assertion of (i).

Conversely, if \(P\) is a principal homogeneous sheaf under \(G\) for
the étale topology, then a finite surjective étale morphism
\(S'\to S\) trivializes \(P\): every étale covering family of a
henselian local scheme is dominated by a family consisting of one
morphism of that form. By the descent hypothesis imposed on \(G\),
\(P\) is representable (SGA~1, VIII, 7.6). This completes the proof of
(i), and shows that (ii) implies (iii). It remains to prove (ii).

\subsubsection*{8.1.3.}

The map in (ii) is injective. Let \(P\) and \(Q\) be two principal
homogeneous bundles under \(G\), trivialized by \(S'\), and consider
the \(S\)-group sheaf
\[
  \mathcal H=
  \underline{\operatorname{Isom}}_{G\text{-}\mathrm{bundles}}(P,Q).
\]
Since \(\mathcal H_{S'}\) is isomorphic to \(G_{S'}\),
\(\mathcal H\) is representable by the second hypothesis on \(G\), as
above. If \(\mathcal H(\kappa(s))\neq\varnothing\), Hensel's lemma
gives \(\mathcal H(S)\neq\varnothing\), and hence \(P\simeq Q\).

\subsubsection*{8.1.4.}

To prove surjectivity in \SGARefLink{sga3:xxiv:proposition:8-1:item:ii}{(ii)}, it is equivalent to prove that
\[
  Z^1(S'/S,G)
  \longrightarrow
  Z^1(S'\otimes_S\kappa(s)/\kappa(s),G_s)
\]
is surjective. Let \(Z^1\) be the \(S\)-functor defined by
\[
  Z^1(T)=Z^1(S'\times_ST/T,G_T).
\]
The preceding map is then
\[
  Z^1(S)\longrightarrow Z^1(\kappa(s)).
\]
Another application of Hensel's lemma reduces the assertion to
showing that \(Z^1\) is representable by a smooth \(S\)-scheme.

\subsubsection*{8.1.5.}

We first show that \(Z^1\) is representable by an \(S\)-scheme locally
of finite presentation. For \(i=0,1,\ldots\), define the \(S\)-functor
\(C_i\) by
\[
  C_i(T)=C_i(S'\times_ST/T,G_T).
\]
Equivalently,
\[
  C_i(T)
  =
  G\bigl((S'/S)^{i+1}\times_ST\bigr),
  \qquad
  C_i
  =
  \underline{\operatorname{Hom}}_S((S'/S)^{i+1},G).
\]
Since \(Z^1\) is obtained from \(C_1\) and \(C_2\) by fiber products,
it is enough to prove that \(C_i\), for \(i=1,2\), is representable by
an \(S\)-scheme locally of finite presentation.

\subsubsection*{8.1.6.}

Let \(T\to S\) be a finite surjective étale morphism which splits
\(S'\). Then
\[
  C_i\times_ST
  =
  \underline{\operatorname{Hom}}_T
    ((S'\times_ST/T)^{i+1},G_T)
\]
is represented by a product of \(n\) copies of \(G_T\), where \(n\)
is the degree of \((S'/S)^{i+1}\). The hypothesis on \(G\), applied
once again, shows that \(C_i\) is represented by an \(S\)-scheme
locally of finite presentation (SGA~1, VIII, \textit{loc. cit.}).

\subsubsection*{8.1.7.}

It remains to prove that \(Z^1\) is smooth. By definition, if \(T\) is
an affine \(S\)-scheme and \(T_0\) is the closed subscheme defined by
a square-zero ideal \(J\), we must show that
\[
  Z^1(T)\longrightarrow Z^1(T_0)
\]
is surjective. Since \(G\) is smooth, the canonical map
\(G(T)\to G(T_0)\) is surjective, and it is enough to prove that
\[
  H^1(S'\times_ST/T,G_T)
  \longrightarrow
  H^1(S'\times_ST_0/T_0,G_{T_0})
\]
is bijective. After a slight change of notation and a harmless
generalization of the hypotheses, this reduces to the following
lemma.

\par\medskip\noindent\SGAInlineRefTarget{sga3:xxiv:lemma:8-1-8}\textbf{Lemma 8.1.8.---}\enspace
\emph{Let \(S\) and \(S'\) be affine schemes, let \(S'\to S\) be
faithfully flat, let \(J\) be a square-zero ideal on \(S\), and let
\(S_0\) be the closed subscheme which it defines. Let \(G\) be a
smooth \(S\)-group. Put
\[
  S'_0=S'\times_SS_0,
  \qquad
  G_0=G\times_SS_0.
\]
Then the canonical map
\[
  H^1(S'/S,G)\longrightarrow H^1(S'_0/S_0,G_0)
\]
is bijective.}

\par\medskip\noindent\SGAInlineRefTarget{sga3:xxiv:remark:8-1-9}\textbf{Remark 8.1.9.---}\enspace
If \(G\) is commutative, the same assertion holds for every
\(H^i\), \(i>0\), with an analogous proof.

\smallskip\noindent\emph{Proof.}
Let \(M_0\) be the quasi-coherent \(\mathcal O_{S_0}\)-module
\[
  M_0
  =
  \Lie(G_0/S_0)\otimes_{\mathcal O_{S_0}}J.
\]
For every \(S_0\)-scheme \(T_0\), put
\[
  M_0(T_0)
  =
  H^0\!\left(
    T_0,
    M_0\otimes_{\mathcal O_{S_0}}\mathcal O_{T_0}
  \right),
\]
and define the \(S\)-group functors
\[
  M=\prod_{S_0/S}M_0,
  \qquad
  \overline G=\prod_{S_0/S}G_0
\]
by
\[
  M(T)=M_0(T_0),
  \qquad
  \overline G(T)=G_0(T_0),
  \qquad
  T_0=T\times_SS_0.
\]
By \SGARefLink{sga3:iii:corollary:0-9}{Exposé III, 0.9} and (\SGARefLink{sga3:iii:remark:0-6:equation:0-6-2}{0.6.2}), for every affine \(S\)-scheme \(T\)
there is a functorial exact sequence
\[
  1\longrightarrow M(T)\longrightarrow G(T)
  \longrightarrow\overline G(T)\longrightarrow1.
\]
We must study the map
\[
  H^1(S'/S,G)
  \longrightarrow
  H^1(S'_0/S_0,G_0)
  =
  H^1(S'/S,\overline G).
\]

Suppose first that \(G\) is commutative. There is then an exact
cohomology sequence
\[
  H^1(S'/S,M)
  \longrightarrow H^1(S'/S,G)
  \longrightarrow H^1(S'/S,\overline G)
  \longrightarrow H^2(S'/S,M).
\]
Moreover,
\[
  H^i(S'/S,M)
  =
  H^i(S'_0/S_0,M_0)
  =
  H^i(S'_0/S_0,\mathcal M_0)
  =
  0
  \qquad(i\neq0),
\]
the last equality following from TDTE I, B, Lemma~1.1.

Now suppose that \(G\) is not commutative. We use the exact sequence
of nonabelian cohomology. If
\(u\in Z^1(S'/S,\overline G)\), the elements of \(H^1(S'/S,G)\) whose
image in \(H^1(S'/S,\overline G)\) equals the class of \(u\) lie in the
image of the corresponding coboundary map
\[
  H^1(S'/S,M_u)\longrightarrow H^1(S'/S,G),
\]
where \(M_u\) is \(M\) twisted by \(u\). Similarly, for
\(v\in Z^1(S'/S,\overline G)\), there is an obstruction
\[
  \Delta(v)\in H^2(S'/S,M_v)
\]
which vanishes exactly when the class of \(v\) is in the image of
\(H^1(S'/S,G)\). It therefore suffices to prove
\[
  H^1(S'/S,M_u)=0,
  \qquad
  H^2(S'/S,M_v)=0.
\]

Recall (\SGARefLink{sga3:iii:proposition:0-8}{Exposé III, 0.8}) that the action of \(G\) on \(M\) defined by
the exact sequence above is induced functorially by the adjoint
representation
\[
  \operatorname{ad}:G_0
  \longrightarrow
  \underline{\operatorname{Aut}}_{\mathcal O_{S_0}}
    (\Lie(G_0/S_0)).
\]
The cocycle \(u\), respectively \(v\), therefore defines a descent
datum on \(\Lie(G_0/S_0)\); write \(L_u\), respectively \(L_v\), for
the resulting quasi-coherent \(\mathcal O_{S_0}\)-module. For
\(T\to S\),
\[
  M_u(T)
  =
  H^0\!\left(
    T_0,
    L_u\otimes_{\mathcal O_{S_0}}J
      \otimes_{\mathcal O_{S_0}}\mathcal O_{T_0}
  \right),
\]
and similarly with \(v\) in place of \(u\). Hence
\[
  H^1(S'/S,M_u)=H^1(S'_0/S_0,L_u\otimes J),
\]
\[
  H^2(S'/S,M_v)=H^2(S'_0/S_0,L_v\otimes J),
\]
and both groups vanish by the result already used.
\hfill\(\square\)

\par\medskip\noindent\SGAInlineRefTarget{sga3:xxiv:proposition:8-2}\textbf{Proposition 8.2.---}\enspace
\emph{Let \(\mathcal C\) be a category with fiber products, endowed
with a topology coarser than the canonical topology. Let \(S'\to S\)
be a morphism of \(\mathcal C\), let \(G'\) be an \(S'\)-group sheaf,
and let
\[
  G=\prod_{S'/S}G'
\]
be the corresponding \(S\)-group sheaf. Let
\[
  H^1_S(S',G')\subset H^1(S',G')
\]
be the subset of classes of principal homogeneous sheaves under
\(G'\) which are trivialized by a sieve on \(S'\) obtained by base
change from a suitable covering sieve on \(S\). The canonical map
\[
  H^1(S,G)\longrightarrow H^1(S',G'),
  \qquad
  P\longmapsto P\times_SS',
\]
induces a bijection
\[
  H^1(S,G)\xrightarrow{\sim}H^1_S(S',G').
\]
Its inverse is induced by
\[
  P'\longmapsto\prod_{S'/S}P'.
\]}

For every object \(X\) of \(\mathcal C/S\), the definition gives an
isomorphism, functorial in \(X\),
\[
  \operatorname{Hom}_S(X,G)
  \xrightarrow{\sim}
  \operatorname{Hom}_{S'}(X\times_SS',G').
\]
Thus, for every \(S\)-object \(T\), there is a bijection functorial in
\(T\),
\[
  H^1(T/S,G)
  \xrightarrow{\sim}
  H^1(T'/S',G'),
  \qquad
  T'=T\times_SS'.
\]
Replacing the single morphism \(T\to S\) by an arbitrary covering
family of \(S\), and then passing to the inductive limit, gives the
first part of the proposition. The second follows immediately.

\par\medskip\noindent\SGAInlineRefTarget{sga3:xxiv:lemma:8-3}\textbf{Lemma 8.3.---}\enspace
\emph{Under the hypotheses of \SGARefLink{sga3:xxiv:proposition:8-2}{8.2}, the equality
\[
  H^1_S(S',G')=H^1(S',G')
\]
is local on \(S\). More precisely, suppose that there is a covering
family \(\{S_i\to S\}\) such that, for every \(i\),
\[
  H^1_{S_i}(S'\times_SS_i,G')
  =
  H^1(S'\times_SS_i,G').
\]
Then \(H^1_S(S',G')=H^1(S',G')\).}

Indeed, let \(P'\) be a principal homogeneous sheaf under \(G'\), and
put
\[
  P'_i=P'\times_{S'}(S'\times_SS_i).
\]
By hypothesis, there is a covering family
\(\{S_{ij}\to S_i\}\) such that, for each \(j\),
\[
  P'\times_{S'}(S'\times_SS_{ij})
\]
has a section. But \(\{S_{ij}\to S\}\) is a covering family of \(S\),
and \(P'\) is therefore trivialized by the covering family of \(S'\)
obtained from it by base change.

\par\medskip\noindent\SGAInlineRefTarget{sga3:xxiv:proposition:8-4}\textbf{Proposition 8.4.---}\enspace
\emph{Let \(S'\to S\) be a finite étale morphism of schemes. Let
\(G'\) be an \(S'\)-group sheaf and let
\[
  G=\prod_{S'/S}G'
\]
be the corresponding \(S\)-group sheaf. For the étale topology,
respectively the local finite étale topology, respectively the
\textup{(fpqc)} topology, the following functors induce mutually
inverse bijections.}

\[
  % BEGIN inputs/XXIV/native_diagrams/exp24/011_weil_restriction_torsor_functors.tex
\begin{tikzcd}[column sep=huge,row sep=large]
\mathcal P
  \arrow[r,mapsto]
&
\mathcal P\times_S S'
\\
\displaystyle\prod_{S'/S}\mathcal P'
&
\mathcal P'
  \arrow[l,mapsto]
\end{tikzcd}

% END inputs/XXIV/native_diagrams/exp24/011_weil_restriction_torsor_functors.tex

\]

\emph{In particular,
\[
  H^1(S,G)\simeq H^1(S',G').
\]}

By \SGARefLink{sga3:xxiv:proposition:8-2}{8.2}, it is enough to show
\[
  H^1_S(S',G')=H^1(S',G').
\]
By \SGARefLink{sga3:xxiv:lemma:8-3}{8.3}, this may be checked locally for the local finite étale
topology. We may therefore suppose that \(S'\) is a finite disjoint
union of copies of \(S\), say \(I_S\), where \(I\) is a finite set.
Then \(G'\) is given by a family \((G_i)_{i\in I}\) of sheaves on
\(S\), and
\[
  H^1(S',G')
  \simeq
  \prod_{i\in I}H^1(S,G_i).
\]
On the other hand,
\[
  H^1(S,G)
  \simeq
  \prod_{i\in I}H^1(S,G_i).
\]
It follows from \SGARefLink{sga3:xxiv:proposition:8-2}{8.2} that
\(H^1_S(S',G')=H^1(S',G')\).
\hfill\(\square\)

\par\medskip\noindent\SGAInlineRefTarget{sga3:xxiv:remark:8-5}\textbf{Remarks 8.5.---}\enspace
\SGARefLink{sga3:xxiv:proposition:8-2}{Propositions 8.2} and \SGARefLink{sga3:xxiv:proposition:8-4}{8.4} may be interpreted through the following
exact sequence, where \(f:S'\to S\) is the given morphism:
\[
  1\longrightarrow H^1(S,f_*G')
  \longrightarrow H^1(S',G')
  \longrightarrow H^0(S,R^1f_*G').
\]
In the commutative case, this exact sequence follows from the Leray
spectral sequence. It remains valid in the noncommutative
case.\footnote{\textit{Editors' note.} See \S V.3.1.4 of
\([{\rm Gi71}]\).}

In this form one sees that the result remains valid when \(f\) is
assumed only finite, or even merely integral, provided that the
topology is the étale topology. Indeed, for every \(G'\),
\[
  R^1f_*(G')=\text{the final sheaf}
\]
by SGA~4, VIII, 5.3.\footnote{\textit{Editors' note.} This reference
also points to \S V.3.1.4 of \([{\rm Gi71}]\).}

The result becomes false for a topology such as \textup{(fpqc)} or
\textup{(fppf)}, even if
\[
  S=\operatorname{Spec}(k),
  \qquad
  S'=\operatorname{Spec}(k[t]/t^2),
\]
where \(k\) is an algebraically closed field of characteristic
\(p\neq0\), and \(G'=\mu_p\) or \(\alpha_p\).

Likewise, 8.2 becomes false, even for the étale topology, if the
finiteness hypothesis on \(f\) is omitted. For example, take \(f\)
to be an open immersion with
\[
  S=\operatorname{Spec}(V),
\]
where \(V\) is a complete discrete valuation ring with algebraically
closed residue field, take \(S'\) to be the open subscheme consisting
of the generic point, and take \(G'\) to be the constant group
\((\mathbb Z/n\mathbb Z)_{S'}\), where \(n\) is prime to the residue
characteristic of \(V\). Then
\[
  H^1(S,G)=0,
  \qquad
  H^1(S',G')\neq0.
\]
Replacing \(S'\) by \(S\amalg S'\) gives an analogous example in which
\(S'\to S\) is étale and surjective, hence a cover for the topology
under consideration.

\section*{Bibliography}

\noindent[BLie]\enspace
N.~Bourbaki, \textit{Groupes et algèbres de Lie}, Chap.~I, Hermann,
1960.\par

\smallskip\noindent[Ch51]\enspace
C.~Chevalley, \textit{Théorie des groupes de Lie}, t.~II,
\textit{Groupes algébriques}, Hermann, 1951.\par

\smallskip\noindent[Di57]\enspace
J.~Dieudonné, Lie groups and Lie hyperalgebras over a field of
characteristic \(p>0\). VI, \textit{Amer. J. Math.} \textbf{79}
(1957), no.~2, 331--388.\footnote{\textit{Editors' note.} The editors
have added, after these three references which appeared in the
original, the references that follow.}\par

\smallskip\noindent[Bo91]\enspace
A.~Borel, \textit{Linear Algebraic Groups}, 2nd ed.,
Springer-Verlag, 1991.\par

\smallskip\noindent[Ch05]\enspace
C.~Chevalley, \textit{Classification des groupes algébriques
semi-simples}, with the collaboration of P.~Cartier,
A.~Grothendieck, and M.~Lazard, \textit{Collected Works}, vol.~3,
Springer, 2005.\par

\smallskip\noindent[DG70]\enspace
M.~Demazure and P.~Gabriel, \textit{Groupes algébriques}, Masson and
North-Holland, 1970.\par

\smallskip\noindent[Gi71]\enspace
J.~Giraud, \textit{Cohomologie non abélienne}, Springer-Verlag,
1971.\par

\smallskip\noindent[Ja87]\enspace
J.~C.~Jantzen, \textit{Representations of Algebraic Groups}, Academic
Press, 1987; 2nd ed., American Mathematical Society, 2003.\par

\smallskip\noindent[Jou83]\enspace
J.-P.~Jouanolou, \textit{Théorèmes de Bertini et applications},
Birkhäuser, 1983.\par

\smallskip\noindent[Ta75]\enspace
M.~Takeuchi, On coverings and hyperalgebras of affine algebraic
groups, \textit{Trans. Amer. Math. Soc.} \textbf{211} (1975),
249--275.\par

% HARD STOP: end of Exposé XXIV authority local page 54.

% END inputs/XXIV/tex/components/13_expose_XXIV_remark631_continuation_through_bibliography_en.tex


